% PhysMon research plan of record: Research Redesign v2, revised edition of 11 November 2026.
% Canonical source. Build the PDF and its searchable text with: scripts/build_plan_pdf.sh
% (tectonic in the separate conda environment `texbuild`). Edit this file, not the PDF.
\documentclass[10pt,a4paper]{article}
\usepackage[margin=2.1cm]{geometry}
\usepackage{fontspec}
\usepackage{amsmath,amssymb}
\usepackage{booktabs,longtable,tabularx,array}
\usepackage{enumitem}
\usepackage{xcolor}
\usepackage{fancyhdr}
\usepackage[htt]{hyphenat}
\usepackage{xurl}
\usepackage[hidelinks]{hyperref}
\setlength{\emergencystretch}{3em}
\setlist{nosep,leftmargin=1.4em}
\setlength{\parskip}{0.35em}
\setlength{\parindent}{0pt}
\renewcommand{\arraystretch}{1.15}
\newcolumntype{L}{>{\raggedright\arraybackslash}X}
\newcolumntype{P}[1]{>{\raggedright\arraybackslash}p{#1}}
\pagestyle{fancy}\fancyhf{}
\lhead{\small PHYSMON \textbar{} RESEARCH REDESIGN v2 (REVISED 11 NOV 2026)}
\rhead{\small\thepage}
\cfoot{\footnotesize Working research protocol -- not completed results}
% source and status tags
\newcommand{\tg}[1]{\textsc{[#1]}}
\newcommand{\NE}{\texttt{NOT\_EVALUATED}}
\newcommand{\gate}[1]{\par\smallskip\noindent\fbox{\parbox{\dimexpr\linewidth-2\fboxsep-2\fboxrule}{\small #1}}\par\smallskip}
\title{\vspace{-1.5em}\textbf{PHYSICS-COUNTERFACTUAL MECHANISMS}\\[0.3em]
\large Benchmark Red-Team, DECOY-v2 Reconstruction, and a Gated Mechanistic Research Programme\\[0.6em]
\normalsize Research redesign v2 --- revised edition (supersedes the edition of 14 October 2026)}
\author{PhysMon working document}
\date{11 November 2026}
\begin{document}
\maketitle
\thispagestyle{fancy}
% ===== s00_front =====
\begin{center}
\fbox{\parbox{0.94\linewidth}{\small
\textbf{DOCUMENT STATUS.} This is a proposed, falsifiable experimental protocol, not a record of experiments
performed. It is the single plan of record for PhysMon from Stage~14 onward. This edition (11 November 2026)
revises the edition of 14 October 2026 after (i) the four ICLR 2027 reviews of the Phase-I paper
\emph{Confidence Without Cause} (submission 50608, withdrawn on 11 November 2026), whose numerical claims were
reproduced from the repository's committed data; (ii) the exploratory results of Stages 14.0--14.20; (iii) a
literature update of 11 November 2026; and (iv) four independent reviews: two of a first draft (one by an agent
with repository access, one by an external language model supplied by the user) and two of the assembled draft by
agents with repository access (fidelity to the October edition and the repository; scientific consistency). The October
edition's identifiers (B0--B7, P0--P11, A01--F20, X01--X12) are kept so that the journal, ledger and decision log
still resolve; Appendix~\ref{app:disposition} gives the disposition of every October block. The October PDF is
preserved in Git history (for example commit \texttt{2cf08e2}).\\[0.3em]
\textbf{Authorship.} Revised by the project's coding agent (Claude) at the user's direction; checked against the
repository by independent agent reviews (Appendix~\ref{app:revisions}); not reviewed by a person outside the project. Revision-sensitive
hypotheses, labels, validation decisions and claim scope require explicit decision-log entries.}}
\end{center}
\medskip

\tableofcontents
\clearpage
% ===== s01_status =====
\section{Status, scope, and rules of evidence}
\label{sec:status}

\textbf{Purpose.} An independently reasoned, falsifiable, executable research programme. The October edition
replaced the June 2026 Phase-II plan and added a mandatory benchmark red-team (Phase~B). This edition keeps that
programme and changes four things, each forced by evidence gathered since: (a) the primary endpoint is the model's
own reasoned answer, read against its sampling noise, not a teacher-forced string likelihood; (b) the scientific
question is narrowed to whether the answer follows a quantity's certified \emph{physical role} across values
(Section~\ref{sec:framework}); (c) a measurement-and-calibration gate precedes every behavioural run, because the
7--8B models are at the floor without reasoning and a 32B thinking model was near the ceiling on easy items; (d)
mechanism work declares which causal claim it tests and is read through an interaction estimand. The target is a
scientifically strong future submission, not a guarantee of acceptance.

\textbf{Starting point.} Phase~I is the 25-page ICLR 2027 paper \emph{Confidence Without Cause: Predicting and
Localizing Physics-Shortcut Sensitivity in Language Model Hidden States} (submission 50608). Its four reviews
(ratings 1, 1, 3, 3; all at confidence~2) and the authors' withdrawal on 11 November 2026 are recorded in
\texttt{docs/reviews/iclr2027\_submission50608/} with a script that reproduces 21 of the reviewers' 23 checkable
claims from committed data (one computed on request, one a code reading). No Phase-I claim survives as stated
(Section~\ref{sec:ledger}); its numbers remain \tg{SUBMITTED} values. The differently titled 29-page \emph{Proved,
Not Annotated} file must not be confused with the submission.

\textbf{Critical epistemic constraint (unchanged).} A symbolic certificate that the correct physical answer is
invariant to an edit does not prove that the model's probability of a particular answer string ought to remain
invariant, nor does it prove use of the quantity in the model's reasoning. Effects on probabilities, internal
representations, correct and incorrect answers, and physically faithful computation are separate estimands.
$S$ denotes an observational sensitivity measure only; \emph{mechanism} requires controlled, specific, predictive
intervention evidence (Section~\ref{sec:causal}); \emph{shortcut reliance} is reserved for evidence that the
model's computation improperly incorporates a quantity. This edition adds a corollary: a quantity whose influence
cancels algebraically, or depends on a regime, may legitimately enter intermediate reasoning; its use there is
not an error (Section~\ref{sec:strata}).

\textbf{Thesis to test, not presume.} Some quantities that the physics makes irrelevant compete with relevant ones
during the model's computation, through a structured and partially generalisable pathway. The alternatives are
legitimate outcomes and the design must be able to support each: (i) answer-string and formatting effects;
(ii) generic difficulty or uncertainty; (iii) numerical copying or answer contamination without a change to the
physical computation; (iv) ordinary context sensitivity of conditional distributions; (v) several unrelated
cue-specific mechanisms; and, added in this edition, (vi) correct role selection, in which case the scientific
object becomes the mechanism of correct selection (H6). Do not design analyses that only accept the preferred
hypothesis.

\textbf{Scope strategy.} Phase~B (Section~\ref{sec:phaseB}) establishes what the benchmark measures. The critical
path is a sequence of gates G0--G7 (Section~\ref{sec:gates}) onto which the October blocks P0--P11 are mapped
(Section~\ref{sec:blocks}). The inventory of 104 registered contrasts (Section~\ref{sec:inventory}) remains a
catalogue of candidates, not a work list. Maximise discriminating evidence per unit of compute and human review;
no unvalidated experiment enters a headline claim after outcomes are seen.

\textbf{Source categories.} \tg{SUBMITTED} = reported in the Phase-I manuscript; \tg{REVIEW} = a criticism in the
ICLR 2027 reviews (or the earlier PAT feedback); \tg{LITERATURE} = an independently checked public work
(Section~\ref{sec:bibliography}); \tg{PROPOSAL} = a new experiment, hypothesis, threshold or allocation;
\tg{V} = verified in the repository's committed files; \tg{S14} = an exploratory Stage~14 result with its journal
identifier; \tg{ENV} = a measured property of the cluster or software stack; \tg{I} = an inference or design choice
of this edition. Numerical targets other than \tg{SUBMITTED}, \tg{V} and \tg{S14} values are planning choices, not
results. Unlike the October edition, this edition was written with the repository visible; every \tg{V} and
\tg{S14} statement cites its source.

\textbf{Rules that apply at all times.}
\begin{enumerate}
\item Data roles (EXPLORATION, GENERATOR\_TUNING, INTERNAL\_CONFIRMATION, EPC\_CONFIRMATION,
ADVERSARIAL\_CHALLENGE) are disjoint and assigned by graph lineage; a shadow test set is inaccessible to all
tuning scripts and agents; a structure never moves from EXPLORATION into a confirmation role.
\item The independent unit is the independently authored symbolic structure or template cluster, not a rendering,
paraphrase, numeric instantiation or decoder sample.
\item \NE{} is not a null and not \texttt{WRONG}. Zero evaluable rows are not a measured zero effect. A record
that did not complete is never relabelled wrong.
\item Validation has tiers and every item states its tier: the authoring agent's self-review is not validation;
a qualified panel of independent LLM reviewers (Section~\ref{sec:llmpanel}) is the interim tier used for
confirmation, disclosed as such; independent human expert review is the strongest tier and is never claimed when
only the panel has run. A symbolic certificate is not a certificate of natural-language correctness.
\item Every threshold records when it was fixed relative to the data it is applied to; outcomes do not license
retroactive relabelling; new metrics are not called preregistered after the fact.
\item Every comparator is tagged \texttt{oracle} (needs the gold or the label's inputs at inference) or
\texttt{deployable}; no oracle comparator is called ``non-activation'', ``one forward pass'' or ``cheap''.
\item Every intervention is read as an interaction against a control prompt (Section~\ref{sec:causal}).
\item Each GPU job needs the user's acceptance of its exact command; no job is resubmitted, cancelled or held
without it (repository \texttt{AGENTS.md} and the Sharanga skill).
\end{enumerate}

% ===== s02_judgment =====
\section{Executive scientific judgment}
\label{sec:judgment}

The Phase-I work produced a reusable asset (a family-based benchmark with machine-checked invariance, open
per-rendering records and a decision log, which the reviewers valued) and a construct that did not survive
scrutiny. The next paper needs one of the following outcomes; which one is decided by the evidence, not chosen in
advance.

\textbf{Track M (mechanism, conditional).} A verified causal account of how a quantity's physical role is used:
which information about the quantity and its relation to the system is represented, where it travels, how it
changes the competition between the physically correct answer and a misuse-derived alternative, and which
intervention changes that dependence while preserving correct physics. Its predictions must hold on unseen
structures, values and phrasings, and replicate in a second model. It must be distinguishable from operation
planning in math word problems (Qu and Gomes), anchoring pathways (Owusu et al.) and positional readout copying
(Liu). A single exceptional component is insufficient. Track~M may also be about \emph{correct} selection (H6) if
the models select correctly.

\textbf{Track E (evaluation and measurement).} A rigorous decomposition of which apparent ``shortcut'' signals are
artifacts of reference-answer scoring, prompt formatting, metric asymmetry or data structure, and which survive
as reproducible violations of certified role-by-value behaviour. Phase~I and Stage~14 already supply much of the
artifact side (Section~\ref{sec:ledger}, Appendix~\ref{app:stage14}). Track~E must not claim a mechanism.

\textbf{Track R (robustness and repair, optional).} A selective intervention that improves certified invariance
or accuracy on unseen data, with little damage to legitimate use of relevant quantities, and an advantage over
cheap baselines (a relevance instruction, self-consistency, a longer reasoning budget).

\textbf{Three questions that precede expensive mechanism work} (sharpened from the October edition): (1) Does the
model's answer follow the certified role-by-value pattern, measured on its own reasoned answers rather than on a
teacher-forced string? (2) Is any departure specific to the physical role, rather than to numeral presence,
plausibility, unit, position, marker wording or generic perturbation, and is it specific to physics rather than
shared with non-physics dependency problems? (3) Does an identified computation generalise beyond the selection
panel? If any answer is no, the corresponding stronger vocabulary is not used.

\subsection{Phase-I ledger after the reviews (final)}
\label{sec:ledger}

The October edition listed the \tg{SUBMITTED} values as ``not replicated here''. Stage~14.1 recomputed them from
committed artifacts (\texttt{results/stage14/x01\_reproduction/}; J14.1-008), and the reviewers' audit was
reproduced (\texttt{results/reviews/iclr2027\_submission50608/attempt\_20261111/}). The values reproduce; what
changed is what they mean.

\begin{center}\small
\begin{longtable}{@{}P{2.8cm}P{5.6cm}P{7.4cm}@{}}
\toprule
Item & \tg{SUBMITTED} value (reproduced \tg{V}) & Status after the reviews \\
\midrule\endhead
Behavioural label & Qwen 70/140, Llama 32/140 at 0.5 nats; ``25 shared positives'' is the overlap of the $\hat S$ labels, the 0.5-nat overlap is 22 & Label retired (user, 2026-10-17, J14.2-011). It scores the bare certified string although all 560 Qwen generations begin with ``Answer:'' (C1a--d); 55 of 70 stable families vary by $\geq0.5$ nats across renderings (C2a); 30 pilot families agree on 17 of 30 labels across two answer instructions, Spearman 0.04 (C2b); a correct formula hint gives 3.7--5.5 nats (C3); 40--49\% (Qwen) and 41--56\% (Llama) of base-0 positives flip with the base (X02, J14.2-006) \\
Threshold & 0.5 nats, described as registered & Chosen after the Stage~4 results on a 30-family pilot whose labels do not transfer across the instruction (C2b); the registration date is not a justification \\
Mean-state monitor & L18 AUROC 0.7596 (four renderings averaged) & Reproduces (C5a); 0.649 [0.46, 0.82] on the 37 panel families with no unresolved sibling versus 0.789 on the 98 clustered (C5d); not evaluated under cluster-held-out folds (activations lost) \\
Single-rendering monitor & L16 AUROC 0.7465 & Reproduces (J14.1-008); same leakage caveat \\
``Inexpensive'' comparator & Fitted entropy 0.7270; Wave~2 analogue 0.7048 & Built from the certified answer's own log-probabilities: an \texttt{oracle} comparator. The oracle-free entropy baseline scores 0.520 (C4) \\
Correctness probe & 0.7657 on its own target & 0.380 when scored against the sensitivity label (C5b) \\
Paired comparison & Monitor minus entropy +0.033, 95\% [$-$0.074, +0.141] & Reproduces; it compared the monitor with an oracle comparator \\
Calibration & Monitor Brier 0.2036 vs 0.2140; ECE 0.1041 vs 0.0907; FNR at FPR 0.10 0.714 vs 0.643 & Reproduces (J14.1-008) \\
Frozen expansion & Pooled 0.820 on 15 families; B $\approx$0.667 (n=5), C $\approx$0.6875 (n=10) & Reproduces; cue-type base-rate confound unresolved \\
Causal intervention & L16 four-head recovery 1.870; head 11 alone 1.848; clipped 0.969 & 17 of 20 families invert (patched $S_{lp}<0$; +4.16 to $-$3.08 nats; C8a); head 11 raises the gold first-token logit by 3.76 on base and 3.77 on sensitive renderings, 20/20 each (C8b): the effect is cue-independent; ``cancels'' and ``necessary'' are not supported \\
Cue C & Pooled with irrelevant cues & 5 of 5 Qwen-sensitive with mean 9.02 nats (C6a); belongs to the equivalence family, not the irrelevant-cue core \\
Proximity & Near-match tail; Pearson $r=0.074$ & Exact-match Cue~B families are \emph{less} often sensitive (20/52 vs 33/48; C7a) \\
Validation & Dual expert validation not completed & All 155 templates \texttt{created\_by: agent}; 140 marked \texttt{pi\_validated} by a delegated agent pass (C9) \\
Structural overlap & 34 clusters; 71\% of family-held-out tests share a cluster & 110 of 155 families touched by unresolved pairs (0.71; C5c). Two panel counts use different definitions: 98 of the 135 panel families touch an unresolved pair (C5c, C5d), while 95/135 = 70.4\% have a cluster sibling exposed in family-held-out training (J14.4) \\
Verifier & Certified invariance & Coulomb templates pass at an absolute tolerance of $10^{-3}$; no Cue~C dimension check (C10) \\
Provenance & 182 runs; coverage 0.703 / 0.451 & Reproduces; claim-to-artifact ledger built in Stage~14.0/14.11 \\
Renamed-donor control & Zero evaluable rows & Unmeasured, not a favourable negative control \\
\bottomrule
\end{longtable}
\end{center}

\textbf{What survives.} Descriptive facts: the teacher-forced gold-string likelihood moves with the choice of base
rendering, with the answer instruction, with in-context answer information and with the scoring string. These
support Track~E's methodological point and nothing stronger.

\subsection{What Stage 14 established (exploratory)}
\label{sec:stage14summary}

All Stage~14 results have data role EXPLORATION and are not paper-eligible; details and identifiers are in
Appendix~\ref{app:stage14}. In brief: \tg{S14} a same-unit numeral in the prompt becomes a far more probable
answer continuation (about 10--13 nats for Qwen2.5-7B, 3.5--5 for Llama-3.1-8B), with no detectable dependence
on proximity at the available precision; matching comparison numerals for plausibility removes about 40\% (Qwen)
and 53\% (Llama) of it, leaves incompatible-unit and metadata effects not distinguishable from plausibility, and
leaves a remainder that depends on unit compatibility (J14.9, J14.14, J14.17). Without reasoning both 7--8B
models are almost never correct (3/211, 6/211); with greedy 2,048-token reasoning Qwen is correct on 96/211 and
Llama on 50/211, but a sentence with no number changes answers about as often as numeric cues (J14.7, J14.13).
Qwen3-32B with thinking answered 63 of 64 graded problems on one sample each; four cue kinds on 32 of those
candidate anchors left 109 of 111 graded answers correct, with 17 of 128 cells \NE{}; the resample that would
have measured its own noise was cancelled, so its cue contrast is \NE{} (J14.19-009 to -015). None of this tests
the role-by-value question; the cues in these runs also carried explicit disclaimers (``unrelated'', ``not
connected'') in the prompt text.

% ===== s03_phaseB =====
\section{A mandatory Phase B: red-team, reconstruct, and independently validate DECOY}\label{sec:phaseB}

This section keeps the October edition's Phase~B (v2.0, pp.~4--13): the board decision, the benchmark
architecture, blocks B0--B7 with their numbered sub-items, and Gates B1--B6. Identifiers are unchanged so that
Stage~14 journal entries (J14.x-nnn), the claim ledger and the first revision (v2.1) still resolve. What has
changed is where each requirement now lives. The construct choice that v2.0 left open (B2) is decided in
Section~\ref{sec:framework}; the pilot, confirmation and external work of B5--B7 are re-sequenced into the
critical path of Section~\ref{sec:gates} (G0 integrity closure; G1 physics construct; G2 measurement and model
calibration; G3 behavioural discovery; G4 independent behavioural confirmation; G5 mechanism; G6 external and
functional confirmation; G7 optional practical use). Where this section and Section~\ref{sec:framework}
disagree, Section~\ref{sec:framework} wins.

\textbf{Status vocabulary.} Each block, and each sub-item whose status differs from its block, carries a line
\textbf{Status (11 Nov 2026):} with one of the following values.
\begin{itemize}
\item \textbf{DONE-EXPLORATORY}: executed in Stage~14 under data role EXPLORATION; the cited journal entries hold
the evidence; nothing is confirmatory or paper-eligible.
\item \textbf{PARTIAL}: part of the requirement is met; the open part is named.
\item \textbf{OPEN}: required and not met.
\item \textbf{REVISED}: requirement kept, content changed; the pointer names the governing subsection of
Section~\ref{sec:framework}.
\item \textbf{DEMOTED}: requirement kept, moved to a later gate.
\item \textbf{REPLACED}: superseded by a named rule; the text is kept only so that old citations resolve.
\end{itemize}

\begin{center}\small
\begin{tabularx}{\linewidth}{@{}P{0.9cm}P{4.2cm}P{3.6cm}L@{}}
\toprule
Block & v2.0 subject & Status (11 Nov 2026) & Now governed by \\
\midrule
B0 & Threat model, estimands, novelty audit & REVISED; PARTIAL & Sections~\ref{sec:framework} and \ref{sec:literature}; threat file in G0 \\
B1 & Reproduce and break DECOY-v1 & DONE-EXPLORATORY (B1.5 OPEN only if v1 is reused) & G0 integrity closure \\
B2 & Construct choice & REVISED (decided) & Section~\ref{sec:framework}; B2.3, B2.4 binding \\
B3 & Modular v2 generator & REVISED (extended) & G1 physics construct \\
B4 & Independent validation & OPEN & Interim tier: qualified LLM panel (Section~\ref{sec:llmpanel}); human study later (D2) \\
B5 & External sets as source material & DEMOTED (split) & Non-physics control in G3; EPC in G6 \\
B6 & Pilot, power, data roles & REVISED (staged) & G2 and G3; data roles unchanged \\
B7 & Deliverables; Phase-B go/no-go & Deliverables kept (OPEN); go/no-go REPLACED & G3 hypothesis scorecard and G4 \\
\bottomrule
\end{tabularx}
\end{center}

%-------------------------------------------------------------------------------------------------
\subsection{The board decision and its reasons}\label{sec:phaseB-decision}

\textbf{Board decision (14 October 2026)} \tg{PROPOSAL}. Do not use the existing 155-family DECOY benchmark
unchanged as the sole scientific foundation of a future PhysMon mechanism claim. Do not discard its Phase-I
evidence, and do not substitute a pre-existing public physics benchmark wholesale. Maintain an immutable
DECOY-v1 (historical reproduction and failure diagnosis); conditionally build a DECOY-v2 experimental suite for
controlled causal discovery; establish a separately sourced, separately certified External Physics
Counterfactuals (EPC) collection for confirmation; and use GSM-DC and GSM-Symbolic as non-physics controls.
Treat unit and re-expression invariance as a distinct task, not interchangeable with irrelevant numerical cues.
Neither a large newly generated dataset nor an improved AUROC makes the work valuable unless these tests
establish a precise, nontrivial model behaviour and an independently validated causal explanation.

\textbf{Three questions that precede expensive mechanism work.} (1) Does the claimed behaviour survive scoring
and task-formulation challenges? (2) Is any internal effect specific to the manipulated physical role rather
than a generic output perturbation? (3) Does the identified computation generalise beyond the selection panel?
If any answer is no, the corresponding stronger vocabulary is not used.

\textbf{Scientific reasons} (each a reported property or a falsification target, not a finding that v2 has
succeeded). (i) The original signed base-relative sensitivity label can change when the same four renderings are
re-indexed. (ii) A certificate about a symbolic answer value does not constrain the probability of a chosen answer
string, and physical assumptions asserted in natural language are not checked by free-symbol elimination.
(iii) Cue~C changes how a governing quantity is expressed; it does not add physically irrelevant information.
(iv) Five original Cue~C families, 34 structure clusters, 2,352 unresolved cross-family rendering pairs, no
completed independent expert validation, and a 135-family monitoring cohort with outcome-related filtering
limit inference \tg{SUBMITTED}. (v) GSM-Symbolic and GSM-DC already establish controlled irrelevant-context
tests; novelty cannot be ``symbolic templates plus distractors'' \tg{LITERATURE}.

\textbf{Status (11 Nov 2026):} in force; reasons (i)--(iv) since verified. (i) X02: of the legacy base-0
positives, 40--49\,\% (Qwen) and 41--56\,\% (Llama) turn negative under another base (J14.2-006); the label also
agrees at chance across two answer instructions on the 30 pilot families (Spearman 0.043, 17 of 30) and scores
the bare certified string rather than the instructed format \tg{REVIEW}. (ii) A correct formula hint raises the
gold-string log-probability by 3.7--5.5 nats on five stable families \tg{REVIEW}; the v1 verifier has no Cue~C
identity or dimension check \tg{REVIEW}. (iii) Cue~C pooling was raised by three of four reviewers (issue~9).
(iv) The 34 clusters were rebuilt from the committed leakage table (J14.4-002). REVISED: the decision's single
notion of an ``irrelevant'' cue is replaced by the four role strata R, D, C, K of Section~\ref{sec:strata}.

%-------------------------------------------------------------------------------------------------
\subsection{Benchmark architecture}\label{sec:phaseB-tracks}

Defined before collecting new outcomes. The track names are kept; their status is new.

\begin{center}\small
\begin{tabularx}{\linewidth}{@{}P{2.3cm}P{4.4cm}P{3.4cm}L@{}}
\toprule
Track & Target and contents (v2.0) & Allowed inference (v2.0) & Status (11 Nov 2026) \\
\midrule
\textbf{V1-R} frozen DECOY-v1 & Reproduction and failure analysis: all 155 families, 140 Phase-I behavioural
families, five monitoring exclusions, 15 later transfer families; original signed labels preserved &
Which historic claims reproduce under their original protocol; not prospective generalisation &
DONE-EXPLORATORY: X01, X02 (J14.1, J14.2), legacy X03 (J14.16-005); legacy label retired as an endpoint
(J14.2-011); Phase-I paper withdrawn on 11 Nov 2026 \\
\textbf{V2-A} causal relevance and numeric interference & Whether use of a numeral is conditioned on physical
relevance: new structures, matched relevant and irrelevant edits, no-cue anchors, non-numeric and numeric
controls, numeral $\times$ unit factorial & Controlled relevance-selectivity and cue-interference effects on
specified outcomes & REVISED: now the role-by-value core with strata R/D/C/K
(Sections~\ref{sec:strata}, \ref{sec:rolevalue}) and supporting conditions (Section~\ref{sec:conditions});
dataset v1 (80 structures, 260 instantiations, 11 Stage-14 cells) is EXPLORATION and has no role-by-value
cells \\
\textbf{V2-B} invariance and equivariance & Equivalent units and coordinate systems, admissible
reference-frame changes, dimension and precision stress tests; old Cue~C here & Transformation-consistency
failure; not by itself irrelevant-cue shortcut reliance & Kept: the convertible-unit cell is a supporting
condition (Section~\ref{sec:conditions}); never pooled with a role stratum \\
\textbf{EPC} external physics counterfactuals & Only validated adaptations of PHYSICS-Verified, UGPhysics,
selected PhysReason and SciBench cases, and investigator-independent expert-authored cases & Generalisation beyond
the v2 template grammar, subject to source-sampling and adaptation limits & DEMOTED to G6; nothing sampled or
adapted (B5) \\
\textbf{MATH-X} non-physics controls & GSM-DC dependency graphs, GSM-Symbolic NoOp, unit-stripped or isomorphic
algebraic analogues & Whether effects or mechanisms are domain-general & Split: a small GSM-DC-style dependency
control with the same role manipulation in G3; full cross-domain falsification in G6 \\
\bottomrule
\end{tabularx}
\end{center}

\textbf{Strict version policy} (unchanged). Any semantic revision of a governing equation, assumption, target
variable, accepted answer or cue role increments the family's (now: the structure's) version; nothing is edited
silently in place. Source dataset revisions and hashes are frozen. Once a confirmatory panel is opened it is
not used to tune generation rules, architecture, scoring, parsers, thresholds, site selection or eligibility.
Negative and ineligible examples stay in an immutable manifest with reasons. \tg{PROPOSAL} Under the
role-by-value design a stratum assignment, a role phrasing, a certified response $y^*_r(z)$ and a misuse-function
list are semantic content of the structure and fall under the same rule.

%-------------------------------------------------------------------------------------------------
\subsection{B0. Benchmark threat model, estimand map and external novelty audit}\label{sec:phaseB-B0}

\textbf{Status (11 Nov 2026):} REVISED and PARTIAL. B0.1 and B0.2 are absorbed into
Section~\ref{sec:framework}; B0.3 is partly covered by the review issue register and Stage~14; B0.4 is now
Section~\ref{sec:literature}; the standalone B0 threat file and red-team protocol are OPEN.

\textbf{B0.1 What the benchmark is not.} The certified statement concerns the physical mapping $y^*=f(r,a)$; it
is not a statement about equal conditional model probabilities $p_\theta(\mathrm{string}(y^*)\mid x_i,q)$. An
answer-invariant pair establishes a correct response target, not the model's computation or its internal
dependence on a nuisance. No label is named shortcut reliance without controlled behavioural and mechanism
evidence. Separate records are kept for (a) physical target invariance, (b) observed gold-string drift,
(c) actual-answer shift, (d) distractor-derived answer preference, (e) erroneous dependency on a physical cue,
(f) a causal internal component or pathway. \emph{Now:} (a) is the stratum certificate
(Section~\ref{sec:strata}); (b) is a retired endpoint (Section~\ref{sec:outcomes}); (c) and (d) are O1--O4, with
misuse curves rather than single decoy matches (Section~\ref{sec:rolevalue}); (e) needs H3 or H4 established under
T2 or T3 (Sections~\ref{sec:hypotheses}, \ref{sec:causal}); (f) is a claim type of Section~\ref{sec:causal}.

\textbf{B0.2 Formal object.} A task stores physical state $s$, assumptions $a$, a dependency graph or program
$g$, target $t$, physical answer $y^*(s,a,t)$, candidate cue $z$, semantic role $R(z)$, perturbation operator
$I$ and rendering $\mathcal{R}$. A nuisance intervention satisfies $y^*(I_z(s),a,t)=y^*(s,a,t)$ within a declared
domain $\mathcal{D}$; a relevant intervention is a positive control whose transformed target is computed,
including cases where the numerical answer happens to coincide despite true dependence. Role is never
established by answer-value coincidence alone: a variable can be causally relevant at a parameter setting with
zero local effect. \emph{Now:} the four strata make this precise: D by graph disconnection, C by an identically
zero derivative, K by a regime predicate with a documented counter-regime value, R by a certified response curve
over a declared range of $m\geq3$ values (Sections~\ref{sec:strata}, \ref{sec:rolevalue}). The last sentence of
B0.2 is the reason K is its own stratum and why a single R value never certifies relevance.

\textbf{B0.3 Preregistered threat taxonomy.} Audits attack: incorrect physical semantics, missing boundary
conditions, excluded-domain exceptions, dimensional and unit errors, answer-rounding ambiguities, canonical
strings and tokens, no-cue anchor dependence, arbitrary base choice, irrelevant numerical value versus extra
text, prompt wrapper and generation position, template or source overlap, near-match selection, cue-condition
class prevalence, sampling and decoder variation, annotator suggestion and automation bias, architecture-specific
scoring incompatibility, adaptive head selection, data contamination and publication novelty. Each risk receives
a test, an accountable owner, a possible failure class and an action.
\textbf{Status:} PARTIAL. Tested on v1 or v2 records: canonical string and wrapper (J14.3-003, J14.10-005;
review check C1 \tg{REVIEW}), base choice (J14.2-006), template overlap (J14.4), near-match selection (review
check C7; J14.17-006), automation bias (agent authorship of all 155 v1 templates, review check C9; agent
verification of the v2 set, J14.19-001). The ICLR issue register lists 20 issues, of which 16 were already
prescribed by v2.0, 3 partly and 1 is writing hygiene \tg{REVIEW}. Not yet tested on v2 items: contamination,
sampling variation under a sampled contract (G2), novelty (Section~\ref{sec:literature}). Threats added since
v2.0: role marker words in the prompt (B2.3 status), selection and estimation on the same samples
(J14.19-014), and budget exhaustion read as a wrong answer (O3, Section~\ref{sec:outcomes}).

\textbf{B0.4 Closest-work matrix.} For each close work record: dependency program, control of nuisance and
relevant edits, formal guarantee, human semantics review, difficulty, output metrics, intervention capability,
public data and derivation, external transfer. The differentiator must hold as a joint contribution
(physical-role interventions, honest behavioural endpoints, mechanistic prediction with causal rescue, external
confirmation); v2 is never called novel because it is larger than v1.
\textbf{Status:} REVISED; now Section~\ref{sec:literature}. v2.0's list (GSM-Symbolic, GSM-DC, UGPhysics,
PhysReason, PHYSICS-Verified, mechanism papers) is kept and extended; the closest named works are now
PRISM-Physics, GSM-DC, Qu and Gomes, Owusu, Wiegreffe and Feldman, and Liu (``The Readout Shortcut''). The absence
of a matching benchmark in our search is not proof of novelty (Section~\ref{sec:question}).

\textbf{Deliverable B0:} a versioned \texttt{benchmark\_claims\_and\_threats.md}, a closest-work table with
accessible URLs, an estimand glossary and an independent red-team protocol. \textbf{Status:} the glossary is
Sections~\ref{sec:outcomes} and \ref{sec:causal}; the table is Section~\ref{sec:literature}; the threat file
(owner and action per threat) and the independent red-team protocol are OPEN and belong with G0
\tg{PROPOSAL}.

\textbf{Failure branch} (kept). If physical relevance cannot be distinguished from GSM-DC-style dependency
selection, the research is framed as generic numerical interference or dependency selection, with physics as a
controlled domain, and no physics-specific relevance computation is claimed. It is now tested early by the
non-physics dependency control of G3 and fully in G6.

%-------------------------------------------------------------------------------------------------
\subsection{B1. Reproduce and deliberately break the original benchmark}\label{sec:phaseB-B1}

\textbf{Status (11 Nov 2026):} DONE-EXPLORATORY (closed). Stage~14 X01, X02 and X03 and the verified ICLR~2027
reviews broke DECOY-v1 as B1 intended: 21 of 23 checks reproduce from committed data (including the paper's own Table~1 and Table~4 values and one
sanity check), one more was computed on request and one is a code reading \tg{REVIEW}. No new claim rests on v1 numbers. B1.5 stays OPEN,
and only if v1 is reused for any new claim.

\textbf{B1.1 Provenance reconstruction.} For every v1 rendering retrieve solver expression and assumptions, units
and target, cue value, base index, prompt bytes, answer string, token ids, native chat template, generation
settings, raw log-probability, raw output, parse status, original label, model commit, split and cluster,
intervention metadata; reproduce every 155/140/135 denominator and Phase-I table without filling missing values;
identify the 15 later families as later confirmation; record authorship, instantiation parentage and inclusion
timestamps. \textbf{Status:} PARTIAL by decision. X01 recomputed every recomputable Phase-I value from committed
per-prompt records (J14.1); the P0 gate was recorded as not passed because values were not regenerated through
the model and hooks could not be tested without activations (J14.1-011). Full regeneration is off the critical
path (open decision D1, side track); the extended parity check at the planned hook sites stays in G0.

\textbf{B1.2 Permutation attack.} Legacy signed sensitivity under all four bases: minimum, maximum, positive count,
flip fraction, rank inversions across models; by cue type, law, answer length, sensitivity quantile and cluster;
compared with the symmetric range and robust dispersion and with anchored no-cue contrasts; threshold history
preserved; new metrics never described as preregistered in June. \textbf{Status:} DONE-EXPLORATORY (J14.2-006):
positives by base 0/1/2/3 are 70/77/82/85 (Qwen) and 32/45/69/53 (Llama); within-model Spearman between the base-0
score and other bases is about zero. The legacy label is retired for v2 analyses (J14.2-008, confirmed by the
user, J14.2-011). Anchored no-cue contrasts are superseded by the role-by-value design and the sampling baseline
of G2.

\textbf{B1.3 Submitted-example sanity test.} The torque family's Qwen gold-string log-probabilities
($-45.531$, $-55.562$, $-45.000$, $-46.469$) give about $+10.031$ nats under base~0 and about $-9.093$ under
base~1 \tg{SUBMITTED}. The order-independent range is absolute dispersion, not a signed risk measure.
\textbf{Status:} DONE: executable test in \texttt{tests/v2/test\_base\_rotation\_x02.py} ($+10.03125$ and
$-9.09375$ exactly; J14.2-005).

\textbf{B1.4 Source audit of generated and ambiguous answers.} Check whether the teacher-forced measure needs
parsing; inspect the five monitoring exclusions; recompute inclusion-agnostic analyses; treat parser failures as
an observable; validate example claims against raw outputs; enumerate correctness transitions without
cherry-picking; test answer variants, precision and units. \textbf{Status:} DONE-EXPLORATORY. The label did not
need parsing; the five exclusions applied a parse-failure filter to negatives only (J14.4-004), changing
non-activation predictor AUROCs by 0.007--0.018 (J14.4-011; the monitor's change is \NE{} without activations).
The scorer received the bare certified string, not the instructed format \tg{REVIEW}; a frozen scoring contract
with seven wrappers exists (J14.3-003). Parser failures are now outcome O3 (Section~\ref{sec:outcomes}). The
legacy transition enumeration and manual scoring of large-effect families close with B1: no claim depends on them.

\textbf{B1.5 Certification stress audit.} Review symbol checks and unit identities under changed assumptions:
implicit friction, non-zero pulley mass, coupling through a shared field, feedback with an ``isolated'' system,
sign and direction conventions, frames and unit conversion, piecewise conditions, singularities, domain
restrictions, physically non-positive parameters; test whether the natural language supplies the assumptions;
use positive mutations that must stay valid and negative mutations that must fail. \textbf{Status:} OPEN if v1 is
reused. The v1 verifier applies an absolute $10^{-3}$ tolerance to Coulomb templates and has no Cue~C identity or
dimension check \tg{REVIEW}. For v2 structures the generator's mutation harness covers part of the list (95
mutants, 0 defects uncaught, 0 equivalent rewrites rejected; J14.6-005). \tg{PROPOSAL} The B1.5 list is carried
into G1 as required mutation types for the new D, C and K certificates, where hidden coupling and regime
conditions are exactly what separates the strata.

\textbf{B1.6 Structural overlap and authorship.} Reconstruct the 34 clusters, 147 family pairs and 2,352
unresolved rendering pairs; review edges; re-evaluate held-out results under cluster, law-class and source
splits; test whether template id, law metadata or crude features predict the label. \textbf{Status:}
DONE-EXPLORATORY (J14.4-002 to J14.4-012): 147 edges, 34 clusters covering 110 of 155 families (a lower bound:
no semantic detector); 71\,\% $=110/155$, 70.4\,\% in the monitoring cohort; template-only predictors lose their
family-held-out signal under cluster-held-out folds while the fitted-entropy features do not (J14.4-009); the 155
families contain only 31 distinct renaming-invariant equation programs (J14.6-003). The monitor scores 0.649 on the
37 panel families without an unresolved sibling and 0.789 on the 98 clustered ones \tg{REVIEW}. OPEN if v1 is reused: an audit of what ``unresolved'' meant in the Phase-I overlap tool; re-evaluation under
law-class and problem-source splits (only the cluster split was done); human adjudication of the 147 pairs; the monitor under cluster-held-out folds needs regenerated
activations and belongs to G7 (monitoring).

\gate{\textbf{Gate B1} (v2.0). Publish an audit dashboard of verified, invalid, semantic-ambiguous,
parser-uncertain and missing records, plus a table of metric, base and wrapper robustness. If the Phase-I effect
remains only a teacher-forced string effect, preserve it as such; no new claim that model physics reasoning was
manipulated. If benchmark artifacts cannot be reproduced, stop all dependent causal claims pending provenance
repair.}

\textbf{Reading (11 Nov 2026).} The second branch applies: the Phase-I effect survives only as a descriptive
teacher-forced string effect (the label moves with base, answer instruction, in-context answer information and the
scoring string), and every causal and monitoring claim as stated is withdrawn with the paper. The dashboard's
content exists in the X01 and X02 attempt summaries, the claim ledger and the committed review-claims check; its
assembly into a Phase-I closure note is a G0 item.

%-------------------------------------------------------------------------------------------------
\subsection{B2. Decide the construct before choosing how many examples to generate}\label{sec:phaseB-B2}

\textbf{Status (11 Nov 2026):} REVISED (decided). The construct is fixed in Section~\ref{sec:framework}
(research question, strata, role-by-value design, outcomes, populations, hypotheses). B2.3 and B2.4 remain
binding rules and are restated below with their new connections.

\textbf{B2.1 Hypothesis families.} v2.0 defined H-P (physical relevance selection), H-N (numerical
interference), H-T (text and interface effects on a single reference string) and H-U (unit and representation
transformations), and forbade pooling H-U with the others. \emph{Now} (Section~\ref{sec:hypotheses}): H-N
becomes H1 and H2; H-P becomes H3, H4 and H6; H-T is the retired string construct plus the position and format
controls; H-U is the convertible-unit family (V2-B). The pooling ban is generalised: results are never pooled
across role strata (Section~\ref{sec:strata}, rule 3).

\textbf{B2.2 Primary behaviour as an interpretable contrast.} v2.0 conditions: (i) no cue, (ii) irrelevant cue of
varying magnitude, (iii) style-matched non-numerical text, (iv) generic numerical metadata, (v) relevant edit with
recomputed gold, (vi) adversarially constructed decoy answer; co-endpoints were decoy-candidate probability and
free-answer correctness, with one primary endpoint chosen after the pilot. \emph{Now:} (ii) becomes the D, C and K
cells at $m\geq3$ values; (v) becomes the R cell with a certified response (and the relevant edit of an existing
input stays as a supporting condition); (vi) becomes the prespecified misuse functions $f_h(z)$ with all matches
kept; (i), (iii), (iv) are supporting conditions (Sections~\ref{sec:rolevalue}, \ref{sec:conditions}). The
primary scientific endpoint is O4; O1, O2, O3 and O5 are always reported; candidate log-odds (O6) exist only
inside a causal claim type; teacher-forced dispersion is a diagnostic (Section~\ref{sec:outcomes}). G4 confirms
the role-by-value family fixed in advance (Section~\ref{sec:gates}); the G3 scorecard may add at most one secondary
family, labelled as chosen on EXPLORATION data.

\textbf{B2.3 Relevance is not a keyword} (binding). Match tasks with explicitly disconnected and physically
connected subsystems, but do not fabricate a single-word role flip that changes other physical assumptions or
allows several plausible readings. Every pair gets its own formal state change, recomputed solution and expert
semantic validation (B4). Add subtle relevance that requires causal tracing (an intermediate variable changes
only through an indirect coupling) beside explicitly stated irrelevance, to tell relevance computation from a
memorised ``isolated'' or ``ignore''. \emph{Connection to the strata:} the explicitly disconnected case is
stratum D; the connected case and the indirect-coupling case are stratum R; strata C and K are the hard
non-keyword cases in which the quantity sits inside the target system (Section~\ref{sec:strata}). The
role-by-value design implements the rule by requiring at least two independently validated phrasings of each role
relation, at least one without marker words such as ``unrelated'' or ``separate'', and agreement across phrasings
before any role claim (Section~\ref{sec:rolevalue}). \textbf{Status:} the Stage~14 items do not meet this rule for
role claims: of 1,464 feasible non-governing cells in dataset v1, 1,085 say the quantity is ``unrelated'', 39
``not connected'', 11 ``takes no part'', 9 ``in a separate experiment'' and the rest use other separating wording,
so the prompt itself states the role (v2.1, \S5.1). Those records bear on presence and unit effects
(J14.14-005, J14.17-006), not on role selection.

\textbf{B2.4 Eligibility strata before scoring} (binding). Report all valid problems (intent-to-evaluate),
format-parseable records (only for metrics that truly need parsing), base-correct items, and model-positive items
(eligible for sensitivity reduction or causal recovery), each with its denominator and selection mechanism. A
result conditional on base correctness is not a general result about all questions; the original signed effect
never substitutes for the question whether the model outputs an incorrect answer. \emph{Connection to the
populations} (Section~\ref{sec:populations}): all valid problems is the population ``all validated tasks'' and is
always reported; format-parseability is no longer a population but the status outcome O3, and a record that is not
\texttt{COMPLETE} is never relabelled wrong; base-correct is ``tasks solved reliably without a cue'', which needs
repeated samples (one correct sample makes a candidate anchor only, J14.19-009); model-positive items are a
sub-population used only inside G5 and G7 with their selection stated. The calibrated intermediate-difficulty
population is new (G2), and its selection samples are never reused for estimation. Every result names both its
role stratum and its population.

\textbf{B2.5 Metric calibration and negative controls.} Kept: absolute differences and robust quantiles; no
ranking of architectures by the share crossing 0.5 nats; a permitted null distribution of cue-free and
formatting-only changes, now the sampling baseline (self-flip floor from fresh no-cue samples) of G2 together with
the style-matched non-numeric sentence; the no-cue reference is not called universally neutral until the
omission-versus-added-text asymmetry is tested, which the non-numeric sentence does. Symmetric range, first-token
and full-sequence scores and varied output strings are diagnostics; output-format sensitivity is also covered by
the three grading views (Section~\ref{sec:outcomes}). The legacy signed label is retired.

\gate{\textbf{Gate B2} (v2.0). If H-P or H-N produces reliable behaviour beyond generic text and answer-string
effects, proceed to a narrow causal benchmark emphasising that phenomenon. If only H-T survives, create a
measurement and LM-interface project and freeze physics-reliance claims. If only H-U survives, split it into a
formal equivariance study, not the core programme. If all effects disappear under realistic controls, archive a
carefully documented negative result and do not manufacture additional distractor conditions to force a positive
headline.}

\textbf{Status:} REPLACED as a decision rule by the G3 hypothesis scorecard over H0--H6
(Section~\ref{sec:gates}). Its branches correspond: H-P or H-N reliable $\rightarrow$ H1--H4 readings; only H-T
$\rightarrow$ readout and interface claims of type T1 only; only H-U $\rightarrow$ the separate equivariance study;
everything absent $\rightarrow$ H0 reported as a null with its equivalence bounds and power. Its last sentence
stays in force without change.

%-------------------------------------------------------------------------------------------------
\subsection{B3. A modular v2 benchmark with mutually discriminating controls}\label{sec:phaseB-B3}

\textbf{Status (11 Nov 2026):} REVISED (extended with the role strata and the role-by-value cells). The typed
generator and 80 structures exist (DONE-EXPLORATORY: J14.6, J14.15, J14.18, J14.19); strata C and K, role-by-value
cells, validated phrasings and misuse functions are OPEN and belong to G1. Target size: 250 structures, then 500
(Section~\ref{sec:datasetplan}).

\textbf{B3.1 Five problem-structure strata} (kept). (i) Single governing equation with a natural nuisance
(calibration and causal simplicity); (ii) multistep dependency graph with the distractor attached to a
non-ancestor node; (iii) interacting subsystems with an explicit or inferable physical disconnection;
(iv) conditional laws or boundary constraints where relevance depends on a stated regime; (v) counterfactual role
reversal, where the same measured quantity is relevant in one verified world and irrelevant in another. Include
inverse problems and sign or vector questions where answer and representation validation are reliable. Do not
assume thermodynamics or multimodal diagrams are easy to certify; admit only validated examples and keep excluded
candidates in a selection log.
\emph{Extension.} Problem structure and role stratum are two axes: each structure records its problem-structure
stratum and which role strata it admits; infeasibility is recorded, never imputed (Section~\ref{sec:rolevalue}).
Typical, not definitional, correspondences: (i) often yields C (e.g.\ a pulley radius that cancels when the
moment of inertia is given as a fraction of $MR^2$) or D; (ii) yields D for a non-ancestor node and R for an
ancestor; (iii) yields R/D twins; (iv) yields K; (v) is the role-by-value design itself (R and D cells of one
structure at the same values). The Stage~14 \texttt{role\_reversal} cell changes the question (the answer minus
system~B's value), so it tests use on request, not role with the question fixed (v2.1, \S5.1); it stays as the
supporting condition ``role reversal (use on request)'' and does not implement (v).

\textbf{B3.2 Typed executable graph generator.} Each node stores dimension, unit, reference frame, sign
convention, quantity type, provenance (measured, computed, assumed), uncertainty where relevant and parent edges;
each structure has a solver program and a small proof object (assumptions, symbolic identities, relevance
certificate, admissible cue domain, tested boundary cases); hard instances are checked by more than one route,
not by several outputs of one symbolic generator. \textbf{Status:} DONE-EXPLORATORY for the v2.0 scope:
dimensional checker, typed DAG, eleven-check certifier, proof object, mutation harness and independently coded
solvers (J14.6-001, J14.6-007); 80 structures in 11 domains with gold traces rendered from the graph
(J14.19-002, J14.19-003). Blind-solver agents are a further route of the same model family, not an independent
one (J14.19-001). REVISED \tg{PROPOSAL}: the Stage~14 generator keeps nuisance quantities out of the graph
(J14.6-001), which can certify D-type absence only. Strata C and K need the quantity inside the governing
relations, with certified symbolic cancellation (C) or a regime predicate and counter-regime value (K); R needs
a certified response $y^*_R(z)$ over the declared range, separated from $y^*_D(z)$ by more than the loosest
tolerance; every D, C and K cell
also needs its frozen misuse functions $f_h(z)$ (Sections~\ref{sec:strata}, \ref{sec:rolevalue}). This is G1
work.

\textbf{B3.3 Core interaction.} v2.0 crossed physical relevance (governing versus non-governing), numerical
proximity (near-match versus a predeclared far band) and unit relationship (same or equivalent versus
incompatible), with diagnostic cells separating digit overlap from magnitude (close magnitudes with little lexical
overlap; far magnitudes with shared leading digits; equal numerals with incompatible units; unit-compatible far
distractor; unit-incompatible but lexically overlapping distractor), cluster-level inference, and empty or
impossible cells recorded rather than filled. \textbf{Status:} REVISED. The primary design is role $\times$ value
$\times$ phrasing within structure (Section~\ref{sec:rolevalue}). Proximity and unit survive as supporting
conditions (Section~\ref{sec:conditions}): proximity as three distinct conditions (exact match, near within the
grading tolerance, tolerance-separated); unit as same, convertible and incompatible with the same digits. The
digit-versus-magnitude cells are secondary diagnostics. The Stage~14 ``near'' cue lies within 1\,\% of the gold
(J14.6-006), inside the relative-error view's tolerance; such cells are flagged source-ambiguous
(Section~\ref{sec:outcomes}). Infeasible cells were recorded (211 of 220 feasible on S01--S10, J14.6-006; 2,764
of 2,860 in dataset v1) and the rule stays.

\textbf{B3.4 Factor families expanded only after the core signal.} Placement and recency, ordering, number of
nuisance values, repeated mentions, prose versus table or diagram, explicitness of isolation, negation,
explanation length, token count, decimal precision, rounding, sign and extremes, time scale, unit scaling,
reference-frame change, instantiation, paraphrase and language, law familiarity, chain length, attachment to an
intermediate versus the final target, competing governing laws, contradiction detection. Each secondary factor
gets a paired negative control and at least one alternative causal explanation; designs use balanced incomplete
blocks or fractional factorials with estimable interactions and a documented design matrix; no full Cartesian
product is claimed. \textbf{Status:} kept; two factors are promoted into the core design: explicitness of
isolation becomes the marker-word phrasing factor (Section~\ref{sec:rolevalue}) and cue position becomes a
supporting condition (Section~\ref{sec:conditions}). The rest wait for G3.

\textbf{B3.5 No-op robustness versus valid sensitivity.} Each structure has at least one irrelevant edit with an
invariant target and, where feasible, a matched relevant edit with a known changed solution; a coupled-system
counterexample makes the formerly irrelevant quantity governing, so a model that ignores every extra number
cannot pass both; relevant positive controls are scored on the new target, not on invariance; ill-posed or
under-specified prompts appear only in a labelled rejection task. \textbf{Status:} kept and built into each
complete role-by-value structure: the D, C, K cells are the no-op edits, the R cell is the matched relevant edit
scored by tracking of $y^*_R(z)$ (O4), and the R twin of a D quantity is the coupled-system counterexample.

\textbf{B3.6 Distributions matter} (unchanged). Keep separately (a) naturalistic distractor sampling, (b) balanced
experimental cells, (c) budget-constrained adversarial optimisation. Freeze (a) and (b) before seeing model
scores; adaptive search examples carry attack budget, query history, target model and seed. Never estimate typical
failure rates from an adversarially optimised subset. Hold out newly generated source graphs and cue-value seeds
for the final test (INTERNAL\_CONFIRMATION and the shadow test set, B6.3).

\textbf{B3.7 Difficulty and ecological validity.} Measure difficulty from physical structure (minimum derivation
steps, dependency nodes, applicable laws, constraints, domain conditions), supplemented by blinded human solve
time; low model accuracy does not establish structural difficulty; keep constructed minimal pairs for mechanism
discovery separate from human-readable external questions for final evidence. \textbf{Status:} kept; REVISED by
model calibration: the calibrated intermediate-difficulty population (G2, Section~\ref{sec:populations}) and the
relevant role's no-cue difficulty as a covariate (Section~\ref{sec:rolevalue}). Blinded human solve-time
assessment is OPEN. Stage~14 levels are the author agents' judgement, and the ``upper-university'' items are
mostly introductory in their number of steps (J14.19-003).

\textbf{Deliverable B3} (v2.0): proof-carrying generator; a fully specified factorial design matrix with
feasibility flags; source and variant lineage; formal invariance and positive-control tests; the first 10 new
certified structures and a discovery-only sample of roughly 40--60 independent structures for the gate pilot.
\textbf{Status:} PARTIAL. Generator, lineage, mutation report and an 11-cell factor manifest with feasibility
flags exist (J14.6-007); 80 certified structures exist (S01--S30, S31--S95), more than 40--60 in number, but all are
UNREVIEWED, all are EXPLORATION, and none has role-by-value cells. OPEN in G1: the strata taxonomy, the
role-by-value generator with phrasings and misuse functions, and feasibility counts of \emph{complete}
role-by-value structures, which are what the calibration gates count (Section~\ref{sec:rolevalue}); the
8--12-structure design-validation pilot of G2 is the first checkpoint for the new design.

\textbf{Failure branch} (kept). If role-matched cells are artificial or impossible, redesign the generative
grammar or narrow the causal question; never loosen verification. Structures are not edited to force feasibility
after the G1 freeze.

%-------------------------------------------------------------------------------------------------
\subsection{B4. Human validation as an independent scientific experiment}\label{sec:phaseB-B4}

\textbf{Status (11 Nov 2026):} OPEN; requirements unchanged. No independent validator exists and none has been
qualified. Validator-qualification materials (v2) and a 33-item planted-defect bank exist (J14.5); the user
reviewed the bank (27 items defective as declared, 6 valid controls; J14.5-013) as one non-blinded reviewer who
did not state physics expertise, so physicist confirmation of the bank is open. S01--S20 had the user's open
walkthrough. On 2026-10-26 the user directed agent verification (author, blind-solver and plan-reviewer agents)
in place of the physicist review (J14.19-001): all 80 v2 structures carry admission category
\texttt{UNREVIEWED} with an agent-verification record, and no text may call them expert-reviewed. Agent checks
share the authors' model family and can share their blind spots (J14.19-001). For v1, all 155 templates were
created by an agent and 140 were marked validated through a delegated agent review \tg{REVIEW}; the reviewers'
objection to this (issue~10) applies equally to the v2 set. \textbf{Consequence:} no structure has passed Gate~B4.

\textbf{Interim validation tier: a qualified LLM review panel} (user direction, 11 November 2026)\label{sec:llmpanel}.
Recruiting independent physicists takes time the project does not want to spend now, so confirmation (G4) uses a
panel of independent LLM reviewers as the interim validation tier, and human expert review is a later study that
measures the panel and can upgrade items.
\begin{enumerate}
\item \emph{Panel.} Two to three judge models from model families other than the authoring model (Claude), and
where possible from different families from each other, at least one a strong reasoning model; the panel
composition, versions and prompts are frozen in G1 (open decision D2 now asks which models). The authoring agent
is never on the panel, and neither are models from the families under test (Qwen, Llama); a judge never sees
another judge's verdict, any model-under-test output or score, or the study hypotheses, in either pass, and does
not see the gold trace in its first pass.
\item \emph{Protocol.} The B4.2 two-pass rubric, run per judge: pass 1 solves the problem and states each
quantity's role stratum from the text alone; pass 2 checks the certificate claims (stratum, invariance or certified
response, units, assumptions, naturalness, whether the renderings differ only as intended) and returns a structured
verdict with severity and reasons.
\item \emph{Qualification before use.} Each judge, with its frozen prompt, is run on the 33-item planted-defect bank
(27 defective and 6 valid controls, labelled by the user, J14.5-013) and must meet thresholds fixed in G1 (planning:
detect at least 85\% of the defective items for the declared reason and accept at least 5 of the 6 valid ones).
Because six valid controls estimate false rejection very imprecisely, G1 adds freshly generated valid items and
new planted defects to the bank before qualification, and holds a portion out unseen; a judge or prompt changed
after seeing bank results is requalified only on that held-out portion or on a fresh planted set. The bank is not
used to tune prompts.
\item \emph{Admission rule.} An item is \texttt{LLM\_PANEL\_VALIDATED} when every qualified judge accepts it with
no high-severity defect; split verdicts go to a recorded adjudication whose adjudicator and tier are stated (the user, who
labelled the bank, is not an independent adjudicator), and adjudicated items are reported separately; items with any
unresolved high-severity flag are excluded and counted. Agreement statistics (raw agreement, Cohen's $\kappa$,
Gwet's AC1) are reported as for human validators.
\item \emph{Claim language.} G4 results on panel-validated items are ``confirmatory under LLM-panel validation'';
no text calls them expert-validated; agent authorship and LLM validation are disclosed.
\item \emph{Human study later.} A stratified sample of panel verdicts (accepted, rejected, split) is later reviewed
by independent humans (G6 or P11) to estimate the panel's error rates; items that pass are upgraded to the
expert tier.
\end{enumerate}
Without a qualified panel, G4 is run but labelled EXPLORATORY-REPLICATION, and no text calls it confirmatory.

\textbf{B4.1 Audit before claiming certified English.} Every v1 family gets at least one qualified blinded
semantics review, and every family used in a publishable confirmatory claim gets two independent reviews with
third-party adjudication where needed; the validation pilot samples across cue types, physics domains, largest
likelihood shifts, negative scores, presumed easy cases, ambiguous coupled systems and high-overlap clusters;
authors never self-certify the semantic naturalness of their own families. \emph{Now:} the v1 part applies only if
v1 is reused; the two-review rule applies to every G4 structure and every adapted EPC item (interim: the qualified
LLM panel of Section~\ref{sec:llmpanel}), and sampling for the
pilot also spans the role strata.

\textbf{B4.2 Two-pass review instrument.} Pass~1, without the authors' equation: solve or outline the physics,
state assumptions and target, whether the cue could affect the target, and whether the wording permits another
physical reading. Pass~2, with the encoded graph and certificate: check graph, units, domain, numerical answer,
cues and any omitted dependency. Validity, clarity, difficulty and naturalness are judged separately;
uncertainty is marked and counterexamples annotated. \tg{PROPOSAL} Under the strata, Pass~1 also asks the reviewer
to assign the quantity's stratum (R, D, C, K) blind and, for K, to state the regime; Pass~2 checks the stratum
certificate, the certified response or invariance range, the misuse functions and each role phrasing separately.

\textbf{B4.3 Qualification and independence.} Screen reviewers with calibrated problems including boundary-condition
traps, unit conversions and misleading ``isolated'' claims; record expertise by domain; balance assignments
across authors, sources, cue types and difficulty; reviewers are blind to model outcomes and to probe or circuit
hypotheses; estimate payment and time in a small pilot, then resource the whole set. Disagreements go to a
domain-qualified independent third reviewer. Agreement is reported as raw agreement, Cohen's kappa on substantive
claims, uncertainty intervals and category-specific defect rates, with unresolved ambiguities listed.

\textbf{B4.4 Adversarial validation tests.} Before release insert predeclared planted defects (hidden coupling,
missing assumption, relevant variable removed from the equation, failed transformation identity, inconvertible or
inconsistent units, invalid parameter domain, insufficient boundary condition, invalid precision, contradictory
distractor statement) and a valid equivalent rewrite; measure detection sensitivity and false-rejection rate by
mutation type; if reviewers miss a defect type, revise the checklist and recalibrate before confirmation. The
planted bank never enters the benchmark as genuine physics.

\textbf{B4.5 Admission categories.} \texttt{CERTIFIED\_AND\_VALIDATED}, \texttt{SOLVER\_VALID\_LANGUAGE\_AMBIGUOUS},
\texttt{SOLVER\_INVALID}, \texttt{ANSWER\_AMBIGUOUS}, \texttt{NATURALNESS\_FLAG}, \texttt{REVISED\_REVALIDATED},
\texttt{EXCLUDED}, \texttt{UNREVIEWED}. Report fractions before and after repair; store a signed reviewer record;
never count repaired examples as untouched independent examples; publish the protocol and aggregate outcomes
without reviewers' personal identifiers. Agent verification is recorded beside these categories and never
changes \texttt{UNREVIEWED}.

\gate{\textbf{Gate B4} (unchanged). Confirmatory physical-relevance claims are blocked until independent expert
semantic validation, version-correct solver checks and the parser criterion for the relevant endpoint pass. No
arbitrary agreement threshold substitutes for reviewing concrete physics defects. High disagreement or
mutation-test failure means a new validation cycle, not optimistic averaging.}

%-------------------------------------------------------------------------------------------------
\subsection{B5. External datasets as source material, not borrowed certificates}\label{sec:phaseB-B5}

\textbf{Status (11 Nov 2026):} DEMOTED and split; nothing has been sampled or adapted. (1) A small non-physics
dependency control (GSM-DC-style problems with the same role manipulation) enters G3 as a supporting condition
(Section~\ref{sec:conditions}), so that a general numeric or dependency phenomenon is recognised before G4.
(2) The external physics sets EPC-D, EPC-C and EPC-N (B5.1--B5.5) and the full cross-domain falsification (B5.6)
belong to G6, with the second model. Rule (Section~\ref{sec:rolevalue}): a counterfactual built on an external
problem is a new item; it needs its own certificate, including its stratum certificate, and its own B4
validation; it never inherits the source's certificate or answer check.

\textbf{Why the hybrid design matters.} Public benchmarks supply diversity and independent authorship but not
DECOY-specific nuisance and relevance interventions or proofs of a modified problem's target. Once a source problem
is rewritten, the adaptation is new benchmark data that needs independent verification. Synthetic v2 and adapted
public questions serve different causal and generalisation purposes.

\begin{center}\small
\begin{tabularx}{\linewidth}{@{}P{2.4cm}P{4.2cm}LP{3.1cm}@{}}
\toprule
Source & Public properties \tg{LITERATURE} (as stated in v2.0) & Adaptation and main danger & v2.0 decision \\
\midrule
PHYSICS-Verified (Yale NLP) & 1,034 PhD-qualifying-exam problems, 2,559 individually scored answers, cleaned
reference solutions, six areas & High realism; source checks do not validate an inserted cue; proof and
multimodal items excluded; small models may fail most advanced items & Primary hard source, screened for
tractable interventions \\
UGPhysics (ICML 2025) & 5,520 undergraduate questions, 13 subjects, English and Chinese, seven answer types & Broad;
licence and derived-work permissions need review; not every item admits a certifiable intervention & Primary
broad source with per-subject selection record \\
PhysReason (ACL 2025) & 1,200 problems; knowledge and reasoning strata; process evaluation & Multistep paths;
reference rationale is not an internal trace and does not certify a nuisance role & Secondary deep-structure
source after reconstruction \\
SciBench (ICML 2024) & College-level scientific tasks from course material & Unsuitable questions or figures;
copyright constraints & Optional, if feasibility is shown \\
GSM-DC (EMNLP 2025) & Symbolic dependency graphs, distracting-context injection, stepwise accuracy & Similar
behavioural confounds; mathematics, not physics & Required competitor and cross-domain control \\
GSM-Symbolic (ICLR 2025) & Symbolic variants and irrelevant NoOp clauses & Shows numeric and context
fragility; establishes no PhysMon computation & Required comparator \\
\bottomrule
\end{tabularx}
\end{center}

PRISM-Physics, now among the closest works (Section~\ref{sec:literature}), is a further possible source under the
same no-borrowed-certificate rule.

\textbf{B5.1 Outcome-independent sampling.} Per source document revision, inclusion criteria, population, seed,
subject and answer-type strata, sampling probabilities, original id and provenance, removed figures and reuse
constraints. Never pick the published problems a model fails and adapt the dramatic ones for a
``representative'' external test.

\textbf{B5.2 Conversion pipeline.} (1) identify the original target and material assumptions; (2) solve
independently or reconstruct an executable dependency program; (3) identify safe intervention points and
candidate quantities (now: candidate quantities for each feasible stratum); (4) derive the output under both edits
(now: the certified response or invariance over the declared values); (5) preserve essential semantics;
(6) independent physics and naturalness review of each adapted pair; (7) audit overlap with v1 and v2 and
training-contamination risk; (8) freeze adapted items and scoring before evaluating models. Reason-coded
rejections are logged at each step.

\textbf{B5.3 Feasibility pilot.} A precommitted, stratified 20--40-candidate sample from each of
PHYSICS-Verified and UGPhysics, optionally a smaller PhysReason sample; measure adaptation yield, reviewer hours,
certification rate, difficulty and baseline solve rate, with intervals on yield and a representativeness
comparison of retained and rejected items. It decides source suitability; it is not an external test.

\textbf{B5.4 Independent external sets.} (i) EPC-D, an adapted development set that may inform methods and
tooling but never confirms an already tuned claim; (ii) EPC-C, a locked confirmation set authored, adapted and
verified by different staff after the target and code are frozen; (iii) EPC-N, new expert-authored problems not
derived from public benchmarks, against contamination and reference leakage. Aim for heterogeneity across law,
source and difficulty rather than more paraphrases. The final external set is opened once; any revision after a
failure needs a new independent set.

\textbf{B5.5 Source- and solvability-aware estimands.} Report all eligible cases, base-correct cases and the
available physically meaningful wrong-answer alternatives; if small models fail hard items, disclose the selection
and the shifted population instead of keeping only successes; use relevant-edit responsiveness (now the R-cell
response, O4) even without baseline correctness; verify native unit and answer scoring for each model before
comparing them.

\textbf{B5.6 Cross-domain falsification.} Select abstractly matched GSM-DC and GSM-Symbolic structures, measure
interference under their own validated protocols, then evaluate a carefully translated subset under PhysMon's
controlled relevance conditions. If a supposed physics relevance gate responds identically to graph-equivalent
grade-school arithmetic, read the mechanism as general dependency filtering or numerical copying unless a truly
physics-specific manipulation separates them. \emph{Split:} the small control in G3 gives the early behavioural
reading; the full version, including mechanism transfer, is G6.

\gate{\textbf{Gate B5} (kept; now applied in G6). If adaptations are rare or ambiguous, preserve source rejection
rates and either restrict the external question types or commission independent expert-authored tasks; never
downgrade certification to inflate numbers. If only synthetic tasks show effects, make that limitation central and
do not claim broad physics reasoning. Confirmed external effects are required for a broad generalisation claim.}

%-------------------------------------------------------------------------------------------------
\subsection{B6. Pilot, power, confirmation and anti-cherry-picking rules}\label{sec:phaseB-B6}

\textbf{Status (11 Nov 2026):} REVISED: staged into G2 (measurement and model calibration) and G3 (behavioural
discovery); the data roles are unchanged. Stage~14 pilot evidence is EXPLORATION only: both 7--8B models were near
the correctness floor without reasoning (J14.9-003, J14.15); reasoning lifted it (J14.13, J14.15-008); Qwen3-32B with
thinking answered 63 of 64 graded no-cue problems correctly with 14 of 78 records \NE{} (9 truncated, 5
unreadable LaTeX-wrapped final lines; J14.19-009); on 128 cue cells over 32 one-sample anchors, 2 were graded wrong
and 17 were \NE{} (J14.19-013, J14.19-014). The sampling-noise resample was cancelled (job 435606, J14.19-015), so
the cue effect against the model's own variation is \NE{}, and the anchor-based baseline is selection-biased by
construction. This is the worked example behind G2's separation of selection and estimation samples.

\textbf{B6.1 Pilot matrix and endpoints.} v2.0: roughly 40--60 new discovery structures to debug construction and
estimate variance, not to prove a small effect; a balanced attainable subset of the core cells with no-cue and
metadata conditions; at least two numeric instantiations; a fixed output grammar; two independently implemented
dense models to start and a stronger model to test persistence beyond 7--8B without presuming it; primary cue
family and endpoint preselected only on the discovery partition. \emph{Now:} G2 holds the grading contract (three
views, parser audit), the sampling baseline (repeated samples; selection samples never reused for estimation), an
8--12-structure design-validation pilot of the role-by-value cells, and the model and difficulty regime; G3 runs
the staged core conditions on calibrated items and scores the hypotheses. ``Two instantiations'' becomes
$m\geq3$ values per role ($m\geq5$ with held-out values for the dependency curve); the fixed output grammar is the
declared evaluation contract (Section~\ref{sec:outcomes}); preselection on discovery data only is kept.

\textbf{B6.2 Precision at the independent unit.} Simulations and blinded design checks with random effects for
structure, source or author, instantiation and model; target intervals for paired differences and the core
interaction, for causal rescue (G5) and for out-of-cluster AUROC margins (now G7, monitoring); a power target and
minimum effect with a scientific interpretation. The October 300--600 band is replaced by the dataset plan of
Section~\ref{sec:datasetplan} (250 structures, then 500); independent dependency programs are worth more than paraphrases (v1 had only 31 distinct
equation programs in 155 families, J14.6-003). \emph{Now:} the interaction targeted is the role-by-value interaction
and the per-stratum tracking-loss, misuse-share and neglect-share estimands (Section~\ref{sec:rolevalue}); the power simulation
uses the variance of those estimands measured in G2, not only no-cue variance.

\textbf{B6.3 Five disjoint data roles} (kept unchanged). \texttt{EXPLORATION} (pilot and unrestricted scans);
\texttt{GENERATOR\_TUNING} (naturalness and control debugging, non-confirmatory); \texttt{INTERNAL\_CONFIRMATION}
(never used for choosing labels, models, heads or thresholds); \texttt{EPC\_CONFIRMATION} (independent external
source); \texttt{ADVERSARIAL\_CHALLENGE} (separate worst-case claim); plus a \emph{shadow test set} inaccessible to
all tuning scripts and agents (v2.0, P3.7). Group by graph lineage and template structure before splitting; where
related groups cannot be separated without losing power, state the limitation and narrow the generalisation claim.
\textbf{Mapping:} the 80 existing structures and all Stage~14 records are EXPLORATION, and G2 and G3 runs on them
remain EXPLORATION; generator and phrasing debugging for the role-by-value design is GENERATOR\_TUNING; new
structures for G4 are INTERNAL\_CONFIRMATION, untouched by any run before G4 opens, with a shadow subset held back by
lineage; EPC-C is EPC\_CONFIRMATION (G6). A structure never moves from EXPLORATION into a confirmation role.

\textbf{B6.4 Primary hypothesis and the factorial null.} v2.0's example primary (a non-governing same-unit
near-match moves correct-versus-decoy log-odds more than a matched far distractor while a governing change is
tracked correctly) is REVISED: the primary endpoint is the role-by-value response O4 under hypotheses H0--H6
(Sections~\ref{sec:rolevalue}, \ref{sec:hypotheses}); log-odds are O6 and live inside a causal claim type. Kept:
the exact estimand, interaction and candidate set are frozen with the model and prompting regime; lexical copying
is a declared alternative (now H1); a predeclared no-practical-effect interval (the equivalence margin); structure-
clustered paired tests, hierarchical regression and within-cluster permutation tests must agree qualitatively
before a finding is called stable.

\textbf{B6.5 Reliability and adversarial budget.} Reuse cue seeds across models; repeat scoring; check numerical
precision and GPU backend variation (Stage~14: bfloat16 reproducible, dtype noise SD 0.09--0.19 nats, while the greedy
answer differed between float32 and bfloat16 in 42 of 211 Llama and 13 of 209 Qwen records, though the correctness
label changed in only one or two; J14.9-003, J14.9-004); sample free generation and report median and
tail outcomes (now required: repeated samples at each model's published defaults, never pooled with greedy
records; Section~\ref{sec:outcomes}). Adaptive cue search has a fixed edit space and per-structure query budget;
attack success is reported apart from fixed-distribution robustness; no monitor is trained and no causal site is
selected on adversarially targeted test items.

\gate{\textbf{Gate B6} (v2.0). Promote to heavy mechanistic research only when a validated causal contrast exists on
held-out structural units, is not dominated by format or tokenisation or by one cue subtype, and is large enough to
support meaningful interventions. If the effect is absent on two independent sources, stop the broad claim and
report boundary conditions. If model output shifts only within teacher-forced string probabilities, restrict the
target to a calibrated likelihood or interface mechanism rather than physical shortcut reliance.}

\textbf{Status:} kept in substance and enforced by the critical path: G4 confirms the role-by-value family fixed
in advance and the G3 scorecard may add at most one secondary family;
G4 supplies the held-out structural units and the format-robustness check; G5 starts only after G4 (for H3, H4, or
H6-RD; H6-extended only if G4 covered C and K). A shift
confined to teacher-forced string probabilities supports at most a T1 readout claim (Section~\ref{sec:causal}).

%-------------------------------------------------------------------------------------------------
\subsection{B7. Benchmark deliverables and acceptance rules for the later mechanism study}\label{sec:phaseB-B7}

\textbf{Status (11 Nov 2026):} the release package, quality-gate checklist and benchmark-only value test are kept
and OPEN (assembled across G1, G4 and G6); the final Phase-B go/no-go is REPLACED by the G3 hypothesis scorecard and
G4.

\textbf{Release package} (path strings illustrative). A frozen v1 manifest (legacy ids and exact original
scoring); v2 structures (typed graphs, assumptions, symbolic and numeric solvers, invariance and relevant-edit
certificates, now stratum certificates and misuse functions); v2 renderings (all cells, generators, lineage, text
hashes); validation (blinded decisions and adjudication); scoring (model, wrapper, answer-token ids, parser
versions); external adaptations (source ids, acceptance and rejection ledger, licences, validator decisions);
splits (lineage and cluster firewall); controls (no-cue, neutral text, generic numerals, role reversals, ill-posed
negatives); results (estimands, eligible denominators, uncertainty and all failed runs); a benchmark card
(intended uses, disallowed inferences, known limits, provenance, contamination).

\textbf{Quality-gate checklist.} Independently correct physics; explicit assumptions and domain; a proof or
independent executable certificate; verified roles (now: per stratum); matched surface manipulations; a validated
parser and answer equivalence; source and licence traceability; independent blinded review; a cluster-independent
split; frozen intervention budget and scoring; balanced reported strata; all intended-to-evaluate records; an
unchanged Phase-I reproduction; one untouched external sample. No ``certified'' claim about natural-language
semantics rests solely on a symbolic free-symbol test.

\textbf{Benchmark-only scientific value test.} If the rebuilt dataset enables no comparison that GSM-DC,
GSM-Symbolic and existing physics benchmarks cannot support, benchmark size or formal verification is not the lead
novelty; the dataset is infrastructure for the causal study. A standalone dataset contribution needs demonstrated
measurement validity, distinctive physics relevance factors, independently validated external tasks, baseline
coverage and usable release tooling; cost and quality are compared, not only counts.

\textbf{Final Phase-B go/no-go (v2.0, kept for citation; REPLACED).} GO--MECHANISM if physically validated,
structure-generalisable, answer-relevant interference is reproducible and sufficiently powered; GO--MEASUREMENT if
a coherent output-likelihood or representation-invariance phenomenon survives but correctness and relevance claims
do not; PIVOT--GENERAL NUMERIC if analogous GSM-DC or physics-stripped tasks explain the effect;
REDESIGN--BENCHMARK if natural-language semantics or scoring cannot be trusted; STOP BROAD CLAIM if no meaningful
effect survives held-out validation.

\textbf{Correspondence for old citations} \tg{PROPOSAL} (the binding rules are in Section~\ref{sec:gates}):
\begin{itemize}
\item GO--MECHANISM $\rightarrow$ G4 establishes H3, H4 or H6-RD (or H6-extended) on the role-by-value family
(Section~\ref{sec:hypotheses}); G5 follows.
\item GO--MEASUREMENT $\rightarrow$ only presence, readout, format or invariance effects survive (H1, H2, or the
convertible-unit family); behavioural report and T1-type readout claims only.
\item PIVOT--GENERAL NUMERIC $\rightarrow$ the non-physics dependency control (G3, then G6) shows the same
behaviour; the result is reported as a general numeric or dependency phenomenon.
\item REDESIGN--BENCHMARK $\rightarrow$ G1 certificates or B4 validation fail; return to G1.
\item STOP BROAD CLAIM $\rightarrow$ H0 shown by equivalence (on the calibrated and the all-validated populations), H5 (cost only), or a G3 effect that G4 does not
replicate; reported as a null with its power, a cost finding, or ``not replicated'', never rescued by new
endpoints.
\end{itemize}

\gate{\textbf{Kept rule.} Gate outcomes do not license retroactive relabelling of confirmatory data. A structure
never moves from EXPLORATION into a confirmation role, and a cell with zero evaluable records is \NE{}, never a
zero effect.}

% ===== s04_framework =====
\section{The falsifiable theoretical framework}
\label{sec:framework}

This section is the construct specification for all work from Stage~15 onward. It replaces the construct
choice left open by the October edition (B2, Section~\ref{sec:phaseB}) with a decision, and it is the
reference that every later section uses for endpoints, strata and causal claims.

\subsection{Research question and scoped claim}
\label{sec:question}

\textbf{Primary question.} When a physics problem contains a quantity whose role in the target computation is
certified, does the model's answer depend on that quantity in the way the physics requires --- invariant when
the quantity is irrelevant, responsive by the certified amount when it is relevant --- beyond what is explained by
the numeral's presence, its value and plausibility, its unit, its position, the wording that marks its role and
generic perturbation of the text? If the dependence departs from the physics, which computation produces the
departure; and if it follows the physics, which computation implements the selection?

\textbf{Secondary question.} Does reasoning (thinking on versus off within the same weights, and the reasoning
budget) change how role is used, and is a quantity's cost paid in correctness, in length, or in failures to
finish?

\textbf{Scoped claim.} PhysMon proposes to distinguish \emph{numerical salience} from \emph{physical dependency}
through independently validated counterfactual manipulations of a quantity's physical role and value, and then to
test whether that distinction has an identifiable causal implementation in a language model. The absence of a
matching benchmark in our literature search (Section~\ref{sec:literature}) is not proof of novelty; the claim is
made against the closest named works: PRISM-Physics (physics formula DAGs and process-level scoring), GSM-DC
(symbolic dependency graphs with irrelevant nodes), Qu and Gomes (a mechanistic account of math distractor
failure), Owusu, Wiegreffe and Feldman (anchoring pathways), and Liu (positional readout copying). A PhysMon
result that one of these already explains is reported as such, not as a new phenomenon.

\subsection{What ``irrelevant'' means: four role strata}
\label{sec:strata}

Physical irrelevance is not one property. A quantity's role is one of four strata, each with its own certificate
and its own correct behaviour. The October edition's single ``non-governing cue'' and the first revision's
``in-system inert quantity'' both conflated the last three.

\begin{center}\small
\begin{tabularx}{\linewidth}{@{}P{2.6cm}LLL@{}}
\toprule
Stratum & Meaning & Certificate & Correct behaviour \\
\midrule
\textbf{R} relevant & The quantity is in the target's dependency structure with non-zero influence &
Solver recomputes the gold; certified response $y^*_R(z)$ over the declared range, separated from the irrelevant-role answer by more than the loosest tolerance &
The answer moves with $z$ as $y^*_R(z)$ prescribes \\
\textbf{D} structurally disconnected & The quantity belongs to a system or event with no path to the target in the typed dependency graph &
Graph disconnection of the quantity's node from the target; answer invariance over the declared range &
The quantity does not enter the target calculation; answer invariant \\
\textbf{C} algebraically cancelling & The quantity appears in the governing relations of the target system and may legitimately enter intermediate steps, but cancels from the final expression &
Symbolic cancellation ($\partial y^*/\partial z=0$ identically) plus numeric invariance; the intermediate relations that contain it are listed &
Intermediate use is \emph{correct}; only the final answer must be invariant \\
\textbf{K} conditionally irrelevant & The quantity's influence depends on a regime, boundary condition or assumption (e.g.\ a static-friction coefficient when the applied force is below the sliding threshold) &
The regime predicate holds at every instance value; invariance within the regime; the counter-regime value at which it would matter is documented &
The model must establish that the condition holds; the answer is invariant inside the regime \\
\bottomrule
\end{tabularx}
\end{center}

\textbf{Rules.} (1) A model's internal representation of a quantity, attention to it, or its use in an
intermediate step is never classified as an error merely because the final answer is invariant to that quantity;
for strata C and K such use can be correct physics. (2) Error claims are made only about outputs (the final
answer's behaviour across values, Section~\ref{sec:rolevalue}) or about intermediate steps whose wrongness is
certified against the graph. (3) Each stratum is a separate benchmark stratum with its own hypotheses
(Section~\ref{sec:hypotheses}); results are never pooled across strata. (4) Stratum D is the only stratum in
which ``the quantity should not enter the computation'' is the expectation.

\subsection{The role-by-value design (scientific centrepiece)}
\label{sec:rolevalue}

A single relevant/irrelevant twin cannot separate role from difficulty (the relevant twin may need more steps)
or from a lexical marker (``separate system'' means ``ignore''). The primary design therefore crosses role and
value within each structure.

For a quantity $q$ in structure $s$, with values $z_1,\dots,z_m$ ($m\ge3$ in every run, so that misuse curves are
identifiable; $m\ge5$ with held-out values for the dependency-curve test), role $r\in\{R, D, C, K\}$ where feasible, and at least two
independently validated phrasings $p$ of each role relation (at least one without marker words such as
``unrelated'' or ``separate''), the generator certifies $y^*_r(z)$ for every cell. Correct physics requires
\[
y^*_D(z_1)=y^*_D(z_2)\ (\text{and likewise for } C, K),\qquad
y^*_R(z_2)-y^*_R(z_1)=\Delta^*_R\neq 0 .
\]
Two questions are then asked at once: is the model invariant where the physics says it must be, and does it
respond by the correct amount where the quantity is relevant?

\textbf{Estimands defined on sampled answers.} Answers are a mixture of correct values, wrong values and
non-complete records, so no single ``answer'' per cell exists; the estimands are built to survive that.
\begin{enumerate}
\item \emph{Tracking} (primary, per sample): $T=1$ if the sample is \texttt{COMPLETE} and its value lies within the
declared tolerance of the certified $y^*_r(z)$ of its own cell, else $T=0$, over all attempted samples; the
complete-case version is a sensitivity analysis. Per structure, $\bar T_r$ averages over values, phrasings and
samples. \emph{Inappropriate dependence} in D (and C, K) is the tracking loss $\bar T_N-\bar T_D$ against the
token-aligned non-numeric sentence $N$ at the same position, so that generic distraction by an added sentence is
not counted as role failure (the loss against the no-cue prompt is reported beside it): a departure in
\emph{either} direction (adding, subtracting, copying or any other use of the quantity) that moves the answer
outside the tolerance lowers tracking, so this measure has no blind sign. \emph{Responsiveness} is $\bar T_R$.
\item \emph{Answer-mode shares} (primary for attribution, per sample): the share of samples whose value matches
each prespecified misuse curve $f_h(z)$ in D, C and K cells (\emph{misuse shares}), and the share of R samples whose
value matches the irrelevant-role answer $y^*_D(z)$, i.e.\ that ignore a quantity that matters (\emph{neglect
share}). A misuse share in a D, C or K cell is compared with the rate at which the non-numeric control $N$ at the
same $z$ produces the same value; a misuse curve is supported only when its share exceeds that rate at two or more
of at least three values. The neglect share cannot use $N$ or the no-cue prompt as comparator (their correct answer
\emph{is} $y^*_D(z)$); it is compared with a placebo match rate in the same R cell, the share of R samples equal to a
prespecified placebo value $2y^*_R(z)-y^*_D(z)$ (as far from the correct answer as the neglect value, on the other
side), certified to be separated from the gold and every foil; where that value is physically impossible (for
example negative for a positive-definite quantity), the generator uses a prespecified multiplicative offset of the
same size instead, so that the placebo is as plausible as the neglect value. The misuse set contains signed variants (add and
subtract the quantity; substitute it; copy it; for C and K the wrong law or wrong regime), each with the direction
of its prediction stated in advance.
\item \emph{Normalised response fraction} (descriptive only): $\lambda_r=[m_r(z_b)-m_r(z_a)]/[y^*_R(z_b)-y^*_R(z_a)]$ for
each value pair, with $m_r$ the Hodges--Lehmann location of the \texttt{COMPLETE} answers of the cell (robust to
isolated unit slips) and the denominator always taken from the R twin, so every role is scaled by the physically
correct response. Targets: $\lambda_R=1$; $\lambda_D=\lambda_C=\lambda_K=0$. The certificate requires
$|y^*_R(z_b)-y^*_R(z_a)|$ above the loosest tolerance, so the denominator is never near zero; signs and zero golds
need no logarithms.
\item \emph{Role-discrimination index} $\Lambda=\lambda_R-\lambda_D$ (target 1), a descriptive summary only: it
is blind to some failures (a model that answers correctly in R and subtracts the quantity in D has $\lambda_D=-1$ and
$\Lambda=2$) and conflates others (ignoring the quantity everywhere and using it everywhere both give $\Lambda=0$).
\end{enumerate}
\textbf{Multimodality guard.} A location summary of a two-mode answer distribution can describe a smooth response
the model never produces (answers at 0 and 100 give a location near 50). Every $\lambda_r$ is therefore reported
with the cell's answer-mode composition (shares of correct, each misuse mode, other and non-complete answers); when
two or more modes each hold at least 20\% of the cell's complete answers, $\lambda_r$ is flagged \texttt{MULTIMODAL}
and not interpreted. Neither $\lambda$ nor $\Lambda$ is ever the sole basis of a claim about physical reasoning.
A cell with no \texttt{COMPLETE} answers has no $m_r$; the structure is \NE{} for $\lambda$ and $\Lambda$ (it still
counts for tracking and answer-mode shares), and the share of such cells is reported by role. The earlier log-elasticity formulation is
withdrawn: it is undefined for zero or sign-changing golds and for mixed answer sets, and its interaction is not
sign-aligned across structures.

\textbf{Misuse hypotheses as response curves.} For each D, C or K cell the generator emits a prespecified set of
misuse functions $f_h(z)$ (copy the quantity; substitute it for the input of the same dimension; add it to a sum
of same-dimension terms; for C and K the most plausible wrong law or wrong regime that keeps it). A misuse
hypothesis is supported only when answers follow $f_h(z)$ across several values of $z$ while controls do not; a
single answer equal to $f_h(z_1)$ is not evidence of that misuse. The generator chooses values so that every pair of
misuse functions, and each misuse function and the certified response, differ by more than the loosest tolerance at
two or more of at least three values; pairs that cannot be separated are recorded as confounded and not
distinguished. Answer matching keeps \emph{all} hypotheses an
answer matches: \texttt{UNIQUE\_MATCH}($h$), \texttt{MULTI\_MATCH}($\{h\}$), \texttt{NO\_MATCH}; a reporting
taxonomy may collapse these, but the overlap is always shown.

\textbf{Difficulty and wording controls.} Value effects are estimated within role (difficulty is roughly constant
across $z$ inside a role); the relevant role's no-cue difficulty is recorded as a covariate. Every estimand is
estimated per phrasing. \emph{Agreement across phrasings}, required before a role claim, means: the
phrasing-by-role interaction on each of the G4 contrasts F1--F3 (Section~\ref{sec:statistics}) is inside its SESOI by two one-sided tests, or, where that test is
underpowered, each phrasing's estimate exceeds the SESOI in the same direction; otherwise the result is reported
per phrasing and called wording-dependent.

\textbf{Feasibility.} Not every structure admits every stratum. Infeasibility is recorded by the generator and
never imputed. A structure is \emph{complete} for the primary contrast only if it has R and D cells at three values
in two phrasings, with $|y^*_R(z)-y^*_D(z)|$ above the loosest tolerance and differing answer-token spans at every
value. Calibration gates count \emph{complete} role-by-value structures, not structures in
general (Section~\ref{sec:gates}).

\textbf{Adapted problems.} A counterfactual built on an external problem (PHYSICS, UGPhysics, PRISM-Physics or any
other source) is a new item: it needs its own certificate and validation; it does not inherit the source's.

\subsection{Supporting conditions}
\label{sec:conditions}

Kept from the October design and the Stage~14 cells, now subordinate to the role-by-value core: no cue;
style-matched non-numeric sentence; same-unit cue at three proximity levels kept as \emph{distinct} conditions
(exact match, near within the grading tolerance, tolerance-separated); incompatible unit with the same digits;
numeric metadata; convertible unit (the invariance/equivariance family, old Cue~C); relevant edit of an existing
input; role reversal (use on request); formula hint (a positive diagnostic: expected to help or to do nothing,
never a gate); an explicit relevance instruction (``identify and ignore irrelevant information''); cue position
(before or after the relevant inputs); and a small \emph{non-physics dependency control} (GSM-DC-style problems
with the same role manipulation) to tell a general numeric or dependency phenomenon from one that needs physical
semantics.

\subsection{Outcomes}
\label{sec:outcomes}

\begin{center}\small
\begin{tabularx}{\linewidth}{@{}P{3.6cm}LP{3.6cm}@{}}
\toprule
Outcome & Definition & Use \\
\midrule
O1 usable-answer success & Share of all attempted samples graded \texttt{CORRECT} under the declared evaluation contract & Operational reliability; no conditioning on evaluability \\
O2 semantic physics correctness & Whether the produced answer is physically correct, where ascertainable: includes answers the contract could not parse but a frozen recovery rule or blinded adjudication can read (e.g.\ a LaTeX-wrapped final line); otherwise \texttt{UNKNOWN} & Reasoning performance \\
O3 completion and grading status & \texttt{COMPLETE}, \texttt{BUDGET\_EXHAUSTED} (generation stopped at the token limit), \texttt{FORMAT\_INVALID} (no end-of-thinking marker or no final-answer line while the budget was \emph{not} exhausted), \texttt{UNPARSEABLE}, \texttt{MULTIPLE\_CANDIDATES}, \texttt{EMPTY} & Failure diagnosis; an outcome in its own right \\
O4 role-by-value response & Tracking $T$ and tracking loss; misuse and neglect shares across values with all matches kept; descriptive response fractions $\lambda_r$ and index $\Lambda$ under the multimodality guard (Section~\ref{sec:rolevalue}) & Primary scientific endpoint \\
O5 cost & Generated tokens (mean and median, never compared across the two), budget exhaustion, thinking-closed share & Length and diversion (H5) \\
O6 candidate log-odds & Defined only within a causal claim type of Section~\ref{sec:causal} & Mechanism metric; secondary behavioural diagnostic \\
\bottomrule
\end{tabularx}
\end{center}

\textbf{Rules.} A record that is not \texttt{COMPLETE} is never relabelled \texttt{WRONG}; it is ``not correct''
for O1 and \texttt{UNKNOWN} or adjudicated for O2. Truncation requires evidence that the budget was exhausted; a
missing end-of-thinking marker alone is \texttt{FORMAT\_INVALID}. Only text after the end-of-thinking marker is
graded by the contract. Units are converted with a dimension check, never stripped. Every record stores its
relative error to the gold. Three grading views are always reported: the declared last-printed-place rule
(primary), the unit-convertible rule, and relative error $\le 1\%$ with dimension-checked conversion (the field
convention of UGPhysics and PHYSICS). Cells whose cue lies within a view's tolerance of the gold stay in O1 and O2;
they are flagged \emph{source-ambiguous} under that view and excluded only from claims about the source of the
answer (copying versus computing). A stratified parser audit reports the parser's false-negative rate; no LLM
judge enters a primary analysis. A cell with zero evaluable records is \NE{}, never a zero effect.

\textbf{Retired:} the legacy signed teacher-forced label ($S_{lp}$, base~0, 0.5~nats) and any single-string
teacher-forced probability as an endpoint. \textbf{Diagnostics only:} the anchor-matched presence boost and
$R_{alt}$ of Stage~14.

\textbf{Sampling.} Repeated samples per cell are required; greedy decoding gives one deterministic answer. The
Stage~14 reasoning runner for Qwen2.5 and Llama is greedy and unbatched, so a new contract version samples at each
model's published defaults (recorded per record) and batches; Qwen3 thinking uses temperature 0.6, top-$p$ 0.95,
top-$k$ 20 as its model card requires. Greedy Stage~14 records are never pooled with sampled ones.

\subsection{Populations}
\label{sec:populations}

Every result names its population; none is silently replaced by another.

\begin{center}\small
\begin{tabularx}{\linewidth}{@{}P{4.2cm}L@{}}
\toprule
Population & Purpose \\
\midrule
All validated tasks & General behavioural and robustness conclusions; always reported \\
Calibrated intermediate-difficulty tasks (per model and budget) & Efficient discovery of correctness-changing effects; selection on samples never reused for estimation \\
Tasks solved reliably without a cue & Induced errors, repair, and the mechanism of \emph{correct} selection \\
\bottomrule
\end{tabularx}
\end{center}

Comparisons between models, or between thinking on and off, use one common item set inside every compared
configuration's band.

\subsection{Hypotheses and their distinguishing predictions}
\label{sec:hypotheses}

All are legitimate outcomes; more than one can hold at once (numeric copying and a selection failure are not
exclusive). They refine the October families: H-N becomes H1 and H2, H-P becomes H3, H4 and H6, H-T is the
retired string construct plus the position and format controls, H-U is the convertible-unit family.

\begin{center}\small
\begin{tabularx}{\linewidth}{@{}P{0.6cm}P{3.6cm}LL@{}}
\toprule
Id & Hypothesis & Predicts (in O4 unless stated) & Evidence against \\
\midrule
H0 & No role or presence effect beyond generic perturbation & D tracking equal to tracking under the non-numeric sentence; no misuse share above its rate under $N$ and no neglect share above its placebo rate; numeric cells equal the non-numeric sentence & Any reliable departure; equivalence must be shown, not assumed \\
H1 & Numeral presence (F-copy, F-answer prior) & Answers follow the copy curve in every role, stratum and unit & Copy tracking absent in D or confined to same-unit cells \\
H2 & Plausibility and unit (F-match, F-unit) & Departures only for same-unit or convertible plausible values; vanish under plausibility matching & Departures with incompatible units or implausible values \\
H3 & Relevance-selection failure (F-role) & D: tracking loss with answers following a prespecified misuse curve (substitute, add or subtract, in its predicted direction). C: final answer follows the wrong-law curve that keeps the quantity. K: final answer follows the out-of-regime law. Moderators (marking, stratum, position) are secondary & Invariance in D, C and K while R responds correctly \\
H4 & Readout contamination & Correct value in the trace, final answer moved toward the quantity; little change under the relevance instruction & Departures already present in the trace's computation \\
H5 & Computational diversion & O5 rises (length, budget exhaustion) while O1 and O4 are equivalent to no-cue & O1 or O4 departures \\
H6 & Correct selection, in two scopes. \emph{H6-RD}: the three G4 contrasts F1 (D tracking loss against $N$), F2 (R neglect share against its placebo rate) and F3 (role $\times$ value interaction on tracking) each shown equivalent to zero within its SESOI (two one-sided tests), with R tracking at or above the G1 criterion, across values and phrasings. \emph{H6-extended}: H6-RD plus the same for C and K (intermediate use in C and K allowed) & Departures in any stratum; an underpowered null is INCONCLUSIVE, not H6 \\
\bottomrule
\end{tabularx}
\end{center}

\subsection{Causal identification standard}
\label{sec:causal}

The withdrawn Phase-I paper failed because its intervention metric did not isolate the claimed function: a
knockout that raised the gold string on base and sensitive renderings alike was read as cue processing. Every
mechanism experiment therefore declares which of three claims it tests, and uses the metric of that claim.

\begin{center}\small
\begin{tabularx}{\linewidth}{@{}P{3.2cm}LL@{}}
\toprule
Claim type & What is held fixed & Metric \\
\midrule
T1 readout given a fixed reasoning state & Only between prompts that share a gold: a D, C or K cell and its token-aligned non-numeric control at the same $z$. One explicitly identified prefix, the same under both prompts, its type declared in G1 per experiment: (i) the gold trace rendered from the typed graph, or (ii) the model's own trace under the control prompt, taken for every sampled trace regardless of its answer (never selected on the behaviour being explained); intervention only at answer positions & $\log p(y^*(z))-\log p(f_h(z))$: the certified answer against a predeclared misuse alternative, both scored over their full value spans under the declared answer-format contract. The prefix's per-token NLL under the recipient prompt is recorded as a declared compatibility stratum; results are reported for all prefixes and by stratum, and no prefix is excluded \\
T2 intermediate computation & A teacher-forced prefix up to the step where the quantity enters; certified intermediate values from the graph & Log-odds or decoding of the certified intermediate value against the prespecified misuse value; for strata C and K intermediate use is expected, not penalised \\
T3 total effect & Nothing downstream: intervene upstream (at the quantity's tokens or a candidate component) and let the model generate & O1, O2 and O4 over $N$ on-policy resampled traces per condition \\
\bottomrule
\end{tabularx}
\end{center}

Only T2 or T3, with the controls below and the role interaction defined next, can support a claim about a
physical-relevance computation; T1 supports a claim about answer readout only. Comparisons between an R twin and a D twin, whose golds differ, are made only
under T2 or T3, never by teacher-forcing one role's trace under the other role's prompt (that trace would carry a
physically wrong value). Candidate foils are fixed by a prespecified physical error hypothesis, never chosen as the
foil the model scores highest; a candidate pair whose values lie within tolerance or share their answer-token span
is T1-infeasible and recorded.

\textbf{Two interaction estimands.} With $\Delta_h m_X=m(X,\text{intervened})-m(X,\text{clean})$ the effect of
intervention $h$ on prompt $X$:
\[
I_{\text{presence}}=\Delta_h m_D-\Delta_h m_N,\qquad I_{\text{role}}=\Delta_h m_D-\Delta_h m_R .
\]
$N$ is the token-aligned non-numeric sentence (so the intervened positions exist in both prompts; the no-cue prompt
cannot serve, since it has no tokens at the quantity's position); $R$ is the relevant-role twin with the
\emph{same numeral, unit and position} at the same $z$. $I_{\text{presence}}$ can use any declared metric, since
$D$ and $N$ share a gold. $I_{\text{role}}$ must use a behaviourally comparable outcome, because the two worlds
have different certified answers: tracking against each world's own certified answer, or the answer-mode shares,
under T3 (or T2 with each world's certified intermediate values); raw gold-string log-odds are never subtracted
across worlds. The clean effects beside each ($m(D)-m(N)$, $m(D)-m(R)$) are always reported. A component with
$I_{\text{presence}}\neq0$ but $I_{\text{role}}\approx0$ is a numeral-cue component; a claim of physical-role
routing requires $I_{\text{role}}\neq0$, holding across values and phrasings, with the controls below. A large equal
effect on all prompts marks a general answer component (the head-11 lesson). $I_h$ elsewhere in this document
means $I_{\text{presence}}$ or $I_{\text{role}}$ as declared.

\textbf{Protocol checklist} (each item cites the source of the practice):
\begin{enumerate}
\item Noising (necessity) and denoising (sufficiency) both, reported separately; they are not complements (Heimersheim and Nanda).
\item Counterfactuals in distribution: role-by-value minimal pairs that change only the role relation or only the value, token-aligned; several corruptions per item; no Gaussian noise; zero ablation never as evidence (Zhang and Nanda).
\item Logit difference or log-odds with each candidate shown; never probability; never normalised recovery (values above one come from negative components, not cancellation).
\item Resample or interchange ablation from matched counterfactuals as primary, mean ablation as a sensitivity; conclusions reported under both (Miller, Chughtai and Saunders; Li and Janson).
\item Direct and total effects, with a frozen-final-LayerNorm control for self-repair (Rushing and Nanda; McGrath et al.).
\item Divergence check: patched activations against the natural distribution at the site, and KL on neutral prompts (Grant et al.).
\item Attribution patching only as a screen; every reported component confirmed by exact patching (Kram\'ar et al.).
\item Selection stability under bootstrap subsampling; discovery and test on disjoint structures (Certified Circuits; review issue~6).
\item Matched random components as nulls; family-wise correction over the components tested.
\item Rescue begins \emph{within structure, value-matched} (under T2 or T3): transfer the candidate state between the R and D twins of one structure at the same $z$, where the principal difference is the role relation; cross-structure donors with independent matching are a later, stronger generalisation test. A transfer must change the computation as predicted while preserving unrelated information (value, unit, answer).
\item A variable claim needs causal abstraction: a role variable localised by DAS (pyvene) with interchange-intervention accuracy on held-out pairs against neuron and random-subspace baselines (MIB), plus the subspace-illusion check (Makelov et al.). A subspace from which role can be decoded is \emph{diagnostic}; a subspace whose interchange changes behaviour as predicted is \emph{causally sufficient}; only the second supports a computation claim. Binding work shows several mechanisms can combine (Gur-Arieh, Geva and Geiger), so a single isolated role feature is not assumed.
\item Reasoning models: T1 on a fixed trace, T3 with on-policy resampling (Thought Branches); never a hand-edited trace as primary evidence; no conclusion from one sample.
\item Probes are diagnostic only (Hidden Error Awareness).
\item Held-out prediction: the sign and size of $I_h$ on held-out structures, values and the second model are committed before the test.
\end{enumerate}

\subsection{Claim hierarchy and evidence gates}
\label{sec:claims}

The October claim levels are kept, with their content restated for the revised construct (they are written
CL1--CL6 here, to avoid confusion with stratum C, review checks C1--C10 and inventory items C01--C18; October
citations of ``C1--C6'' mean these levels).
\begin{itemize}
\item \textbf{CL1 (benchmark):} certified role strata and responses plus audited natural-language fidelity
(B4), independence, coverage and well-specified exclusions.
\item \textbf{CL2 (behavioural):} a reproducible role-by-value pattern (O4) on unseen template clusters (G4),
reported across endpoints O1--O5 rather than one answer-string likelihood.
\item \textbf{CL3 (monitor, G7):} correctly scoped single-input monitoring with paired incremental value beyond
deployable baselines and enough precision to resolve deployment trade-offs.
\item \textbf{CL4 (causal localisation):} several valid intervention families identify specific components under
T2 or T3 with the interaction estimand, replicated on held-out structures and not explained by model damage.
\item \textbf{CL5 (mechanistic explanation):} a falsifiable functional account with specific necessity,
controlled sufficiency, upstream/downstream path tests, causal abstraction and held-out prediction
(Section~\ref{sec:causal}); generalisation beyond the selection panel (G6).
\item \textbf{CL6 (repair, G7):} a selective, generalising intervention that beats appropriate non-mechanistic
alternatives on a documented utility and cost frontier.
\end{itemize}
\textbf{Promotion rule (kept).} A claim advances one level only when the corresponding tests pass; a highly
significant CL3 monitor does not establish CL4 or CL5; a T1 readout result supports at most CL4 for readout. The next
paper needs either a convincing CL5, or an exceptionally well-validated CL1--CL2 finding with novelty beyond the
literature of Section~\ref{sec:literature} (Track~E), preferably with CL6 if feasible. ``Shortcut reliance'' is
reserved for evidence that the computation improperly uses the quantity (H3 or H4 established under T2 or T3). The
repository's forbidden-phrase check stays on.

% ===== s05a_gates =====
\section{Ordered execution: the critical path}
\label{sec:execution}

\subsection{Gates G0--G7}
\label{sec:gates}

Each gate names its deliverable, its passing condition and what happens otherwise; each is closed by a committed
gate report, never by declaration. Section~\ref{sec:blocks} maps the October blocks P0--P11 onto these gates.

\begin{center}\small
\begin{tabularx}{\linewidth}{@{}P{3.4cm}LP{4.6cm}@{}}
\toprule
Gate & Key deliverable & Decision \\
\midrule
G0 integrity closure & Scoring, provenance, comparator inventory with oracle/deployable tags, claim registry, review-derived checks & Can future results be trusted and reproduced? \\
G1 physics construct & Strata taxonomy R/D/C/K with certificates; role-by-value generator with phrasings and misuse curves; frozen outcomes and analysis & Are the intended interventions physically well defined? \\
G2 measurement and calibration & Grading contract, sampling baseline and self-noise floor, design-validation pilot, model and difficulty regime & Can the intended behaviour be measured with enough precision? \\
G3 behavioural discovery & Core role-by-value factorial plus minimal controls; hypothesis scorecard; non-physics control & Which phenomenon, if any, deserves mechanism work? \\
G4 independent behavioural confirmation & New validated structural clusters; frozen analysis of the role-by-value family & Is the effect real beyond development data? \\
G5 mechanism & Function hypothesis; T1/T2/T3 experiments with the interaction estimand, rescue and held-out prediction & What computation explains the behaviour? \\
G6 external and functional confirmation & Independently sourced physics items (re-certified), non-physics transfer, second model & Does the explanation generalise as claimed? \\
G7 practical use (optional) & Monitoring or selective repair against cheap alternatives & Does the explanation support a useful intervention? \\
\bottomrule
\end{tabularx}
\end{center}

\subsubsection*{G0. Integrity closure (CPU, plus a short GPU parity check)}
\begin{enumerate}
\item Status sync: the Stage~14 status and roadmap through Stage~14.20; ledger rows and decision-log entries for job
435605 and the cancellation of 435606; the stale journal-index rows.
\item The seven review-derived additions (\texttt{docs/reviews/iclr2027\_submission50608/analysis.md}, section~5):
(a) cross-instruction stability of any label on the pilot families (CPU where records allow, otherwise \NE{});
(b) the formula-hint cell; (c) the interaction rule for every intervention (Section~\ref{sec:causal}); (d)
oracle/deployable tags on every comparator, propagated to the Wave~2 record and the claim ledger; (e) a mechanical
claim-wording check between any manuscript draft and the claim registry; (f) a provenance-consistency check that
\texttt{created\_by} and validation fields match any text about authorship and validation; (g) a dated note that
the 0.5-nat threshold was chosen after the Stage~4 results.
\item P0, narrowed (open decision D1): full regeneration of Phase-I values is off the critical path because the
paper is withdrawn and no new claim rests on its numbers. HF-versus-TransformerLens parity already passed for
Qwen2.5-7B-Instruct and Llama-3.1-8B-Instruct on the smoke prompt (job 399217: top-1 equal, cosine 0.99998 and
0.99999); G0 adds parity at the hook sites and trace lengths G5 will use.
\item Phase-I closure note: which statements survive (Section~\ref{sec:ledger}).
\item Qwen3-32B: identify its revision and write a manifest; earlier runs (417416, 425243, 435605) stay
\texttt{model\_commit NOT\_RECORDED} unless the file hashes match a published revision, recorded as an
after-the-fact identification.
\end{enumerate}
\gate{\textbf{Gate G0.} Passes when the items are committed. If extended parity fails for a model or backend: no
mechanism work on it; behaviour can proceed.}

\subsubsection*{G1. Physics construct (CPU)}
Freeze before any G2 output: the strata taxonomy and certificates (Section~\ref{sec:strata}); the role-by-value
generator with at least two validated phrasings per role relation, misuse functions $f_h(z)$ and feasibility
records (Section~\ref{sec:rolevalue}); proximity conditions regenerated so that exact-match, near-within-tolerance
and tolerance-separated cues are distinct conditions; the outcome definitions, grading views and O3 status
categories (Section~\ref{sec:outcomes}); a sampled, batched reasoning-contract version (prospective only; the
greedy Stage~14 records are never pooled with it); analysis and power-simulation scripts tested on synthetic
records; the scorecard criteria of G3 with numeric thresholds; the minimum R-tracking criterion checked in G2;
the T1 prefix type per experiment and its compatibility strata; equivalence margins with their scientific
justification (Section~\ref{sec:statistics}); experiment cards; the LLM review panel (models chosen under D2, prompts
frozen, qualified on the extended planted-defect bank, Section~\ref{sec:llmpanel}); the M0 pilot structures; and
the dataset admission standard and data-role allocation of Section~\ref{sec:datasetplan}.
\gate{\textbf{Gate G1.} Passes when certificates, mutation tests and analysis tests pass, the LLM panel is
qualified and the cards are committed. Nothing is run before. The M250 build (Section~\ref{sec:datasetplan}) follows
the G2 design-validation pilot and precedes G2 calibration.}

\subsubsection*{G2. Measurement and calibration (GPU, small)}
\begin{enumerate}
\item \textbf{Design-validation pilot} (8--12 structures, core role-by-value cells, a few samples): are the
contrasts unambiguous, feasible, correctly certified and correctly graded? Structures are checked by the best
available reviewer (the qualified LLM panel once it exists, plus the user's walkthrough; each labelled as such).
Defects found here change the generator, never the analysis of later data.
\item \textbf{Calibration} per candidate model and budget: $k_{sel}$ no-cue samples to select items, then $k_{est}$
fresh samples on the selected items to estimate the baseline and the self-flip floor (never the selection
samples: the 14.20 preliminary contrast had a baseline of 1.000 by construction, J14.19-014); on a subset of the
selected items, a few samples of the R and D cells at three values in one phrasing and of the non-numeric sentence,
so that the variance of the actual primary estimands (tracking loss, misuse and neglect shares) is estimated, not only a
no-cue variance; tokens, budget exhaustion and format failures.
\item \textbf{R-tracking check.} A minimum tracking share on R cells, fixed numerically in G1, is checked here: a
model or budget below it does not use added relevant quantities reliably enough for any role reading, and is
changed or dropped at this gate, never after G3.
\item \textbf{Primary-model rule}, frozen at the end of G2 and before any G3 output: the primary model is the first
in the prespecified order Qwen2.5-7B-Instruct, Llama-3.1-8B-Instruct, Qwen3-8B (if downloaded, D3) that meets the
G2 criteria (complete structures in band, R-tracking criterion, extended parity in the study dtype for G5
instrumentation); the next model that meets them is the second model. The choice never depends on observed cue
effects.
\end{enumerate}
\gate{\textbf{Gate G2.} Passes when enough \emph{complete} role-by-value structures (Section~\ref{sec:rolevalue}:
R and D cells at three values in two phrasings) fall in the band to meet the power target of
Section~\ref{sec:statistics} (planning: at least 24), the R-tracking check passes, and the cost fits
Section~\ref{sec:compute}. Otherwise: add harder or easier variants or change the budget, and repeat G2; a model
that cannot be brought into the band is reported on the all-validated population only and its band-based discovery
is not run. The band and the R criterion are not changed after seeing G2 results.}

\subsubsection*{G3. Behavioural discovery (GPU, staged)}
\begin{enumerate}
\item \textbf{Core run.} 24--40 complete calibrated structures from the 130-structure EXPLORATION pool of M250 (Section~\ref{sec:datasetplan}); core conditions: no cue; R and D roles at three
values each in two phrasings; the non-numeric sentence; plus the small non-physics dependency control set. Sampling
may be adaptive only by a rule fixed in G1. Results are reported on the calibrated population and, for the same
cells where run, on all validated structures.
\item \textbf{Targeted explanatory ablations}, added only when the core run shows something to explain: strata C
and K, marking and system membership, unit and proximity, position, relevance instruction, formula hint,
convertible units.
\item \textbf{Hypothesis scorecard.} For each of H0--H6 the prespecified evidence for and against (Section~\ref{sec:hypotheses})
is estimated with intervals and given one status: SUPPORTED, CONTRADICTED, INCONCLUSIVE or \NE{}. Several
hypotheses may be supported at once; hypotheses are not exclusive rows and are not read by first match. An
interval that contains zero but is not inside the equivalence margin is INCONCLUSIVE, never H0; H0 and H6 are
SUPPORTED only by two one-sided tests, which G3's size usually cannot power, so in G3 they are normally
INCONCLUSIVE. A null on the small non-physics control is INCONCLUSIVE unless it is powered, and is not evidence of
physics-specificity.
\item \textbf{Optional exploratory mechanism pilot} on the primary 7B model, if it helps discriminate hypotheses;
its results are never confirmatory and its site selection is separated from G5's held-out tests.
\end{enumerate}
\gate{\textbf{Gate G3.} The G4 question is fixed now, not chosen from G3's largest estimate: G4 confirms the
role-by-value family (Section~\ref{sec:statistics}) whatever G3 shows, provided G2 passed. G3 can add at most one
secondary family to G4, chosen by a rule fixed in G1 and labelled as chosen on EXPLORATION data. If R tracking in
G3 falls below the G1 criterion, the G3 role result is reported as uninterpretable for role (itself a finding);
items and band are not changed on the same data. A formula hint that does not help is a diagnostic, not a gate.}

\subsubsection*{G4. Independent behavioural confirmation (GPU, medium)}
Structures from the sealed INTERNAL\_CONFIRMATION pool of M250 and M500 (270 structures, Section~\ref{sec:datasetplan}), in template clusters no earlier run has touched, opened in the number the power simulation requires (shadow set held back),
validated by the qualified LLM panel (Section~\ref{sec:llmpanel}; human B4 review where available) or, failing that,
labelled EXPLORATORY-REPLICATION. Preregistered: the role-by-value family with its
estimands, decision rules and sample size from the power simulation; the primary and second models fixed by the
G2 primary-model rule, the second by fixed sequence; two declared final-answer formats on a random half of items, each with its effect estimate and
interval, and a stated rule for how much format dependence changes the interpretation (planning: the reading
changes if the format-specific effects differ by more than a factor of two or have opposite signs with both
intervals excluding zero). Opened once.
\gate{\textbf{Gate G4.} Passes when the primary family meets its rule on the primary model. If not, the G3 pattern
is reported as not replicated.}

\subsubsection*{G5. Mechanism (GPU)}
Gated on G4 establishing H3 or H4, or H6-RD (the mechanism of correct selection between relevant and structurally
disconnected quantities, when the model selects correctly and errors are rare; it needs the equivalence tests that
only G4's size can power). A mechanism-of-correctness claim covers only the strata G4 confirmed: H6-extended (C and K
as well) needs those strata in G4's preregistered family; under H1/H2 only a narrow study of
unit or number handling, framed against Liu and Owusu et al. The exploratory mechanism pilot of G3 is the only
mechanism work before G4. Every experiment declares claim type T1, T2 or T3 and
follows Section~\ref{sec:causal}. Starts on the primary and second models fixed by the G2 rule (planning order Qwen2.5-7B-Instruct, Llama-3.1-8B-Instruct; both listed as fully verified by TransformerLens~3.4;
locally, HF-versus-TL parity has passed only on one smoke prompt in float16 with parameters on CPU, job 399217, so
the G0 extended parity in the study dtype is required first); Qwen3-8B with thinking afterwards (if downloaded, D3); Qwen3-32B only for a frozen-function
transfer test.
\gate{\textbf{Gate G5.} A role-computation claim needs T2 or T3 evidence with $I_{\text{role}}\neq0$ (not only
$I_{\text{presence}}$), both directions, rescue within structure, causal abstraction for any variable claim, and a correct held-out
prediction. Otherwise the result is reported at the level it reached (Section~\ref{sec:claims}).}

\subsubsection*{G6. External and functional confirmation}
Independently sourced physics problems (EPC; adaptations of PHYSICS, UGPhysics or PRISM-Physics where feasible,
each re-certified and re-validated as a new item); non-physics transfer (GSM-DC or GSM-Symbolic NoOp with the same
role-by-value logic where it maps); a second model family; frozen function, not frozen component index. Required
for any broad generalisation claim.

\subsubsection*{G7. Practical use (optional)}
Monitoring (the Phase-I monitor under cluster-held-out folds is a side track, D6) and selective repair against a
cheap ladder (relevance instruction, self-consistency, longer budget), with strength-versus-damage curves.

% ===== s05b_blocks =====
\subsection{Dataset plan: 250 structures, then 500}
\label{sec:datasetplan}

\textbf{Target (user decision, 11 November 2026).} DECOY-v2 dataset~v2 is built to \textbf{250 independent
structures} first and expanded to \textbf{500 structures} after G3 and before G4. The count is in independent
structures, the unit of inference (Section~\ref{sec:statistics}): each structure is its own template cluster and
its own symbolic program, and yields several numeric values, role-by-value cells and phrasings, so 250 structures
mean roughly a thousand or more problem instances. This replaces the open-ended ``only what the gates need'' sizing
of the earlier revision and the October edition's tentative 300--600 band.

\textbf{Starting point} \tg{V}. Dataset~v1 (Stage~14.19, \texttt{data/v2/dataset/v1}) has 80 structures and 260
instances (50 structures with four instances, 30 with two; 80 distinct template clusters and programs; 11 domains;
levels 104 high-school, 124 introductory-university, 32 upper-university instances). All are agent-verified and
\texttt{UNREVIEWED}, carry the eleven Stage~14 distractor cells (2,764 of 2,860 feasible) and no role-by-value
cells. Models have already been run on them, so they can only ever be EXPLORATION data.

\textbf{Admission standard (``extreme high quality'').} A structure counts towards 250 or 500 only when all of
these hold, and each is recorded in its provenance:
\begin{enumerate}
\item typed dependency graph with a solver certificate for every cell (strata R and D required; C and K where
feasible, Section~\ref{sec:strata}), certified responses $y^*_r(z)$ at three or more values, separation of
$y^*_R$ from $y^*_D$ and of every misuse curve and placebo value beyond the loosest tolerance;
\item two or more independently validated phrasings per role relation, one without marker words;
\item mutation tests passed (P3.3), including broken cancellation and out-of-regime mutations for C and K;
\item acceptance by the qualified LLM panel (Section~\ref{sec:llmpanel}) on every phrasing, with no unresolved
high-severity flag;
\item a blind naturalness check and a grading dry run (the gold answers parse and grade as \texttt{CORRECT} under
every grading view);
\item no overlap with any other structure of dataset~v1, v2 or DECOY-v1 under the P3.5 detectors; a new template
cluster and a new symbolic program.
\end{enumerate}
A structure that fails is replaced, never edited after any model output on it exists. Planning assumes about 1.3 to
1.5 generated structures per admitted one.

\textbf{Composition targets} \tg{PROPOSAL}, checked at each milestone and reported, not forced by editing items:
at least 10 domains with no domain above 25\% of structures; at least 20\% upper-university structures and
multi-step variants (B3.1 strata), so that models are not at ceiling; stratum C feasible in at least 30\% and K in
at least 20\% of structures, so that H6-extended and the C/K forms of H3 can be powered; difficulty spread recorded
before any model run.

\textbf{Milestones and data roles} (assigned by lineage, before any outcome is seen; Section~\ref{sec:status},
rule 1):
\begin{center}\small
\begin{tabularx}{\linewidth}{@{}P{2.3cm}P{3.0cm}P{2.6cm}L@{}}
\toprule
Milestone & When & Structures & Role and use \\
\midrule
M0 & G1 & 8--12 & EXPLORATION: design-validation pilot of G2 (may include upgraded v1 structures) \\
M250 & After the G2 design-validation pilot has confirmed the generator and grading, before G2 calibration & 250 in total & 130 EXPLORATION (the 80 v1 structures upgraded with role-by-value cells and re-validated, plus 50 new) for the G2 pilot, calibration and G3 discovery; 120 new INTERNAL\_CONFIRMATION, sealed until G4 \\
M500 & After G3, before G4 opens & 500 in total & 250 more new structures: 150 INTERNAL\_CONFIRMATION (G4 pool 270 in total), 50 shadow set (inaccessible to all tuning scripts and agents), 50 ADVERSARIAL\_CHALLENGE and G6 reserve \\
\bottomrule
\end{tabularx}
\end{center}
G4 opens only as many sealed structures as its power simulation requires (planning at least 50 complete and
in-band per model, Section~\ref{sec:statistics}), chosen by a rule fixed before opening; the rest stay sealed for
G6. Because calibration bands differ by model, a large sealed pool is what lets two models and the C/K strata reach
their targets. External (EPC) items are additional and are not counted in the 500. If generation or validation
falls short at a milestone, the shortfall is reported and the gate proceeds only if its own power target is met;
admission standards are never relaxed to reach a count.

\textbf{Effort.} Generation and certification are agent and CPU work. LLM-panel validation costs about 2--3 judges
$\times$ 2 passes per structure and phrasing (about 3,000--6,000 judge runs for 500 structures; local GPU-hours or
an API budget, decided with D2). The later human accuracy study samples the panel's verdicts, not all 500.

\subsection{The original critical-path blocks P0--P11 under the revised gates}
\label{sec:blocks}

The October edition ordered its work as twelve critical-path blocks, P0 to P11. They are kept here, with their
identifiers, questions, sub-items, hard gates and failure responses, because they remain the most detailed
statement of what each kind of evidence requires. They no longer set the order of execution: that is fixed by the
gates G0--G7 (Section~\ref{sec:gates}), and the construct, strata, endpoints, hypotheses and causal standard are
those of Section~\ref{sec:framework}. Where a block's text conflicts with either, Sections~\ref{sec:framework}
and~\ref{sec:gates} win; the conflicts known at this revision are resolved in the block text below and listed in
its \emph{Disposition} line. Block text that names the legacy constructs ($S_{lp}$, Cue~A/B/C, the four-rendering
family, L16~H11) is rewritten in terms of the role strata R/D/C/K, the role-by-value design and the outcomes
O1--O6 wherever the requirement survives.

Three rules of the October dependency paragraph stay in force. (1) Each gate is about correctness, not dates; the
plan is not calendar-driven. (2) A claim ledger is updated continuously (P11) but final claims are frozen only when
the branches they need have closed. (3) Mechanism work does not start before trustworthy partitions and a frozen
behavioural contrast exist; in the revised plan that means after G4, for H3, H4 or H6-RD (Section~\ref{sec:gates}). The October
allowance for ``exploratory H11 checks'' is replaced by a narrower one: a small exploratory mechanistic pilot may
run during G3, marked exploratory and never confirmatory, and no claim may rest on it.

\begin{center}\small
\begin{tabularx}{\linewidth}{@{}P{0.9cm}P{4.1cm}P{1.7cm}L@{}}
\toprule
Block & October title (short) & Gate(s) & Disposition at this revision \\
\midrule
P0 & Evidence reproduction, provenance, code audit & G0 & Narrowed (open decision D1): full regeneration of Phase-I values off the critical path; parity, registers and a conditional spot check kept \\
P1 & Measurement validity of the sensitivity construct & G0, G2 & Largely done in Stage~14; likelihood findings are diagnostics; open sub-items optional \\
P2 & Scale the DECOY-v2 design & G1 (M250), after G3 (M500) & Revised: 250 structures, then 500 before G4 (Section~\ref{sec:datasetplan}) \\
P3 & Expert validation, parser, leakage firewall & G1, G2, G4 & Kept; parser audit extended to LaTeX final lines and the O3 status categories \\
P4 & Behavioural mechanism discrimination and power & G3, G4 & Central; recast as the role-by-value design \\
P5 & Monitoring study & G7 & Demoted to an optional side track; comparators carry oracle/deployable tags \\
P6 & Re-localisation under valid causal metrics & G5 & Rules kept; the H11 lesion (card X07) retired as stated; components found anew under the interaction estimand \\
P7 & What the components compute & G5 & Kept; F-hypotheses mapped onto H1--H3 and the position control; H4 and H6 added \\
P8 & Upstream/downstream computation & G5 & Kept; gated on G4 \\
P9 & Cross-model, cross-domain, post-training generalisation & G5, G6 & Kept; gated on G4; freeze the function, not the head index \\
P10 & Mechanism-guided repair & G7 & Kept as optional practical use; cheap ladder first \\
P11 & Confirmatory closure and claim freeze & G6, all & Kept; external cohort from EPC/external physics; venue after the result \\
\bottomrule
\end{tabularx}
\end{center}

Every \textbf{Status} line below is dated 11~November~2026. Nothing done in Stage~14 is confirmatory: all of it is
data role EXPLORATION and not paper-eligible. Numbers marked \tg{SUBMITTED} are from the withdrawn Phase-I
manuscript; planning figures are \tg{PROPOSAL} and are choices, not results.

%----------------------------------------------------------------------
\subsubsection{P0. Independent evidence reproduction, release provenance and code audit}
\label{sec:blockP0}

\textbf{Question (October).} Which of Phase~I's reported numbers follow from the exact datasets, versions, model
weights and scoring code? A corrupted label, leaked fold or misaligned hook invalidates downstream mechanistic
interpretation.

\textbf{Revised scope (G0).} The Phase-I paper was withdrawn on 11~November~2026, and no new claim rests on its
numbers. Full bottom-up regeneration of Phase-I values is therefore off the critical path (open decision D1); if the
full P0 is kept, it runs as a side track and blocks none of G1--G4. What remains on the critical path is the part
that protects \emph{new} evidence: hooks that touch the tensors and positions they claim, and registers that tie
every number to bytes.

\begin{description}[leftmargin=1.2em,style=nextline]
\item[P0.1 Immutable artifact and claim registry (kept).] Every family or structure instance gets an identifier;
hashes are frozen for the symbolic template, rendered text, response prefix, reference answer, model revision,
tokenizer, generation configuration, scoring script, fold, intervention site, donor and target identities,
scheduler job, git SHA and seed. A machine-readable map links each \texttt{claim\_id} to artifact paths, generator
code, checksum, eligibility rule and source table or figure, and marks exploration, internal confirmation and
untouched confirmation. Missing, cancelled and zero-row runs are registered.
\item[P0.2 Bottom-up reproduction of the submission tables (off the critical path).] Regenerating the
$155\to140\to135$ inclusion flow, per-cue rates, label overlaps, AUROC/AUPRC, Brier, ECE, FNR@FPR, out-of-fold
predictions, bootstrap distributions, the layer-16 head ranking, causal per-family numerators and denominators and
the 15-family transfer membership \tg{SUBMITTED} is part of the D1 side track only.
\item[P0.3 End-to-end re-run (narrowed).] Two pieces stay. (a) \emph{Extended parity}: Hugging Face against
TransformerLens at the hook sites G5 will use, on traces of the planned length, for every model that enters
mechanism work. (b) A forward-pass spot check of 10 legacy families' stored log-probabilities, stratified as in
October (including \texttt{CM\_B\_UM\_059} and the adversarial edge cases), run \emph{only if} a Phase-I number is
cited again.
\item[P0.4 Leakage and cohort-selection code inspection (rules carried forward).] The Phase-I code audit belongs
to the side track. Its rules apply to all new work: scalers, features, hyperparameters, thresholds, layers,
calibrators, stacking and residualisation inside the nested training loop; no candidate selected on future
confirmation data; no model-generated parse used as an inclusion filter; no structurally related template across
a split; one observation traced by hand by a second investigator before any predictor result is reported.
\item[P0.5 Provenance discrepancies (logged, not resolved).] The renamed-donor control (zero evaluable rows in the
manuscript, ``completed with zero recovery'' in an older record) and the exploratory Llama head-2 sweep
($\approx1.85$, incomplete provenance) \tg{SUBMITTED} stay in the discrepancy log with that status; neither is
upgraded or used without a complete clean replication.
\end{description}

\textbf{Deliverables (kept where in scope).} Artifact and job registers, claim map, frozen manifests, environment
lock, discrepancy log including unfixed defects, and the parity report. A Phase-I value that is cited again is
labelled verified, different but explained, unreproducible, or not enough source data.

\gate{\textbf{Gate P0 (hard stop; part of G0).} Mechanism work on a model or backend proceeds only if the extended
parity check passes there, i.e.\ the intervention hooks demonstrably affect the claimed tensors and token
positions. If it fails, no mechanism work on that model or backend; behaviour may proceed. A Phase-I number may be
cited only after the spot check; if it changes materially, a versioned erratum goes into the internal ledger and
derived findings are invalidated. Never patch new experiments onto a broken denominator.}

\textbf{Status (11 Nov 2026):} HF-versus-TransformerLens parity \emph{passed} for Qwen2.5-7B-Instruct and
Llama-3.1-8B-Instruct on the smoke prompt (job 399217, 17~October~2026: top-1 equal, cosine 0.99998 and 0.99999);
extended parity at the planned hook sites and trace lengths is not yet run (one short GPU job, on acceptance).
The append-only artifact registry, Slurm job register and claim-ledger rows exist since Stage~14.11 and are
extended with every run. P0.2 and the full P0.3--P0.5 are not started (D1 open). Qwen3-32B has no pinned revision
or manifest; its runs stay \texttt{model\_commit NOT\_RECORDED} unless the scratch hashes are matched to a
published revision after the fact.

%----------------------------------------------------------------------
\subsubsection{P1. Measurement validity: challenge the entire sensitivity construct}
\label{sec:blockP1}

\textbf{Question (October).} Does the observed signal reflect improper computation by the model, or mainly a
property of the string and the scoring interface?

\textbf{Revised scope (G0, G2).} The question has been answered for the legacy construct well enough to retire it
(Section~\ref{sec:outcomes}): $S_{lp}$ and any single-string teacher-forced probability are no longer endpoints, and
the anchor-matched presence boost and $R_{alt}$ are diagnostics only. P1's sub-items therefore survive in two
forms: as completed Stage~14 evidence, and as measurement checks that G2 applies to the new contract (noise floor,
format robustness, parse rates).

\begin{description}[leftmargin=1.2em,style=nextline]
\item[P1.1 Anchor-free rescoring and base permutation.] Done for the legacy label (X02). No further work.
\item[P1.2 No-cue and non-numeric controls.] Cue absent, irrelevant numeric distractor, style-matched non-numeric
sentence, numeric metadata and a re-certified relevant edit, compared as matched differences-in-differences. These
are now standing conditions of the design (Section~\ref{sec:conditions}).
\item[P1.3 Output-wrapper multiverse.] Done on the v2 structures under seven wrappers. Carried forward: exact
tokenizer IDs for every scored sequence and prefix, no prompt or padding tokens in a scored span, and the two
declared final-answer formats that G4 runs.
\item[P1.4 Length and tokenization.] Tokenizer audit done. \emph{Open, optional:} first-token versus full-sequence
scores and value-only versus value-plus-unit likelihood (computable from stored per-token log-probabilities).
\item[P1.5 Proper likelihood comparisons.] Correct-versus-decoy log-odds measured as anchor-matched answer
competition; now a diagnostic. Within a mechanism claim the candidate pair is the one fixed by
Section~\ref{sec:causal}, never the foil the model scores highest.
\item[P1.6 Free generation.] Done with reasoning (Stages~14.13, 14.15, 14.19--14.20). The October requirement of
repeated stochastic samples per cell is now part of the contract (Section~\ref{sec:outcomes}); the transition
categories (correct$\to$wrong, wrong$\to$different wrong, invalid) survive inside O1--O3.
\item[P1.7 Metric reliability.] Done for dtype (bfloat16 reproducible; noise SD 0.09--0.19 nats). Carried
forward: G2 estimates each model's own sampling and self-flip floor on fresh samples.
\item[P1.8 Selective prediction.] \emph{Open, optional:} confidence bins for final-answer correctness, ECE/Brier
with cluster intervals, stable versus unstable structures at matched difficulty.
\end{description}

\gate{\textbf{Gate P1 (October), resolved.} The October branches have been taken: the legacy $S_{lp}$ is not called
a shortcut construct; Cue~C is separated as the convertible-unit (invariance/equivariance) family; and the thesis
moved from likelihood sensitivity to the role-by-value response in free generation (O4). Any new likelihood finding
is reported as a diagnostic and carries no safety or correctness claim on its own.}

\textbf{Status (11 Nov 2026):} largely done in Stage~14, all exploratory. X02: 40--49\,\% (Qwen) and 41--56\,\%
(Llama) of base-0 positives flip with the choice of base. Stage~14.10: the sign of the cue effect persists in all
seven wrappers, its size depends on format for Qwen, and parse rates depend heavily on format. Stages~14.14 and
14.17: every number in a prompt receives an answer-presence boost, a cue more than a relevant input; plausibility
matching removes about half of the same-unit cue effect, and the remainder depends on unit compatibility. Open
and optional: P1.4 value-only and first-token scores, P1.8 selective prediction. The cross-instruction stability
check of any label on the pilot families is a G0 item (Section~\ref{sec:gates}).

%----------------------------------------------------------------------
\subsubsection{P2. Scale the approved DECOY-v2 design (not the old benchmark)}
\label{sec:blockP2}

\textbf{Question (October).} Can the proposed mechanism's causal ingredients be manipulated independently while
preserving physics, with genuinely new structures held out?

\textbf{Revised scope (G1 for M250; after G3 for M500).} P2 builds DECOY-v2 dataset~v2 to 250 independent
structures in G1 and to 500 before G4, under the admission standard and data-role allocation of
Section~\ref{sec:datasetplan} (user decision, 11 November 2026); the October range of 300--600 is replaced. The
generator requirements stay, re-expressed in the role-by-value design.

\begin{description}[leftmargin=1.2em,style=nextline]
\item[P2.1 Explicit generative grammar.] Typed quantities, dimensions, roles and assumptions; the equation as a
directed computation graph from inputs through intermediates to the target. Each quantity carries the certificate
of its stratum (R, D, C or K; Section~\ref{sec:strata}), and an independent semantic checklist verifies that the
rendered text matches the graph. Negative-control structures with a planted coupling must be rejected by the
validators. Free-symbol absence alone is not a certificate when the omitted quantity changes assumptions, frames or
boundary conditions.
\item[P2.2 Factorial design (recast).] The October six axes map onto the design of
Sections~\ref{sec:rolevalue} and~\ref{sec:conditions}: physical status becomes the role strata; numerical
closeness becomes the three proximity levels kept as distinct conditions, plus values $z_1,\dots,z_m$ for the
dependency curve; units become same-unit, incompatible-unit-same-digits and convertible-unit cells; semantic
relationship becomes at least two validated phrasings per role relation, one without marker words; placement
becomes the cue-position condition; task structure becomes the difficulty variants. Fractional or blocked designs
are prespecified, never chosen after outcomes.
\item[P2.3 Crucial deconfounding cells (kept).] Equal cue digits with different physical role; equal role with
different digit overlap; matched answer value under a different law; same law at a different instantiation; same
numeral and unit with a different role phrasing. The role-by-value cells provide most of these by construction.
\item[P2.4 Quantity and independence (kept, target withdrawn).] Count distinct computation graphs, not
paraphrases. Complete template clusters, law combinations and source collections are reserved for confirmation. If
there are too few genuinely different graphs, prefer fewer valid structures and a narrower stated scope.
\item[P2.5 Domain extension.] Each new domain is a separate cohort, not pooled into a favourable aggregate; domain
difficulty is a measured covariate.
\item[P2.6 Out-of-distribution sources.] Moved to G6 (external sets, below under P11). Adapted external problems are
new items and need their own certificates (Section~\ref{sec:rolevalue}).
\item[P2.7 Adversarially plausible wording.] Independent LLMs may flag ambiguity and suggest tricky distractors but
never assign physics ground truth; admission requires rewriting and symbolic validation; rejection reasons and
excluded structures are published.
\item[P2.8 Matched negative and positive sets.] Now the supporting conditions of Section~\ref{sec:conditions},
including role reversal (use on request), which separates a selective relevance computation from indiscriminate
digit suppression.
\item[P2.9 Domain-specific robustness.] Dimension conversions, significant-figure conventions, sign and frame
conventions, rounding, near-zero denominators and degenerate limits are checked; any tacit real-world assumption
gets a stated-assumption audit and may exclude the structure from the core.
\end{description}

\gate{\textbf{Gate P2.} Admission requires a verifier pass, an independent human semantic pass, no open leakage
collision and a declared structural parent or cluster; a structure contributes to the primary contrast only if it
is complete: R and D cells at three values in two phrasings (Section~\ref{sec:rolevalue}); infeasible cells are
recorded, never imputed, and structures are not edited after G1 to force feasibility. If factorial creation yields
artificial text, run a blind naturalness comparison and improve the grammar; do not keep artificial items because
their effects are large. If law breadth cannot be achieved, narrow the generalisation claim.}

\textbf{Status (11 Nov 2026):} typed generator with certificates and a mutation harness (Stage~14.6); 80
certified structures exist (S01--S30 and the Stage~14.19 gold dataset of 260 samples), all EXPLORATION;
S01--S20 have the user's open walkthrough approval, the rest agent verification only, certificates UNREVIEWED.
The role-by-value cells for strata R, D, C and K, the phrasings, the misuse functions and the regenerated
proximity conditions (exact match, near within tolerance and tolerance-separated kept as distinct conditions) are
G1 work, not yet done. No structure exists yet for INTERNAL\_CONFIRMATION.

%----------------------------------------------------------------------
\subsubsection{P3. Independent expert validation, parser reliability and leakage firewall}
\label{sec:blockP3}

\textbf{Question.} Is every confirmatory item's natural language correct, is every graded answer read correctly,
and is no confirmation structure related to anything used in discovery? A positive model result never waives these
requirements. Kept in full; mapped to G1 (rules and mutation tests frozen), G2 (parser audit on the new contract)
and G4 (validation of confirmation structures, cohort freeze).

\begin{description}[leftmargin=1.2em,style=nextline]
\item[P3.1 Personnel and blinding.] Two independent physics-qualified validators per confirmatory structure and a
third adjudicator (interim: the qualified LLM panel of Section~\ref{sec:llmpanel}, under the same blinding), calibrated on easy, adversarial and ambiguous cases, blind to model outputs, scores and the
hypothesis. Each property asked separately: equation correctness, sufficient assumptions, the role certificate of
each quantity (Section~\ref{sec:strata}), answer invariance or certified response, naturalness, missing
interactions, units, rounding, and whether paired renderings differ only as intended. Confidence and reasons are
recorded.
\item[P3.2 Disagreement.] Predefined exclusion and adjudication rules; raw agreement, Cohen's $\kappa$ and a
prevalence-aware alternative under imbalance, structure-level bootstrap intervals, domain failure rates and
adjudication edits and changed-label counts. The October aspiration of $\kappa=0.80$ stays an aspiration, never a
licence to drop hard cases.
\item[P3.3 Mutational verification.] Deliberate failure mutations (quantity leaks into the equation, changed
assumption, removed isolation, unit error, sign flip, changed relevant variable, renamed role, swapped target, hidden dependence, missing
boundary condition) and algebraically equivalent positive mutations; sensitivity and specificity by mutator; invalid
items must fail for the intended reason. For strata C and K the mutations include breaking the cancellation and
leaving the regime.
\item[P3.4 Parser study.] A blinded, stratified sample across models, conditions and formats, including refusals,
multiple numbers, intermediate answers, hallucinated units, LaTeX-wrapped final lines, scientific notation,
significant figures, symbolic answers and truncations; two human adjudicators; coverage, error rate, coverage
conditional on outcome, answer-equivalence false positives, and conservative bounds on behavioural rates. The O3
status categories (\texttt{COMPLETE}, \texttt{BUDGET\_EXHAUSTED}, \texttt{FORMAT\_INVALID}, \texttt{UNPARSEABLE},
\texttt{MULTIPLE\_CANDIDATES}, \texttt{EMPTY}) are audited as categories in their own right; ambiguous responses
stay explicit. The audit reports the parser's false-negative rate; no LLM judge enters a primary analysis.
\item[P3.5 Multi-level overlap detection.] Union-find clusters over prompt hashes, number-masked text, symbolic
AST, equation-graph isomorphism, governing law, frame, renamed graph, cue grammar, near-duplicate embeddings,
reviewed semantic relations and identical solution programs; a frozen clustering policy with a manually
adjudicated sample; transitive siblings never straddle splits. Same law is held out in a stricter transfer test.
\item[P3.6 Audit of the original 34 clusters.] Review of the 2,352 unresolved rendering pairs and 147 family pairs,
and re-evaluation of the original LOOCV by cluster-held-out and by distinct symbolic program, is needed only if v1
families are reused. Always required: a report of overlap between v1 and v2 structures and every new locked
confirmation set (G4, G6); newly generated paraphrases of old computation graphs stay in discovery.
\item[P3.7 Cohort freeze.] Discovery, internal confirmation, EPC confirmation, adversarial challenge, the
human-written source set and the shadow set are manifested, and every structure and variant gets an immutable split before any outcome is seen; a removed
structure keeps a tombstone with its reason; intent-to-evaluate is reported beside eligible results. A structure
never moves from EXPLORATION into a confirmation role.
\end{description}

\gate{\textbf{Gate P3.} No confirmatory behavioural, hidden-state or causal generalisation claim until each
confirmatory structure's natural language has passed independent validation and its leakage risk is audited. If
the certificate proves the equations but validators reject a sizeable share of the prompts' semantics, revise the
benchmark and restart the affected measurements; this is not a statistical nuisance covariate. Until human
review exists, independent validation for G4 is provided by the qualified LLM panel of
Section~\ref{sec:llmpanel}, disclosed as such; without a qualified panel, G4 runs as EXPLORATORY-REPLICATION and no
text may call it confirmatory.}

\textbf{Status (11 Nov 2026):} validator-qualification materials and a 33-item planted-defect bank exist
(Stage~14.5), reviewed by the user as one non-blinded reviewer; no validator is qualified, and agent verification
replaced the physicist review by the user's decision of 26~October~2026 (the reviewers' authorship and validation
objection applies to these items). Structure-overlap clusters rebuilt (Stage~14.4: 34 clusters from 147 family
edges; human review of the 2,352 rendering pairs not attempted). Parser evidence: LaTeX-wrapped final lines made 5
of 78 Stage~A records unparseable (J14.19-009); the LaTeX final-line parse enters the next contract version
prospectively, never applied to old records. Validation for G4: the qualified LLM panel (Section~\ref{sec:llmpanel}); the panel models are open decision D2.

%----------------------------------------------------------------------
\subsubsection{P4. Definitive behavioural discrimination and power planning}
\label{sec:blockP4}

\textbf{Question.} Which of H0--H6 (Section~\ref{sec:hypotheses}) describes the model's behaviour on certified
role-by-value structures, and with what precision? P4 is now the scientific centre of the plan: its exploratory
half is G3 and its preregistered half is G4.

\begin{description}[leftmargin=1.2em,style=nextline]
\item[P4.1 A small preregistered core experiment (recast).] The October primary contrast (proximity $\times$ unit
compatibility $\times$ relevance on answer contamination, correct-versus-decoy log-odds and likelihood dispersion)
is replaced by the role-by-value interaction on O4 (Section~\ref{sec:rolevalue}), with O1, O2, O3 and O5 reported
beside it; proximity, unit and marking are moderators. G4 preregisters the role-by-value family fixed in advance
(Section~\ref{sec:gates}), with its estimands, decision rules and sample size, on held-out structures in untouched
clusters, and opens it once; G3 may add at most one secondary family, labelled as chosen on EXPLORATION data. Broad multiverse tables are secondary.
\item[P4.2 Honest power.] Simulation from discovery-only data with cluster dependence and within-structure
correlation, including the variance of a cue \emph{difference}, not only the no-cue variance; sample size for the
smallest scientifically meaningful effect; attrition from parser and validator failures recorded; a
precision-limited conclusion when the interval cannot be narrowed.
\item[P4.3 Multi-level models.] Structure random intercepts and, where estimable, random slopes, or predeclared
cluster-robust models; only predeclared interactions (role $\times$ value, role $\times$ marking, proximity $\times$
digit overlap, unit compatibility, position, model); family-wise or FDR control over any wider scan, with
confirmation.
\item[P4.4 Near-match profile (kept).] The Phase-I Pearson 0.074 \tg{SUBMITTED} is a warning. Monotone distance,
threshold or tail, and token-overlap models are compared by out-of-cluster predictive density; counterexamples to
each (close but digit-disjoint, digit-similar but far, exact value with incompatible units). Misuse hypotheses are
tested as response curves over several values of $z$, never from a single matching answer.
\item[P4.5 Same panel across models.] Identical certified items on every compared model, chosen inside every
configuration's calibration band (Section~\ref{sec:populations}); thinking on versus off within the same weights
answers the secondary question; exact model revisions and chat templates are recorded.
\item[P4.6 Severity and validity strata.] Positive-structure rate, median and tail effects, answer changes,
correctness transitions, high-confidence errors, O3 statuses and token cost reported separately; effects must
survive dropping the one to three most extreme structures; means and robust quantiles both reported.
\item[P4.7 Core negative tests (kept, restated).] If the departure is reproduced by an irrelevant bibliographic
numeral, is unchanged by unit incompatibility, does not differ between role twins, vanishes under equivalent output
formats, or tracks only token overlap, the evidence counts against H3 and for H1 or H2 in the scorecard (hypotheses are
not exclusive, Section~\ref{sec:gates}); H6-RD is supported only when F1--F3 (D and R) are each shown equivalent to zero;
H6-extended adds C and K. These are results, not
failures to hide; a certified null is reported with its power.
\end{description}

\gate{\textbf{Gate P4 (now G3 then G4).} Mechanism work proceeds only after the role-by-value family has
passed G4 (Section~\ref{sec:gates}); the contrast fixes the claim type and metric of Section~\ref{sec:causal}, and
the legacy normalised recovery is never reused. If the only effect is format-sensitive confidence drift or cost
(H5), report it as such. If strata or conditions behave differently, choose one primary contrast for mechanism and
treat the others as explicit scope tests. A cell with zero evaluable records is \NE{}, never a zero effect.}

\textbf{Status (11 Nov 2026):} G3 and G4 not started; G1 (construct freeze) comes first. The Stage~14 pilots
predate the role-by-value design and contain no matched role twins. The Stage~14.20 cue cells on 32 Qwen3-32B
candidate anchors (job 435605, one sample per cell) gave 2 wrong and 17 not evaluated of 128 cells; the
sampling-noise baseline (job 435606) was cancelled on the user's instruction before it started, so the cue-effect
contrast is \NE{}; the preliminary analysis is selection-biased by construction and read descriptively only. Its
most informative signal is length (role-reversal cells about 45\,\% more tokens than the no-cue anchors).

%----------------------------------------------------------------------
\subsubsection{P5. A definitive monitoring study: does access to hidden states help?}
\label{sec:blockP5}

\textbf{Thesis (kept).} A monitor is an operational advance only if it is deployable under its claimed input
budget, calibrated for its target, statistically and practically better than cheaper alternatives on unseen
structures, and robust to shift. A predictor that reads four renderings is a counterfactual-family diagnostic, not a
one-input detector.

\textbf{Revised scope (G7).} Demoted to an optional side track: the primary question no longer needs a monitor.
The monitored target must be a frozen endpoint of Section~\ref{sec:outcomes}, not the retired $S_{lp}$. Every
comparator carries an \emph{oracle} or \emph{deployable} tag stating what it needs at inference (in particular,
whether it uses the gold answer); no comparator built from the label's own inputs is called ``non-activation'' or
``one forward pass''.

\begin{description}[leftmargin=1.2em,style=nextline]
\item[P5.1 Three separate tasks.] S (single prompt, no other family member at inference), F (full family, compared
with a multi-query behavioural measurement), V (variant-local, per-variant labels); no pooled AUROC across them.
\item[P5.2 Three frozen feature sets.] One predeclared residual layer from discovery, the embedding layer as
lexical control, next-token logits/entropy as cost-matched comparator; other sites are discovery scans.
\item[P5.3 Comparator ladder.] Constant and base-rate predictors; prompt-only text features; frozen-encoder
embeddings; one-pass output uncertainty; generic correctness and semantic-entropy predictors; black-box multi-query
consistency with real pass counts; post-generation critics (not pre-generation competitors); a fully tuned
non-activation stack. Each with its oracle/deployable tag.
\item[P5.4--P5.6 Probe capacity, generalisation level, leakage controls.] Capacity-matched probe families with
random-label controls; prediction across instantiations, clusters, laws, conditions and checkpoints (a signal that
needs the same scaffold in training is template recognition); permuted-label and irrelevant-metadata controls.
\item[P5.7--P5.9 Deconfounding, selection accounting, calibration.] Incremental value over correctness and entropy
by nested stacking and paired loss on the external set; layer chosen on discovery and frozen; calibrators fitted on
development groups only; FNR at FPR 0.05 and 0.10 \tg{PROPOSAL} with full curves and intervals; prevalence shift
evaluated separately.
\item[P5.10--P5.11 Information bounds and cost.] Oracle features are non-deployable upper bounds; wall-clock,
memory and throughput per decision on the same hardware.
\end{description}

\gate{\textbf{Gate P5 (kept).} A monitor is promoted only if a frozen single-input specification shows a
practically meaningful, statistically supported incremental advantage on a held-out structural cohort under a
preregistered endpoint. Otherwise monitorability is a descriptive diagnostic and the null incremental value is
reported. An AUROC above 0.5 is not incremental value.}

\textbf{Status (11 Nov 2026):} side track, not scheduled (open decision D6). The Phase-I monitor under
cluster-held-out folds is \NE{}: its activation tensors were lost and must be regenerated. Known from committed
files: against the oracle-free entropy baseline (0.520) the paper's comparator (0.727, fitted on the label's own
inputs) \tg{SUBMITTED} was oracle-like; the same applies to the Wave~2 comparator. The oracle/deployable tag and its
propagation to the Wave~2 record and the claim ledger are G0 items.

%----------------------------------------------------------------------
\subsubsection{P6. Re-localise the phenomenon under valid causal metrics and selection rules}
\label{sec:blockP6}

\textbf{Question (revised).} Which components, if any, have an effect on the role-by-value contrast that passed G4
--- an interaction $I_h\neq0$ under the declared claim type (T1, T2 or T3; Section~\ref{sec:causal}) --- rather than
a general effect on the answer?

\textbf{Retired as stated: the H11 lesion (card X07).} The October question asked whether Qwen layer-16 head~11 is
a causal component. The Phase-I evidence does not support starting there: the four-head L16 knockout inverted the
signed drop (normalised recovery above one) in 17 of 20 families, mean recovery 1.87, and removing head~11 alone
raised the certified answer's first-token logit by 3.76 on base and 3.77 on sensitive renderings, in 20 of 20
families each \tg{SUBMITTED}, verified from the released files. An effect equal on both prompts marks a general
answer component. Components are found anew under the interaction estimand; the legacy L16 sites, and the Llama
L21~H2 candidate, have no privileged status.

\begin{description}[leftmargin=1.2em,style=nextline]
\item[P6.1 Pairs chosen prospectively.] The counterfactuals are the role-by-value minimal pairs (R and D twins of one
structure at the same $z$; the same role at two values), token-aligned, several per item, labelled by design and
not by observed sensitivity; stratified by correctness and by O3 status. Where alignment is impossible, only
semantically matched positions are patched, with position controls. Selection on a discovery split; tests on
disjoint held-out structures.
\item[P6.2 Frozen response variable.] The metric of the declared claim type (Section~\ref{sec:causal}): T1
$\log p(y^*(z))-\log p(f_h(z))$ on the declared prefix, between prompts that share a gold, each candidate's logit
shown; T2 the certified intermediate against the prespecified
misuse value; T3 O1, O2 and O4 over on-policy resampled traces. Original score, intervened score and absolute
difference are always reported, for $D$, $N$ and (under T2/T3) $R$, with the clean effects beside
$I_{\text{presence}}$ and $I_{\text{role}}$; never a normalised
recovery and never a ratio over near-zero denominators.
\item[P6.3 A new experiment, not a reproduction.] Candidates come from a discovery screen on EXPLORATION
structures, with flanking layers, norm- and entropy-matched heads and random sites as controls; per-structure
before/after values reported.
\item[P6.4 Several ablation semantics.] Resample or interchange ablation from matched counterfactuals as primary,
mean ablation as sensitivity, the conclusion stated under both; projection removal and rescaling as further
sensitivities; zero ablation never as evidence. A method-by-site stability table, not one intervention flavour.
\item[P6.5 Specific necessity.] Noising a component must change the role-specific effect while preserving, within
a preregistered tolerance, the response to relevant inputs (the R twin's certified response), unrelated problems
and role reversal. No global claim of logical necessity: redundant pathways are expected.
\item[P6.6 Sufficiency with the right variables.] Denoising and transfer begin within structure, value-matched,
between the R and D twins; cross-structure donors come later. Dose sweeps in both directions; effects on generated
answers, not only on logits; content-preserving swaps from same-answer, same-law and same-unit irrelevant donors.
An injection that changes a likelihood score but not answers or the certified-versus-misuse competition is not
called sufficiency.
\item[P6.7 Off-target damage and self-repair.] An independent control suite (no-cue physics, relevant edits,
unrelated numeric and lexical tasks, common natural-language QA and syntactic completion); loss and KL shifts on
neutral prompts, next-token entropy and output length; compensatory changes in nearby heads and MLPs; local
replacement against global ablation; divergence of patched activations from the natural
distribution; direct and total effects with a frozen-final-LayerNorm control.
\item[P6.8 Second model.] Llama-3.1-8B-Instruct is localised afresh on its own discovery split; a site that misses
triggers a discovery sweep on exploratory data, never repeated scans of confirmation structures. Reasoning models
follow the T1/T3 split of Section~\ref{sec:causal}. October's DeepSeek requirement (confirm the prompt boundary,
reasoning trace and final-answer scoring, and match denominators, before concluding that a reasoning model lacks a
pathway) applies to any reasoning model studied; DeepSeek-R1-Distill itself is not planned (Appendix~\ref{app:disposition}).
\end{description}

\gate{\textbf{Gate P6 (part of G5).} A component advances to P7 as a role component only if $I_{\text{role}}\neq0$ under T2 or
T3 with the controls of Section~\ref{sec:causal}, robust across ablation semantics, on held-out structures; with
$I_{\text{presence}}\neq0$ alone it is a numeral-cue component. An effect equal on $D$, $N$ and $R$ marks a general
answer component and is reported as such. An effect shown only under T1 supports a readout claim,
not a role-computation claim. An effect on log-odds but not on O1, O2 or O4 is a cause of answer-competition change,
not of success or failure in physics.}

\textbf{Status (11 Nov 2026):} not started; gated on G4. No activations exist for the v2 items. Parity for the two
7--8B models is in hand (P0); extended parity is not.

%----------------------------------------------------------------------
\subsubsection{P7. Identify what the causal components compute}
\label{sec:blockP7}

\textbf{Question.} What function connects the quantity's tokens to the model's dependence, or independence, of
its answer on that quantity? Localisation without function is not a mechanism claim.

\begin{description}[leftmargin=1.2em,style=nextline]
\item[P7.1 Anatomy of each surviving component.] (Was: head~11 at L16.) For every relevant position (inputs, the
quantity's value and unit, the role phrasing, the question, the last prompt position, answer positions): attention
logits and patterns, output norm, OV contribution and direct attribution across the role-by-value pairs. Attention
maps are not causal evidence; attention probabilities and value vectors are patched separately.
\item[P7.2 Candidate functions tested head-to-head (kept), mapped to the revised hypotheses.]
\begin{itemize}[leftmargin=1.2em]
\item \emph{F-copy} (retrieves or copies a distractor numeral) and \emph{F-answer prior} (shifts one output format
or length without representing the quantity) $\to$ H1;
\item \emph{F-match} (equality or proximity of numeral and target) and \emph{F-unit} (dimension or unit matching)
$\to$ H2;
\item \emph{F-role} (semantic disconnection or role relation suppresses an attractive candidate) $\to$ H3;
\item \emph{F-position} (near-query tokens or sentence order) $\to$ the cue-position control;
\item added: \emph{readout} (correct value computed in the trace, final answer moved toward the quantity) $\to$ H4,
tested under T1 at answer positions on a fixed trace;
\item added: \emph{correct selection} (D, C, K invariant and R correct) $\to$ H6, whose mechanism is studied when
errors are rare.
\end{itemize}
Each function needs at least two interventions whose predictions differ from the alternatives: same digits with
an incompatible unit (F-copy yes, F-unit no); same magnitude, digit-disjoint form (F-match yes, F-copy no); same
tokens moved (F-position changes, F-match preserved); the R twin or role reversal (F-role reverses, F-copy need
not). These are cells of Sections~\ref{sec:rolevalue} and~\ref{sec:conditions}, evaluated on a held-out grid, not
on anecdotes.
\item[P7.3 Activation decoding with interventions.] Nested linear probes for the numeral, target magnitude, digit
overlap, unit compatibility, role and relevant inputs are diagnostic only; subspace interventions that remove one
decoded variable while preserving others test behaviour; causal directions are classified as independent, oblique
or redundant. A decodable subspace is diagnostic; only a subspace whose interchange changes behaviour as predicted
is causally sufficient.
\item[P7.4 Representation geometry.] Matched differences $\Delta(z_1,z_2)$ and $\Delta(\mathrm{R},\mathrm{D})$ at key
sites; cross-structure alignment (CKA, CCA, principal angles, projection energy) against equal-dimension random and
nearest-neighbour nulls; subspaces fitted on discovery only and used to predict held-out causal effects.
\item[P7.5 Direct logit attribution.] On the prespecified candidate pair only, followed by exact counterfactual
patching; decomposed into own-role candidate, other-role candidate, both, or generic length.
\item[P7.6 Temporal localisation.] Now expressed through the claim types: T2 on a teacher-forced prefix up to the
step where the quantity enters (for C and K intermediate use is expected, not penalised); T1 at answer positions on
a fixed trace; free explanations aligned only by semantic events and descriptive unless resampling makes positions
comparable.
\item[P7.7 Invariance and role reversal.] The proposed representation is patched between the R and D twins at the
same $z$; both the normal and the role-reversed response are predicted. A transfer must change the computation as
predicted while preserving value, unit and unrelated answer content.
\item[P7.8 Hypothesis scorecard.] For each function: predeclared sign predictions over the cells, held-out
prediction score (sign and size of $I_h$ committed before the test), selective patch effect, falsifying control,
the alternative it beats, and failure cases. The simplest function that explains held-out intervention effects is
preferred.
\end{description}

\gate{\textbf{Gate P7 (part of G5).} Advance only if at least one explicit functional hypothesis predicts effect
direction across independent manipulations and survives direct causal tests under T2 or T3. Otherwise report a
localised but uninterpreted causal handle, try larger multi-component representations (P8), or pivot to Track~E
(the behavioural and measurement result). No activation feature is called ``physics-relevance'' because its top examples sound
physics-related.}

\textbf{Status (11 Nov 2026):} not started; gated on P6. Under H1 or H2 after G3, only the optional
unit-matching study is in scope, framed against the closest prior accounts (Section~\ref{sec:question}).

%----------------------------------------------------------------------
\subsubsection{P8. Recover and validate upstream and downstream computation, not one head}
\label{sec:blockP8}

\textbf{Question.} Which source-to-mediator-to-sink path carries the role information, and does an explicit
high-level algorithm agree with the network under interchange? Kept; gated on G4 and on P6/P7.

\begin{description}[leftmargin=1.2em,style=nextline]
\item[P8.1 Sources and sinks declared first.] Sources: the quantity's value and unit tokens, the role phrasing,
relevant inputs, law words, the format instruction. Sinks: mediators that passed P6/P7 and the metric of the
declared claim type. Contrasts: the R/D twins at the same $z$ (primary) and two values within a role; proximity
contrasts as moderators.
\item[P8.2 Two discovery procedures.] Attribution patching (with integrated-gradient refinement where it saturates)
as a ranking screen, plus targeted activation patching; optional mask optimisation. Every reported component is
confirmed by exact patching.
\item[P8.3 Path patching.] Receiver-conditioned path patching for source to mediator, mediator to downstream and
downstream to answer, with checks of hook order, graph direction, tolerance and reverse interventions; the path must
matter only when the mediator receives the source information.
\item[P8.4 Knockout--rescue ordering.] Knock out the source and measure the loss; rescue with the mediator's clean
activation; knock out the mediator and predict that upstream rescue fails; control with unrelated and same-answer
donors; rescue a role-reversed task to show the path is not a destructive filter; contrast against rescuing a
generic next-token loss change.
\item[P8.5 Completeness and minimality.] Model run with only the hypothesised edges against matched replacement of
the complement; sufficiency, comprehensiveness, selectivity, stability over corruptions; effect-versus-size and
specificity-versus-size curves in absolute metric units with each candidate shown; random matched graphs and a
whole-residual patch as baselines; residual unexplained effect reported.
\item[P8.6 Causal abstraction.] An explicit high-level algorithm (e.g.\ parse inputs, parse candidate values,
evaluate role and unit, choose the answer) with a map from neural subspaces to its variables; a role variable
localised by DAS with interchange-intervention accuracy on held-out pairs against neuron and random-subspace
baselines, plus the subspace-illusion check (Section~\ref{sec:causal}). One high-level variable is varied at a
time while correlated variables are held fixed. Poor interchange agreement falsifies the
account whatever the knockout sizes.
\item[P8.7 Bottleneck hypotheses.] Does the mediator store the raw value, normalised magnitude, a role flag, a
source identifier or a merged candidate; does it act as a routing gate or a logit amplifier; confirmed on held-out
laws, magnitudes, prompts and positions. Several binding mechanisms may combine; one isolated role feature is not
assumed.
\item[P8.8 Sparse features (optional).] SAEs, transcoders, DAS and neurons compared at equal compute and feature
budget; reconstruction error, seed stability and causal interchange evaluated; SAEs are a method under test.
\item[P8.9 Synthetic ground truth.] A small model trained with a planted numeric rule, and a matched model without
it, tests whether the pipeline recovers a known route; recovering it does not prove the same route operates naturally
in the studied models, but failure to recover it questions the method.
\end{description}

\gate{\textbf{Gate P8 (part of G5).} A mechanism claim needs a compact algorithmic account with validated
source-to-mediator-to-sink effects on untouched structures and held-out predictions that a simpler lexical or generic
alternative cannot explain (the mechanism rung of Section~\ref{sec:claims}). Redundant routes are reported as a
distributed or partially redundant mechanism, not forced into a single-head story. If only correlational graphs
emerge, conclusions stop at localisation.}

\textbf{Status (11 Nov 2026):} not started; gated on G4.

%----------------------------------------------------------------------
\subsubsection{P9. Cross-model, cross-domain and post-training generalisation}
\label{sec:blockP9}

\textbf{Question.} Does the identified function recur beyond the discovery model, domain and format? Kept; gated
on G4 and on a frozen function from P7/P8. The within-model replication is part of G5; external and cross-domain
confirmation is G6.

\begin{description}[leftmargin=1.2em,style=nextline]
\item[P9.1 Cross-architecture confirmation.] Freeze the function, not the head index. Fresh discovery samples
localise functional analogues in the second architecture; the hypothesis is evaluated on untouched clusters with the
same claim type and metric. A functional analogue in a different head is replication; a failed matched-layer
knockout is not a complete null.
\item[P9.2 Scaling ladder.] Conditional on compute, small to large within one family; behaviour across scale before
expensive causal cases; only models with parity-tested hooks and answer scoring. In this plan Qwen3-32B enters only
for a frozen-function transfer test.
\item[P9.3 Same-backbone post-training.] Base, instruct and reasoning variants with prompting controlled; thinking
on versus off within the same weights. Differences are not attributed to RL unless comparable training runs isolate
it.
\item[P9.4 Familiar-to-novel-law transfer.] Valid symbolic worlds or dimensionally coherent toy laws with explicit
rules and difficulty matched by solver-assisted pilots; together with the non-physics dependency control
(Section~\ref{sec:conditions}) this separates a general numeric or dependency phenomenon from one that needs
physical semantics.
\item[P9.5 Source and language transfer.] Human-written problems, new formats, terse equations, narratives and
validated bilingual renderings; English-only results do not imply cross-lingual generality.
\item[P9.6 Composition and adversarial boundary.] Two or three simultaneous quantities of different roles;
additive, super-additive or suppressive interactions in behaviour and in the causal effect; adversarial search under
a fixed budget, retested on freshly authored confirmation items.
\item[P9.7 Persistence across tools and precision.] Two hook backends where feasible; bf16 against fp16 and modest
quantisation behind a numerical agreement gate; an effect that depends on an implementation bug is retired.
\end{description}

\gate{\textbf{Gate P9 (G5 for the second model, G6 beyond).} Separate conclusions for behaviour, monitor, function
and anatomy. If the function recurs but component identities differ, the heterogeneity is part of the result. If
only one model shows it, a deep single-model case study needs stronger within-model causal proof and a stated reason
why it generalises conceptually; no universal claim about language models is made.}

\textbf{Status (11 Nov 2026):} not started; gated on G4. Qwen3-8B (thinking on/off) is not downloaded (open
decision D3).

%----------------------------------------------------------------------
\subsubsection{P10. Mechanism-guided repair and comparative efficacy}
\label{sec:blockP10}

\textbf{Question.} Does knowledge of the mechanism yield a repair that beats simple non-mechanistic methods?
Kept as optional practical use (G7); it starts only after P7 yields a specific intervention. A cheap repair that
works is a result in its own right.

\begin{description}[leftmargin=1.2em,style=nextline]
\item[P10.1 Operational loss.] Weighted O1/O2 error, role-by-value departures (O4), calibration and OOD robustness,
subject to a maximum degradation on no-cue and relevant-input tasks; never a teacher-forced gold-likelihood
statistic alone.
\item[P10.2 Non-mechanistic ladder (first).] The explicit relevance instruction (``identify and ignore irrelevant
information'', already a G3 condition), self-consistency or majority vote, a longer reasoning budget, then a
deterministic semantic filter, rephrasing into governing variables, a small symbolic solver, two-query consistency
and targeted counterfactual augmentation; latency, compute and the damage from wrongly ignoring relevant inputs
recorded. This ladder sets the bar.
\item[P10.3 Training-free circuit interventions.] Conditional replacement, projection of a causal direction, gating
of a proven edge, small direction edits, patching from a matched no-cue rendering; intensity tuned on development
only, with a strength-versus-damage curve and checks on prompts where the same quantity is relevant (the R twin).
\item[P10.4 Monitor-triggered repair.] Only if P5 passes; triggered against universal application under realistic
base rates; an oracle trigger is a non-deployable upper bound.
\item[P10.5 Training-based mitigation.] If justified, certificate-backed paired consistency training with a
correctness loss and relevance-preservation controls, against random distractor augmentation, strong-distractor
curricula and generic instruction tuning; whole law clusters held out; follow whether the mechanism weakens,
relocates or is bypassed.
\item[P10.6 Attack and defence.] Certified adversarial edits under declared attacker access and budget; adaptation
curves and fresh-attack transfer.
\item[P10.7 Causal linkage.] Targeted edit against equal-norm random, nearby-layer, alternative-component and
prompt-only baselines and no edit; a repair that works without the mechanism explaining why is an empirical
intervention, not a causal account.
\end{description}

\gate{\textbf{Gate P10 (G7).} A positive result needs held-out paired improvement in a practically meaningful
outcome (preferably O1, O2 or O4), statistically resolved harm bounds, and an advantage or distinctive trade-off
against the non-mechanistic ladder. If cue removal or the relevance instruction works better with less damage, the
neural edit is not called operationally superior; it may still be a mechanistic test.}

\textbf{Status (11 Nov 2026):} not started. A Phase-I reviewer asked for exactly this downstream use; the first
rung (relevance instruction) is measured in G3 at no extra cost.

%----------------------------------------------------------------------
\subsubsection{P11. Confirmatory closure, adversarial review and claim freeze}
\label{sec:blockP11}

\textbf{Question.} Can an independent reader reproduce the primary result and state the causal account and the best
attempt to falsify it? Kept; it closes every gate and owns the external confirmation of G6.

\begin{description}[leftmargin=1.2em,style=nextline]
\item[P11.1 One causal chain, not the lab log.] Figures follow the claim ladder of Section~\ref{sec:claims}: the
role-by-value design and hypotheses; the measurement audit (generated answers against likelihood diagnostics); the
role-by-value effect on independent structures; interventions with absolute effects; the path with its
counterfactual predictions; replication or repair. Failures appear in compact explicit tables.
\item[P11.2 Mandatory final audit.] Independent recomputation of the main numbers, validation signatures, code and
model versions, contamination audit, preregistration deviations, multiple-testing ledger, negative and zero-row
results, exact intervention sites, intervals, baselines, alternative metrics and the statistical unit; a mechanical
check of every draft's wording against the claim registry's prohibited wording; a check that provenance fields
(authorship, validation) match what the text says. An unaffiliated peer attacks the strongest claim from the raw
evidence and the released protocol.
\item[P11.3 Two-stage confirmation.] The internal hold-out is used for one explicit confirmatory decision (G4). A
fully untouched external cohort, from the EPC external physics sets or other newly written, law-held-out material,
is opened once in G6, after the hypothesis, functional representation, causal site and inference code are frozen.
Any tuning after opening turns it into development; new data are then needed.
\item[P11.4 Reproducible release.] Prompts, symbolic programs, certificates, validator labels, all eligible outputs
and log-probabilities, inclusion metadata, raw before/after intervention values, code, environment, seeds, hardware
and checksums; passing and failing runs with an explicit authority hierarchy.
\item[P11.5 Venue chosen after the result.] A mechanism narrative needs a substantial new insight; a measurement
result takes an evaluation or benchmark framing; a mitigation result a robustness framing. The narratives are not
kept at equal weight unless the evidence forms one non-inflated chain.
\end{description}

\gate{\textbf{Gate P11.} No submission until an independent reader can reproduce the primary figures and state, in
one paragraph, the causal explanation and the best attempt to falsify it. Disconfirming evidence enters the headline
scope, not future work. Each claim states its validation tier (LLM panel or human expert); without either, results stay exploratory;
any manuscript discloses agent authorship and LLM validation.}

\textbf{Status (11 Nov 2026):} the Phase-I paper was withdrawn on 11~November~2026; its reviews and a verification of
the reviewers' claims are recorded, and no Phase-I claim survives as stated. The claim ledger exists. The EPC
external cohort is not built (October B5, demoted to G6). No venue is chosen.

% ===== s06_inventory =====
\section{Complete downstream experimental and ablation inventory}
\label{sec:inventory}

This section keeps the October edition's catalogue of 104 registered contrasts (A01--A18, B01--B14, C01--C18,
D01--D16, E01--E18, F01--F20) with its identifiers, so that the Stage~14 journal and ledger still resolve. The
items are \emph{registered candidates, not commitments}: the critical path is the gate sequence of
Section~\ref{sec:gates}, and the Phase-B and gate contrasts defined there are additional to this list. As in the
October edition, the catalogue is not 104 independent hypothesis tests: contrasts are grouped into
preregistered families, screening is separated from confirmation, a confirmatory subset is selected before the
corresponding test cohort is opened, and low-value branches are stopped early. A negative result removes the
branch's supported claim, not the datum.

\textbf{Columns.} \emph{ID} carries the October tier under the identifier: R = required on the critical path,
C = conditional on earlier evidence, X = exploratory or high-upside. \emph{Contrast} is a condensed but faithful
restatement of the October item (what is manipulated and what is read out). \emph{Gate} is the gate at which the
contrast would run (G0 integrity closure, G1 physics construct, G2 measurement and calibration, G3 behavioural
discovery, G4 independent behavioural confirmation, G5 mechanism, G6 external and functional confirmation, G7
optional practical use). \emph{Status} is one of:
\begin{itemize}
\item \textbf{DONE-EXPLORATORY} --- run in Stage~14 with the cited journal entry; exploratory evidence (data role
EXPLORATION), never confirmation; ``in part'' says which part was not run;
\item \textbf{OPEN} --- not run; kept as specified, with the readout of Section~\ref{sec:outcomes};
\item \textbf{REVISED} --- kept, but its target, outcome or object is changed by the revised core
(Section~\ref{sec:framework}); the replacement is stated;
\item \textbf{RETIRED} --- not run as specified; the reason is stated;
\item \textbf{SUPERSEDED} by a named core element that answers the same question more directly.
\end{itemize}
Every item that rests on the retired legacy label ($S_{lp}$, base~0, 0.5~nats; retired for v2 analyses in
J14.2-011) or on normalised recovery is marked REVISED (with its replacement outcome among O1--O6) or RETIRED.
A cell or contrast with zero evaluable records is \NE{}, never a zero effect.

% ---------------------------------------------------------------- A
\subsection{A. Input and task semantics (A01--A18)}
\label{sec:inv-A}

\begingroup\small
\begin{longtable}{@{}P{1.2cm}P{7.2cm}P{1.3cm}P{5.6cm}@{}}
\toprule
ID & Contrast (manipulated / read out) & Gate & Status \\
\midrule
\endfirsthead
\toprule
ID & Contrast (manipulated / read out) & Gate & Status \\
\midrule
\endhead
\bottomrule
\endfoot
A01 \newline R & No cue vs unrelated neutral sentence vs unrelated numeral; isolate nuisance effects from mere insertion. Paired by law, length and target. & G3 &
\textbf{DONE-EXPLORATORY} (J14.8-002; J14.14-005), likelihood and greedy answers on ten structures. Kept: no cue and the style-matched non-numeric sentence are supporting conditions (Section~\ref{sec:conditions}), read in O1/O4 with repeated samples. \\
A02 \newline R & Numeric near-match vs far-match at fixed units and cue grammar; nonlinear response and tail conditional effects. & G3 &
\textbf{DONE-EXPLORATORY} (J14.8-002; J14.9-004), likelihood level only. For G3 the proximity levels are three distinct conditions (exact, near within tolerance, tolerance-separated); value dependence is read as O4 response curves; cues within a view's tolerance are flagged source-ambiguous. \\
A03 \newline R & Exact printed-number match vs mathematical equality under alternate decimal/scientific notation; separate string identity from magnitude. & G3 &
\textbf{OPEN}. Read as O4 copy tracking under the three grading views. Would test the notation account of the Qwen attenuation under the scientific wrapper, which J14.10-005/-007 leave unseparated. \\
A04 \newline R & Digit overlap matched while numerical proximity differs; lexical-copy mechanism. & G3 &
\textbf{OPEN}. H1 copy-curve prediction (O4). \\
A05 \newline R & Numerical proximity matched while digit/token overlap differs; magnitude-comparison mechanism. & G3 &
\textbf{OPEN}. Stage~14 near-match cells with high and low leading-digit overlap exist (J14.6-006), with descriptive copy counts only (J14.8-002). \\
A06 \newline R & Same-unit vs dimensionally incompatible distractor with identical numeral; dimensional reasoning or gating. & G3 &
\textbf{DONE-EXPLORATORY} (J14.12-002; J14.14-005; J14.17-006): the likelihood boost depends on unit compatibility. Kept as a supporting condition for H2, read in O1/O4. \\
A07 \newline R & Same-unit vs exactly convertible-unit distractor of identical physical magnitude; text-unit vs physical-unit coding. & G3 &
\textbf{OPEN}. Convertible-unit cell exists (J14.6-006), copy counts only (J14.12-002). Kept as the invariance/equivariance family (old Cue~C); never pooled with D cells. \\
A08 \newline R & Explicitly mechanically isolated vs implicitly separate vs now coupled system, the coupled case re-certified for relevance; semantic causal boundary. & G1, G3, G4 &
\textbf{SUPERSEDED} by the role-by-value design (Section~\ref{sec:rolevalue}): R vs D (and C, K) cells at $\geq3$ values in $\geq2$ validated phrasings, one without marker words; O4 (tracking loss, misuse and neglect shares; $\lambda$, $\Lambda$ descriptive). \\
A09 \newline R & Cue irrelevant to the original target but relevant to a counterfactual target question; selective relevance vs global numerical suppression. & G3 &
\textbf{SUPERSEDED} by the R stratum of the role-by-value design (question fixed, role changed). Role reversal (question changed) stays as the use-on-request supporting condition; exploratory cells J14.12-002, J14.19-014. \\
A10 \newline R & Cue before relevant variables, after them, and adjacent to the target question; recency/position effect. & G3 &
\textbf{OPEN}. Kept as cue position (before/after the relevant inputs, Section~\ref{sec:conditions}); position was fixed in all Stage~14 cells (J14.6-006). \\
A11 \newline R & One vs two vs three distractor appearances, total prompt length matched; repetition amplification. & G3 &
\textbf{OPEN}. \\
A12 \newline R & Distractor as factual sentence vs quoted example vs suggested answer; discourse role. & G3 &
\textbf{OPEN}. The suggested-answer arm is the anchoring setting of Owusu, Wiegreffe and Feldman; results are framed against it. \\
A13 \newline C & Cue as table cell, bullet, equation, prose and code-like JSON; representational interface without changing physics. & G3 &
\textbf{OPEN}. \\
A14 \newline C & Correct target zero/near-zero, very large/small, positive/negative and rounded; identify pathological relative-proximity definitions. & G1 &
\textbf{OPEN}. Fixes the proximity definition and the grading views (relative error stored per record) before G2. \\
A15 \newline C & Target precision changed (exact irrational, 3 decimal places, 1 significant figure); answer-rendering ambiguity. & G1 &
\textbf{OPEN}. Bears on the last-printed-place rule and on O2. \\
A16 \newline C & Unrelated intermediate quantity inserted near a non-final intermediate answer while the final answer is far; where numerical competition occurs. & G3, G5 &
\textbf{OPEN}. Behaviour in O4; mechanism only under claim type T2 (certified intermediate values, Section~\ref{sec:causal}). \\
A17 \newline C & One-step vs multi-step problem with the same final answer and comparable length; is internal computation depth causal. & G2, G3 &
\textbf{OPEN}. G2 difficulty variants record step count; value effects are estimated within role, difficulty as covariate. \\
A18 \newline X & Contradictory external sentence or malicious decoy; not an invariant family unless the contradiction's semantics are resolved; out-of-scope robustness quantified separately. & G7 &
\textbf{OPEN}. Never pooled with the role strata. \\
\end{longtable}
\endgroup

% ---------------------------------------------------------------- B
\subsection{B. Dataset integrity, ground truth and representation of physics (B01--B14)}
\label{sec:inv-B}

\begingroup\small
\begin{longtable}{@{}P{1.2cm}P{7.2cm}P{1.3cm}P{5.6cm}@{}}
\toprule
ID & Contrast (manipulated / read out) & Gate & Status \\
\midrule
\endfirsthead
\toprule
ID & Contrast (manipulated / read out) & Gate & Status \\
\midrule
\endhead
\bottomrule
\endfoot
B01 \newline R & Free-symbol certificate vs independently programmed equation-solver certificate on a sample; disagreements to manual adjudication. & G1 &
\textbf{DONE-EXPLORATORY} for the v2 structures (J14.6-007; J14.19-002): each carries an independent solver route; agent-verified only, UNREVIEWED. To be extended to the stratum certificates (graph disconnection for D, symbolic cancellation for C, regime predicate for K). \\
B02 \newline R & Symbolic mutation that introduces cue dependence; verifier catch rate, not only acceptance of valid equations. & G1 &
\textbf{DONE-EXPLORATORY} (J14.6-005; mutation harness rerun in J14.19-002). Same extension to C and K certificates. \\
B03 \newline R & Physical-assumption mutation removing ``massless'', ``frictionless'', ``mechanically isolated'', ideal-source or equilibrium constraints; human and formal checks. & G1 &
\textbf{DONE-EXPLORATORY} in part: formal (dropped-assumption mutator, J14.6-005); human (planted bank reviewed by the user, J14.5-010; physicist confirmation open). Central to K-stratum regime predicates. \\
B04 \newline R & Same target with a dimensionally incorrect unit; validator and parser must flag the error rather than silently normalise it. & G1 &
\textbf{OPEN}. Core rule already adopted: units converted with a dimension check, never stripped (Section~\ref{sec:outcomes}). \\
B05 \newline R & Exact prompt hashes vs AST/graph-normalised structural duplicates; re-estimate the main probe on a cluster-held-out split. & G0, G1 &
\textbf{DONE-EXPLORATORY} in part (J14.4-002; J14.4-009): clusters rebuilt, non-activation predictors refit. The probe is \NE{} (no activations), a G7 item. New structures: program-hash overlap checks (J14.6-002; J14.19-002). \\
B06 \newline R & Dual expert human validation blinded to generated outputs vs single expert or agent audit; disagreement and exclusions. & G1, G4 &
\textbf{OPEN}. Human tier; for G4 the interim tier is the qualified LLM panel (Section~\ref{sec:llmpanel}); Stage~14 items are agent-verified and UNREVIEWED (J14.19-002). \\
B07 \newline R & Native answer parser vs blinded human consensus; failure modes and impact on sample admission. & G2 &
\textbf{OPEN}. Becomes the stratified parser audit (false-negative rate); LaTeX-wrapped final lines already observed (J14.19-009). \\
B08 \newline R & Family selected for scoring before vs after parseability; selective attrition of the five historically removed stable families. & G0 &
\textbf{DONE-EXPLORATORY} for legacy predictors (J14.4-004; J14.4-011). Rule absorbed into O1/O3: no conditioning on evaluability. \\
B09 \newline R & Held-out same-law new numbers vs held-out new laws/graphs; structural generalisation gap. & G4, G6 &
\textbf{OPEN}. G4 uses new template clusters; law-held-out external items in G6. \\
B10 \newline C & Cue-type-balanced vs natural-prevalence panels; base-rate-driven changes of pooled AUROC. & G7 &
\textbf{REVISED}. Strata are never pooled (Section~\ref{sec:strata}, rule~3); the prevalence question survives only for a G7 monitor with an O1/O4 target. \\
B11 \newline C & Blinded human-naturalness preference against curated real physics exercises; reject conspicuously synthetic cues. & G1 &
\textbf{OPEN}. Supports the unmarked-phrasing requirement: most Stage~14 cue cells carry an explicit ``unrelated''-type disclaimer. \\
B12 \newline C & Parallel physics solvers or independent human derivations for subtle thermodynamics, frame and circuit cases. & G1 &
\textbf{OPEN}. Only agent blind solvers so far (J14.19-002). \\
B13 \newline C & Construct-validation error bounds: pessimistic/optimistic imputation of all unresolved semantics; bound the primary effects. & G4 &
\textbf{OPEN}. \\
B14 \newline X & Distinct-model content-contamination probes and source-provenance checks; limits of inferring pretraining membership. & G6 &
\textbf{OPEN}. \\
\end{longtable}
\endgroup

% ---------------------------------------------------------------- C
\subsection{C. Scoring, calibration and behavioural endpoints (C01--C18)}
\label{sec:inv-C}

\begingroup\small
\begin{longtable}{@{}P{1.2cm}P{7.2cm}P{1.3cm}P{5.6cm}@{}}
\toprule
ID & Contrast (manipulated / read out) & Gate & Status \\
\midrule
\endfirsthead
\toprule
ID & Contrast (manipulated / read out) & Gate & Status \\
\midrule
\endhead
\bottomrule
\endfoot
C01 \newline R & Arbitrary base 0 vs all four rotated bases; count changed sensitivity labels without new model calls. & G0 &
\textbf{DONE-EXPLORATORY} (X02, J14.2-006): 40--56\% of legacy positives flip with the base; legacy $S_{lp}$ retired (J14.2-011). \\
C02 \newline R & Legacy signed max drop vs symmetric range/variance/MAD and no-cue-anchored signed drop; construct stability. & G0 &
\textbf{RETIRED} as an endpoint question. Stored-value part done (J14.2-006). Teacher-forced single-string scores are diagnostics only; replaced by O1, O2 and O4. \\
C03 \newline R & Raw summed log-probability vs per-token average, value-only and value+unit; answer-length and token effects. & G0 &
\textbf{RETIRED}: a single-string teacher-forced probability is not an endpoint. Length check on legacy values in J14.3-005; candidate length matters only for O6 inside a claim type. \\
C04 \newline R & Same physical answer under several equivalent serialised answers; matched within-wrapper contrasts, not raw score pooling. & G2, G4 &
\textbf{DONE-EXPLORATORY} (X03: J14.10-005; J14.16-005). Continues in G4: two declared final-answer formats, each with its estimate and interval, and the format-dependence rule of Section~\ref{sec:gates}. \\
C05 \newline R & Chat template, prefix, assistant headers and exact token slicing; loss covers only target output tokens. & G2, G5 &
\textbf{DONE-EXPLORATORY} (J14.3-007): tokenizer audit; double BOS (Llama) and missing \texttt{Answer:} prefix found. Manual teacher-forced cases \NE{}. Required again for any O6 run. \\
C06 \newline R & Gold-string likelihood vs restricted gold-vs-decoy likelihood ratio; does a distractor gain relative probability. & G5 &
\textbf{REVISED} to O6: log-odds of the certified value against a misuse value fixed by a physical error hypothesis, only inside a declared claim type. Stage~14 $R_{alt}$ (J14.9-004) is a diagnostic. \\
C07 \newline R & Teacher-forced score vs deterministic final-answer correctness and answer flip; expose likelihood-only outcomes. & G2 &
\textbf{SUPERSEDED} by O1--O3: free-generation outcomes are primary, from sampled, not deterministic, answers. \\
C08 \newline R & Multiple decoding temperatures and seeds; answer distribution and reproducibility of behavioural effects. & G2 &
\textbf{SUPERSEDED} by the sampling rule (Section~\ref{sec:outcomes}): repeated samples per cell at model defaults; G2 estimates the self-flip floor on fresh samples. Greedy Stage~14 records are not pooled. \\
C09 \newline R & Prompt-length-matched no-cue edit vs actual cue edit; normalise against ordinary context variation. & G2, G3 &
\textbf{SUPERSEDED} by the style-matched non-numeric sentence and H0's generic-perturbation term, read against the G2 noise floor. Exploratory: J14.8-002. \\
C10 \newline R & Next-token entropy and response-token confidence vs correctness calibration; entropy not conflated with knowing the answer. & G7 &
\textbf{OPEN}. Every comparator tagged oracle or deployable; entropy built from the gold log-probability is oracle (review check C4). \\
C11 \newline R & Fixed FPR .05/.10, precision under shifted base rate and cost-sensitive utility; operational rather than ranking value. & G7 &
\textbf{REVISED}: the monitored event is an O1/O4 departure on role-by-value cells, not the $S_{lp}$ label. \\
C12 \newline C & Average of valid numerical answer paraphrases vs max; selection optimism and logsumexp semantic grouping. & G5 &
\textbf{REVISED}: answer equivalence is handled by the grading views for O1/O2; for O6 each candidate's equivalent serialisations are declared before scoring. \\
C13 \newline C & Restricted answer-distribution JSD vs generated semantic answer divergence; truncation and sampling approximation. & G3 &
\textbf{SUPERSEDED} by O4 on sampled answers (per-cell distributions, misuse-curve matches). \\
C14 \newline C & Confidence shift at the first answer token vs along the answer chain vs at the first intermediate quantity; locate the interface dependency. & G5 &
\textbf{SUPERSEDED} by claim types T1 (answer positions on a fixed trace) and T2 (certified intermediate value). \\
C15 \newline C & bf16/fp16, different batches and repeated exact prompts; numerical accuracy of score computation. & G2 &
\textbf{DONE-EXPLORATORY} in part (J14.9-004: dtype noise floor and run-to-run identity; J14.17-006). Batching untested; needed for the batched sampling contract. \\
C16 \newline C & Different no-cue anchors matched for length/style; the preferred anchor is not just another arbitrary base. & G3 &
\textbf{REVISED}: O4 compares values within a role and roles at the same value, not against one anchor; no-cue and non-numeric sentence are both reported. \\
C17 \newline C & Exclude near-zero original causal sensitivities vs include them with an absolute-effect alternative; ratio instability. & G5 &
\textbf{RETIRED}: normalised recovery is retired (Section~\ref{sec:causal}, checklist item~3); absolute effects and $I_h$ replace it. \\
C18 \newline X & Alternative proper-scoring rules and semantic choice sets on enumerable answers; sensitivity to answer-set selection. & G5 &
\textbf{REVISED}: diagnostic inside O6 only; candidate sets fixed by prespecified misuse hypotheses, never by the model's scores. \\
\end{longtable}
\endgroup

% ---------------------------------------------------------------- D
\subsection{D. Monitoring, statistical selectivity and model comparison (D01--D16)}
\label{sec:inv-D}

All sixteen D contrasts were written for a monitor of the legacy sensitivity label. That label is retired
(J14.2-011) and monitoring is a side track (G7). Each D contrast therefore survives only as \textbf{REVISED}: its
target is an O1 or O4 event on role-by-value cells, and every comparator carries an oracle or deployable tag.
Legacy-label versions count only for the Phase-I closure note (G0).

\begingroup\small
\begin{longtable}{@{}P{1.2cm}P{7.2cm}P{1.3cm}P{5.6cm}@{}}
\toprule
ID & Contrast (manipulated / read out) & Gate & Status \\
\midrule
\endfirsthead
\toprule
ID & Contrast (manipulated / read out) & Gate & Status \\
\midrule
\endhead
\bottomrule
\endfoot
D01 \newline R & Four-rendering averaged L18 family probe vs strictly one-rendering deployment probe; matched query budgets. & G7 &
\textbf{REVISED} (O1/O4 target). The legacy probe needs regenerated activations. \\
D02 \newline R & Within-family LOOCV vs the 34 withheld structural clusters and entirely unseen symbolic graph classes. & G7 &
\textbf{REVISED}. Legacy: non-activation predictors done (J14.4-009); monitor \NE{}. \\
D03 \newline R & Fitted entropy vs probe vs calibrated combined model under identical nested folds; paired-difference CI. & G7 &
\textbf{REVISED}. The legacy fitted entropy is an oracle comparator (built from the gold log-probability; review check C4). \\
D04 \newline R & Generic correctness probe plus entropy/surface vs sensitivity probe; incremental information by held-out log loss. & G7 &
\textbf{REVISED}: O1/O4 target on role-by-value cells, comparators tagged oracle or deployable; G7. \\
D05 \newline R & Embedding layer vs mid-layer vs final residual; preselected discovery layer vs best-of-all-layers selection. & G7 &
\textbf{REVISED}; the layer is preselected. \\
D06 \newline R & Digit-count, digit-identity, unit-lexicon, equation-AST and prompt-length baselines, fairly tuned. & G7 &
\textbf{REVISED}. Legacy surface baselines in J14.4-009. \\
D07 \newline R & Class prevalence matched across cue types vs natural cue mixture; pooled transfer inflation. & G7 &
\textbf{REVISED}: strata are never pooled. \\
D08 \newline R & Single prompted family vulnerability vs variant-specific event risk; no copying of the family label into variant-level claims. & G7 &
\textbf{REVISED}: the structure is the unit; claims are made per cell. \\
D09 \newline C & Logistic vs sparse probe vs shallow nonlinear MLP at controlled regularisation and capacity. & G7 &
\textbf{REVISED}: O1/O4 target on role-by-value cells, comparators tagged oracle or deployable; G7. \\
D10 \newline C & Hidden-state covariance and representation-stability uncertainty baselines, including multi-rendering costs. & G7 &
\textbf{REVISED}: O1/O4 target on role-by-value cells, comparators tagged oracle or deployable; G7. \\
D11 \newline C & Calibration shift to novel laws, domains and numeral ranges; recalibration from training data only. & G7 &
\textbf{REVISED}: O1/O4 target on role-by-value cells, comparators tagged oracle or deployable; G7. \\
D12 \newline C & Independently selected model families, reasoning-tuned checkpoints and same-backbone fine-tunes. & G3, G4, G6 &
\textbf{REVISED}: becomes the model and thinking-mode comparison of behavioural endpoints on one common item set (Section~\ref{sec:populations}); further families in G6. \\
D13 \newline C & Prompt-site cue-token vs last-prompt vs answer-prefix states; semantics and token positions matched. & G5 &
\textbf{REVISED}: probes are diagnostic only (checklist item~13); site comparison under T1/T2. \\
D14 \newline C & Probe directions trained on Cue~A tested on B/C and vice versa; family-cluster holdout and base-rate control. & G7 &
\textbf{REVISED}: transfer across the strata R/D/C/K, never pooled. \\
D15 \newline C & Probe complexity and selectivity vs shuffled labels within structural strata and artificial metadata controls. & G5, G7 &
\textbf{REVISED}: O1/O4 target on role-by-value cells, comparators tagged oracle or deployable; G7. \\
D16 \newline X & Learning curve and label efficiency over growing numbers of distinct template clusters, not repeated paraphrases. & G7 &
\textbf{REVISED}: O1/O4 target on role-by-value cells, comparators tagged oracle or deployable; G7. \\
\end{longtable}
\endgroup

% ---------------------------------------------------------------- E
\subsection{E. Direct causal localisation, intervention semantics and specificity (E01--E18)}
\label{sec:inv-E}

The October E and F groups were written around the Phase-I layer-16 head~11 (H11) and its four-head set, which
were selected on the panel they were tested on and read through normalised recovery of $S_{lp}$. In the revision
every E and F contrast runs at G5, after G4 establishes H3, H4 or H6-RD, on components found anew. Each declares its claim type
(T1/T2/T3) and reports the interaction estimand $I_h$ with the clean cue effect beside it
(Section~\ref{sec:causal}).

\begingroup\small
\begin{longtable}{@{}P{1.2cm}P{7.2cm}P{1.3cm}P{5.6cm}@{}}
\toprule
ID & Contrast (manipulated / read out) & Gate & Status \\
\midrule
\endfirsthead
\toprule
ID & Contrast (manipulated / read out) & Gate & Status \\
\midrule
\endhead
\bottomrule
\endfoot
E01 \newline R & Qwen L16 H11 only vs the original four-head set vs no intervention vs nearby layers/heads, on untouched causal families; absolute effects. & G5 &
\textbf{RETIRED} as stated (card X07, Section~\ref{sec:cards}); replaced by components found anew under $I_h$ on untouched structures. \\
E02 \newline R & Zero vs mean vs donor replacement vs projected subspace deletion; detect artifacts of the intervention distribution. & G5 &
\textbf{SUPERSEDED} by checklist items~2 and~4: resample/interchange primary, mean ablation as sensitivity, zero ablation never as evidence. \\
E03 \newline R & Same-layer, norm-matched random direction and random head; compare dose and expected change. & G5 &
\textbf{OPEN} (checklist item~9: matched random components as nulls). \\
E04 \newline R & Patch at a random prompt position, relevant variable, cue numeral, cue unit, isolation token and question token; localise source-to-target flow. & G5 &
\textbf{OPEN}. ``Isolation token'' becomes the role-relation phrase tokens of the R/D twins. \\
E05 \newline R & Matched donors: same-answer, cue-present stable, unrelated, same-structure, different-structure; specificity. & G5 &
\textbf{SUPERSEDED} by checklist item~10: within-structure, value-matched R/D donors first; cross-structure later. \\
E06 \newline R & Correct vs wrong vs stable recipient; stratify by pre-intervention behaviour and score. & G5 &
\textbf{REVISED}: recipients stratified by O1/O4 behaviour and population (Section~\ref{sec:populations}), not the legacy stable/sensitive label. \\
E07 \newline R & Cue near-to-far and far-to-near donor transplantation; predicted reciprocal intervention sign. & G5 &
\textbf{OPEN}. Value-only minimal pairs within a role; sign committed in advance (checklist item~14). \\
E08 \newline R & Knockout with and without rescue of the H11 activation; reversal of the intervention effect on the chosen endpoint. & G5 &
\textbf{REVISED}: rescue within structure and value-matched (checklist item~10), for components found anew. \\
E09 \newline R & Full-head output vs QK pattern vs V/OV path and source-specific contribution; separate head functions. & G5 &
\textbf{OPEN} (for components found anew). \\
E10 \newline R & Accuracy and decoy log-odds vs unrelated-task loss/KL; constrain off-target damage. & G5 &
\textbf{OPEN} (checklist item~6). Decoy log-odds become O6 (certified value against a prespecified misuse value). \\
E11 \newline R & Edit the irrelevant cue vs a relevant numeric input that changes $y^*$; does the lesion destroy legitimate computation. & G5 &
\textbf{SUPERSEDED} by $I_{\text{role}}$ with the R twin at the same $z$ (Section~\ref{sec:causal}). \\
E12 \newline R & Original 20 high-$S$ families vs new unbiased and near-threshold causal panels; selection effect. & G5 &
\textbf{REVISED}: the $S_{lp}$-selected panel is retired; discovery and test on disjoint structures, selected by G4 behaviour (checklist item~8). \\
E13 \newline C & Model self-repair: ablate and measure downstream compensation in MLPs and alternate heads. & G5 &
\textbf{OPEN} (checklist item~5: direct and total effects, frozen-LayerNorm control). \\
E14 \newline C & Joint knockouts of redundant head sets vs H11 alone; nonadditivity and substitute routes. & G5 &
\textbf{REVISED}: nonadditivity among components found anew. \\
E15 \newline C & Direction injection into stable recipients at several scales and both signs, with off-target outcome measures. & G5 &
\textbf{REVISED}: recipients from the reliably solved population; reported as denoising (checklist item~1). \\
E16 \newline C & Cross-family activation donation aligned by physical role, source output and tokenizer position. & G5, G6 &
\textbf{OPEN}: the later cross-structure generalisation test of checklist item~10. \\
E17 \newline C & Causal ablations during teacher-forced intermediate-quantity prediction and during generation-time answering. & G5 &
\textbf{SUPERSEDED} by claim types T2 and T3. \\
E18 \newline X & Matched learned steering vectors, gradient-based erasure, linear subspace projection and feature clamping. & G5, G7 &
\textbf{OPEN} (exploratory). \\
\end{longtable}
\endgroup

% ---------------------------------------------------------------- F
\subsection{F. Functional mechanism, circuits, synthetic controls and mitigation (F01--F20)}
\label{sec:inv-F}

\begingroup\small
\begin{longtable}{@{}P{1.2cm}P{7.2cm}P{1.3cm}P{5.6cm}@{}}
\toprule
ID & Contrast (manipulated / read out) & Gate & Status \\
\midrule
\endfirsthead
\toprule
ID & Contrast (manipulated / read out) & Gate & Status \\
\midrule
\endhead
\bottomrule
\endfoot
F01 \newline R & The head's near-match magnitude code vs the same digits with a distant numerical value; causal decoding and swap. & G5 &
\textbf{REVISED}: function test on components found anew; H1/H2 predictions on value-only and digit-only minimal pairs. \\
F02 \newline R & Unit representation vs token overlap at a fixed numeral; function-discriminating causal patch. & G5 &
\textbf{OPEN} (H2). \\
F03 \newline R & Semantic isolation role vs no-isolation semantic variant; does upstream text gate the H11 effect. & G5 &
\textbf{REVISED}: R vs D phrasings (marked and unmarked) gating components found anew; role variable by DAS (checklist item~11). \\
F04 \newline R & Attention map vs source-token value patch; no causal inference from attention visuals alone. & G5 &
\textbf{OPEN}. \\
F05 \newline R & Direct logit attribution vs true nonlinear patching; estimated vs actual intervention effects. & G5 &
\textbf{OPEN} (checklist item~7). \\
F06 \newline R & Source-to-H11 edge patch, H11-to-downstream patch, knockout/rescue; validate the directed causal pathway. & G5 &
\textbf{REVISED}: the pathway through components found anew (card X09). \\
F07 \newline R & Whole-residual patch vs path-specific patch; causal specificity and reconstruction loss. & G5 &
\textbf{OPEN}. \\
F08 \newline R & Retain-only-circuit vs remove-circuit; faithfulness/comprehensiveness, robustness to method choice, edge-count frontier. & G5 &
\textbf{REVISED}: faithfulness reported as absolute log-odds and $I_h$, never as normalised recovery (checklist item~3). \\
F09 \newline R & Interchange intervention on the high-level variable ``cue relevant?'' vs a low-level subspace; causal-abstraction agreement. & G5 &
\textbf{SUPERSEDED} by checklist item~11: role variable by DAS with interchange accuracy on held-out pairs, baselines and the subspace-illusion check; decodable is diagnostic, only causally sufficient supports a computation claim. \\
F10 \newline R & Held-out-law prediction of effect direction from the discovered function; no refitting. & G5, G6 &
\textbf{OPEN} (checklist item~14, extended to size and the second model). \\
F11 \newline C & Functionally aligned candidate pathway in a second architecture, not the same layer/head index. & G6 &
\textbf{OPEN} (card X10). \\
F12 \newline C & Supervised distributed alignment search vs sparse-autoencoder features vs neuron/low-rank baseline; validate the interchange effect. & G5 &
\textbf{OPEN} (checklist item~11 baselines). \\
F13 \newline C & Synthetic model with a planted shortcut vs clean-trained synthetic model; sensitivity and specificity of causal discovery. & G5 &
\textbf{OPEN}. \\
F14 \newline C & Near-match-only vs multi-cue cases; additive or competing routes. & G5, G6 &
\textbf{OPEN}. \\
F15 \newline C & Mechanism-guided patch or head gate vs plain cue removal and prompt editing; accuracy and compute matched. & G7 &
\textbf{OPEN} (card X11). \\
F16 \newline C & Mechanism-guided training vs random distractor augmentation vs generic fine-tune, on an entire law-held-out test. & G7 &
\textbf{OPEN}. \\
F17 \newline C & Counterfactual-consistency loss plus correctness loss vs either alone; invariance/accuracy trade-off. & G7 &
\textbf{OPEN}. \\
F18 \newline C & Adversarially generated certified cue attacks against the frozen monitor and repair, with fresh-validation transfer. & G7 &
\textbf{OPEN}. \\
F19 \newline X & Small-model scale/curriculum scan for the emergence of numeral-copy and relevance-gating pathways. & G6 &
\textbf{OPEN} (exploratory). \\
F20 \newline X & Open cross-layer transcoder attribution graph and feature editing; faithfulness under reconstruction error against an alternative-feature baseline. & G5 &
\textbf{OPEN}; a screen only (checklist item~7). \\
\end{longtable}
\endgroup

\subsection{Status summary}
\label{sec:inv-summary}

\begin{center}\small
\begin{tabular}{@{}lrrrrrr@{}}
\toprule
Group & DONE-EXPLORATORY & OPEN & REVISED & RETIRED & SUPERSEDED & Total \\
\midrule
A & 3 & 13 & 0 & 0 & 2 & 18 \\
B & 5 & 8 & 1 & 0 & 0 & 14 \\
C & 4 & 1 & 5 & 3 & 5 & 18 \\
D & 0 & 0 & 16 & 0 & 0 & 16 \\
E & 0 & 8 & 5 & 1 & 4 & 18 \\
F & 0 & 15 & 4 & 0 & 1 & 20 \\
\midrule
All & 12 & 45 & 31 & 4 & 12 & 104 \\
\bottomrule
\end{tabular}
\end{center}

``DONE-EXPLORATORY'' includes items done only in part; the cell text says which part. No item is confirmed.

\subsection{Contrasts added by the revision}
\label{sec:inv-added}

These core elements of Section~\ref{sec:framework} have no October identifier; they are registered here with
their gates.
\begin{enumerate}
\item \textbf{Role-by-value cells per stratum} (R, D, C, K where feasible; $\geq3$ values, $\geq5$ with held-out
values for the dependency curve; infeasibility recorded, never imputed). Certificates at G1, discovery at G3,
the preregistered role-by-value family at G4.
\item \textbf{Multiple phrasings} of each role relation ($\geq2$ independently validated, at least one without
marker words such as ``unrelated'' or ``separate''); effects must agree across phrasings. G1, G3, G4.
\item \textbf{Misuse response curves} $f_h(z)$ for D, C and K cells, with answer matching that keeps every match
(\texttt{UNIQUE\_MATCH}, \texttt{MULTI\_MATCH}, \texttt{NO\_MATCH}). Frozen with the item at G1; read at G3/G4.
\item \textbf{Formula hint} as a positive diagnostic (expected to help or to do nothing; never a gate). G3; also a
G0 review-derived addition.
\item \textbf{Relevance instruction} (``identify and ignore irrelevant information''), with and without a cue: G3;
also the cheapest mitigation baseline at G7.
\item \textbf{Non-physics dependency control} (GSM-DC-style problems with the same role manipulation), to tell a
general numeric or dependency phenomenon from one needing physical semantics. Small control at G3; wider
non-physics transfer at G6.
\item \textbf{O3 completion and grading status categories} (\texttt{COMPLETE}, \texttt{BUDGET\_EXHAUSTED},
\texttt{FORMAT\_INVALID}, \texttt{UNPARSEABLE}, \texttt{MULTIPLE\_CANDIDATES}, \texttt{EMPTY}) as an outcome in its
own right. Frozen at G1; measured from G2 on.
\item \textbf{Three grading views and a stratified parser audit} (last-printed-place rule, unit-convertible rule,
relative error $\leq1\%$ with dimension-checked conversion). Frozen at G1; audit at G2.
\item \textbf{Thinking on versus off and reasoning budget} within the same weights (secondary question, H5 cost via
O5). G2 calibration, G3 discovery.
\item \textbf{Two interaction estimands on every intervention}: $I_{\text{presence}}$ against the non-numeric
sentence and $I_{\text{role}}$ against the R twin, with the clean effects beside them. G5.
\end{enumerate}

% ===== s07_cards =====
\section{Twelve decisive experiment cards}
\label{sec:cards}

The October edition specified twelve cards in priority order, as minimal specifications for an engineer or
research agent. Each has a purpose, a primary outcome, the competing comparators, a no-go response and an
artifact. Those fields are kept below in condensed form. Two fields are added: \emph{Status}, from the Stage~14
record, and \emph{Revision}, from the disposition of the first revision and the revised core
(Section~\ref{sec:framework}). Where the two disagree, the Revision field governs.

\textbf{Rule kept from the October edition.} An agent must not invent a different estimand or panel mid-run. The
estimands and panels stated in a card's Revision field are frozen at G1 or when the card's gate opens; from then
on the same prohibition applies to them. A card with zero evaluable records reports \NE{}, never a zero effect.

\begin{center}\small
\begin{tabularx}{\linewidth}{@{}P{1.0cm}P{4.3cm}P{1.4cm}L@{}}
\toprule
Card & October title & Gate & Status in brief \\
\midrule
X01 & Authoritative reproduction & G0 & Done (exploratory) for committed records; full regeneration narrowed \\
X02 & Base-variant invariance challenge & G0 & Done (exploratory); no-go met; legacy label retired \\
X03 & Format/tokenization falsification & G2 & Done (exploratory) in two versions; format robustness moves into G4 \\
X04 & Physics-semantic factorial & G3, G4 & Open; becomes the role-by-value factorial \\
X05 & Structural leakage and real OOD monitoring & G7 & Partly done (exploratory); side track \\
X06 & Frozen external behavioural replication & G4, G6 & Open; internal part folded into G4, external part in G6 \\
X07 & H11 lesion/patch specificity & G5 & Retired as stated; replaced by components found anew \\
X08 & Mechanistic function discrimination & G5 & Open; gated on G4 \\
X09 & Source-to-sink circuit and rescue & G5 & Open; gated on G4 \\
X10 & Architecture functional replication & G6 & Open; gated on G4 and G5 \\
X11 & Selective mitigation challenge & G7 & Open; gated on G4 \\
X12 & External audit and claim freeze & G6 & Open; gated on G4 \\
\bottomrule
\end{tabularx}
\end{center}

% ------------------------------------------------------------------
\subsection{Card X01 --- Authoritative reproduction}
\label{card:X01}
\begin{description}
\item[Purpose.] Independent recomputation, from the committed Phase-I raw logs, benchmark manifest and
model/tokenizer checksums, of labels, panels, out-of-fold probe scores, bootstrap distributions and H11 effects;
validation of 10+ hand-selected families directly against model forward passes.
\item[Primary outcome.] Equality to the reported panel denominators, checksums and endpoint values within a
declared numerical precision.
\item[Comparators and controls.] Independent implementation; data and label lineage; no circular calculations
from summary tables.
\item[No-go response.] Block every dependent experiment; localise the fault and rerun upstream.
\item[Artifact.] Signed reconciliation ledger and replayable tests.
\item[Gate.] G0.
\item[Status.] \textbf{DONE-EXPLORATORY} for the committed-record scope (J14.1-008): the headline denominators,
labels and endpoint values reproduced from committed records, with the exceptions recorded there (one count
differs under the 0.5-nat definition; one gain is not recomputable from committed predictions; the 71\% figure
was later reconciled in J14.4-003). The forward-pass check of 10+ families and the hook check were not run, so
the P0 gate was recorded as not passed (J14.1-011). The reproduced four-head ``recovery'' is an over-reversal
of the legacy statistic, not a restoration (J14.1-009).
\item[Revision.] Full regeneration of Phase-I values is off the critical path, because the paper is withdrawn
and no new claim rests on its numbers. G0 keeps the extended parity check, plus a forward-pass spot check of 10
legacy families if a Phase-I number is ever cited again. Full regeneration is a side track if the user keeps it
(decision D1). If the parity check fails for a model or backend, no mechanism work runs on it.
\end{description}

% ------------------------------------------------------------------
\subsection{Card X02 --- Base-variant invariance challenge}
\label{card:X02}
\begin{description}
\item[Purpose.] Using all stored four-value log-probability arrays and no new GPU work, rotate the baseline
through all four renderings. Compare the legacy classification with the base-free range and with a no-cue anchor
in selected examples, across cue types and models.
\item[Primary outcome.] Reclassification matrix and model rank concordance.
\item[Comparators.] The four base choices; base-free range; no-cue anchor.
\item[No-go response.] If sensitivity depends heavily on the arbitrary variant~0, retire the reference-based
construct as a primary endpoint and reanalyse Phase~I under a preregistered replacement. Keep the history rather
than rewriting the original preregistration retroactively.
\item[Artifact.] Per-family metric multiverse with score versions.
\item[Gate.] G0.
\item[Status.] \textbf{DONE-EXPLORATORY} (J14.2-006). Of the base-0 positives, 40--49\% (Qwen) and 41--56\%
(Llama) turn negative under another base. The no-go condition is met (J14.2-008), and the legacy $S_{lp}$ label is
retired for v2 analyses on the user's confirmation (J14.2-011). The no-cue anchored contrasts needed new
evaluations and were not part of it.
\item[Revision.] The replacement is not another teacher-forced label. It is the outcome set O1--O6
(Section~\ref{sec:outcomes}), with O4 primary for the scientific question; range and anchor-matched scores are
diagnostics only. Phase~I is not re-scored under a substitute label; the G0 closure note records which
descriptive Phase-I statements survive. The 2026-06-14 preregistration block stays unedited.
\end{description}

% ------------------------------------------------------------------
\subsection{Card X03 --- Format/tokenization falsification}
\label{card:X03}
\begin{description}
\item[Purpose.] A stratified 60--100-family design with paired wrapper variants on both primary dense models:
4--6 equal-meaning output wrappers with exact answer-token spans; matched cue and no-cue.
\item[Primary outcome.] Within-family nuisance effect vs the wrapper $\times$ nuisance interaction on gold/decoy
log-odds and on the reference-free range.
\item[Comparators.] Each wrapper against the reference wrapper; tokenizer audit.
\item[No-go response.] If the wrapper dominates, any causal-mechanism interpretation is tied to output
elicitation and cannot be presented as a universal physics vulnerability.
\item[Artifact.] Tokenizer audit and hierarchical paired contrasts.
\item[Gate.] G2.
\item[Status.] \textbf{DONE-EXPLORATORY} in two versions. (a) Ten v2 structures under seven wrappers
(J14.10-005): the sign of the same-unit near-cell shift persists in every wrapper and both models. For Qwen the
\texttt{scientific} wrapper dominates under the frozen rule, so the effect size cannot be stated independently of
format. Parse rates depend strongly on format. (b) 60 stratified legacy families under four wrappers
(J14.16-005): no wrapper dominates for either model, under a looser control (no no-cue prompt). Not done:
first-token vs full-sequence score; value-only vs value-plus-unit likelihood.
\item[Revision.] Both versions are likelihood-level results; the Stage~14.3 scoring contract stays in force for
any O6 scoring. Format robustness of the behavioural endpoints moves into G4: the final-answer line is run in two
declared formats on a random half of the items, each with its estimate and interval, under the format-dependence
rule of Section~\ref{sec:gates}. Format
failures are O3 categories (\texttt{FORMAT\_INVALID}, \texttt{UNPARSEABLE}), never \texttt{WRONG}.
\end{description}

% ------------------------------------------------------------------
\subsection{Card X04 --- Physics-semantic factorial (becomes the role-by-value factorial)}
\label{card:X04}
\begin{description}
\item[Purpose.] Novel certified structures in $2\times2\times2\times$ position blocks: physical relevance
$\times$ unit compatibility $\times$ digit overlap, with proximity manipulated independently where possible.
Randomised cells blocked by family/structure, human-validated, with an exact output protocol.
\item[Primary outcome.] A preregistered interaction in generated answer selection and in gold/decoy log-odds.
\item[Comparators.] The factorial cells themselves (relevance vs copying/priming accounts).
\item[No-go response.] If there is no relevance-dependent effect, prefer a copying/priming mechanism and stop
relevance-gating claims.
\item[Artifact.] Factorial table plus held-out conditional predicted effects.
\item[Gate.] G3 (discovery on EXPLORATION structures), then G4 (one preregistered contrast on new held-out
template clusters).
\item[Status.] \textbf{OPEN}. The Stage~14 records contain no matched relevant-role control (J14.12-002); the
Stage~14 cue cells are not this factorial.
\item[Revision.] Becomes the role-by-value factorial of Section~\ref{sec:rolevalue}:
\begin{itemize}
\item \emph{Design.} Role $r\in\{R,D,C,K\}$ where feasible, crossed with value $z$ ($m\geq3$; $m\geq5$ with
held-out values for the curve) and with $\geq2$ validated phrasings per role relation, at least one without marker
words. Unit compatibility, the three proximity levels, digit overlap and position become supporting conditions, as
do the formula hint, the relevance instruction and the non-physics dependency control
(Section~\ref{sec:conditions}). A structure enters the primary contrast only with R and D cells at three values in
two phrasings.
\item \emph{Primary outcome.} O4: the G4 family F1 (D tracking loss with misuse attribution), F2 (R neglect share)
and F3 (role $\times$ value interaction on tracking), with all matches kept; $\lambda$ and $\Lambda$ are descriptive. O1--O3 and O5 are always reported. Strata are never pooled.
\item \emph{No-go, restated in the hypotheses of Section~\ref{sec:hypotheses}.} Hypotheses are scored, not exclusive. Copy tracking in every role
supports H1. Departures confined to same-unit plausible values support H2. H6-RD (correct selection between R and D) needs F1, F2 and F3 each shown
equivalent to zero (Section~\ref{sec:statistics}); H6-extended adds the same for C and K. Relevance-selection language (H3) is
used only if a D tracking loss is attributed to a prespecified misuse curve (substitute, add or subtract, in its
predicted direction; for C and K a wrong-law curve) across values.
\item \emph{Selection.} G4 confirms the role-by-value family fixed in advance; the G3 scorecard may add at most one
secondary family (Section~\ref{sec:gates}).
\end{itemize}
\end{description}

% ------------------------------------------------------------------
\subsection{Card X05 --- Structural leakage and real OOD monitoring}
\label{card:X05}
\begin{description}
\item[Purpose.] On the old 135 monitoring families, the 34 adjudicated clusters and newly authored structure
clusters: nested train-by-cluster cross-validation vs the original family LOOCV, then a final frozen external test.
\item[Primary outcome.] Paired cluster-held-out AUROC and log loss vs fitted entropy and a surface/AST baseline,
with intervals corrected for selection.
\item[Comparators.] Fitted entropy; surface/AST baseline; family LOOCV.
\item[No-go response.] If performance collapses, the monitor reads the template or editing grammar, and no
general out-of-structure deployment claim is made.
\item[Artifact.] Exact fold manifest, predictions, and a fairness/cost table.
\item[Gate.] G7 (side track).
\item[Status.] \textbf{Partly done, exploratory} (J14.4-009, J14.4-010, J14.4-011; card file in
\texttt{docs/experiment\_cards/v2/}). Clusters were rebuilt. Template-only predictors lose their family-held-out
signal under cluster-held-out folds; the fitted entropy features do not. The monitor itself under cluster-held-out
folds is \NE{}: its activations no longer exist.
\item[Revision.] A side track off the critical path. The legacy-label monitor under cluster folds belongs only to
the Phase-I closure, and only if activations are regenerated. Any new monitor targets an O1/O4 event on
role-by-value cells. Every comparator is tagged oracle or deployable: the legacy fitted entropy is oracle, being
built from the gold log-probability (review check C4).
\end{description}

% ------------------------------------------------------------------
\subsection{Card X06 --- Frozen external behavioural replication}
\label{card:X06}
\begin{description}
\item[Purpose.] An external panel from a distinct source and distinct clusters, independently authored and
verified, on two dense models. The output format is locked; reference-free and free-generation endpoints.
\item[Primary outcome.] Interval for the cue-specific effect and the model interaction, tested before any
mechanistic selection.
\item[Comparators.] No-cue and matched control conditions; the second model.
\item[No-go response.] If there is no effect, narrow the main thesis to the historic panel and investigate the
distribution mismatch rather than cherry-picking rare positives.
\item[Artifact.] Blind external results and a null-inclusive effect distribution.
\item[Gate.] G4 (internal part) and G6 (external part).
\item[Status.] \textbf{OPEN}.
\item[Revision.] The internal part is folded into G4. The role-by-value family is preregistered with its
estimands, decision rule and simulation-based sample size, on new template clusters no earlier run has touched
(INTERNAL\_CONFIRMATION). The primary and second models are those fixed by the G2 primary-model rule (planning
order Qwen2.5-7B-Instruct, Llama-3.1-8B-Instruct); the second is tested only if the primary passes (fixed sequence). Two final-answer formats are used, and the cohort is opened once. Items are validated by the
qualified LLM panel (Section~\ref{sec:llmpanel}); without it the run is labelled EXPLORATORY-REPLICATION. The external part (independently authored items, e.g.
adapted from EPC sources, each with its own certificate) stays in G6. The endpoint is O4 with O1--O3, not a
reference-free likelihood.
\end{description}

% ------------------------------------------------------------------
\subsection{Card X07 --- H11 lesion/patch specificity (retired as stated)}
\label{card:X07}
\begin{description}
\item[Purpose.] On an untouched Cue-B panel including near and far cues plus role-reversal positive controls:
H11 zero, mean, clean/corrupt donor and norm-matched controls, recording all logits and free outputs.
\item[Primary outcome.] Restoration of the absolute correct-vs-decoy logit contrast, with limited effect on
genuinely relevant edits.
\item[Comparators.] Ablation variants; norm-matched controls; relevant-edit controls.
\item[No-go response.] If only the signed ratio improves, the causal shortcut explanation is rejected.
\item[Artifact.] Per-family before/after table and a layer/head stability plot.
\item[Gate.] G5.
\item[Status.] \textbf{RETIRED as stated.}
\item[Why.] The committed Phase-I intervention records already show that H11 is the wrong object for this card.
\begin{itemize}
\item The four-head layer-16 knockout produced unclipped normalised recovery above~1 (above 100\%) in 17 of 20
families (J14.1-009; review check C8a). It over-reversed the legacy signed statistic rather than restoring a
baseline.
\item Removing head~11 alone raised the gold answer's first-token logit about equally on base and sensitive
renderings (3.76 and 3.77, 20 of 20 families each; review check C8b). The effect is cue-independent: H11 behaves
as a general answer component, not as cue processing.
\item H11 was selected on the panel it was tested on (review issue~6).
\end{itemize}
The record points the way the card's no-go anticipated. It is not a fresh test, because the panel was not
untouched. The retirement rests on construct grounds: a single-head target chosen on its own test panel and read
through a ratio of a retired label.
\item[Revision (replacement).] Components are found anew at G5, after G4 establishes H3, H4 or H6-RD, on untouched
structures and under the interaction estimand $I_h$ of Section~\ref{sec:causal}. Each experiment declares its
claim type (T1/T2/T3) and reports the clean cue effect beside $I_h$. Noising and denoising are reported
separately. Resample/interchange ablation is primary. Matched random components serve as nulls. The sign and size
on held-out structures are committed before the test. H11 has no special status. Two parts of X07 survive: the
role-reversal and relevant-edit positive controls become the R-twin control prompt of $I_{\text{role}}$, and ``record all
logits and free outputs'' becomes O6 with each candidate shown plus T3 on-policy outcomes.
\end{description}

% ------------------------------------------------------------------
\subsection{Card X08 --- Mechanistic function discrimination}
\label{card:X08}
\begin{description}
\item[Purpose.] Held-out factorial cells of numeric proximity, digits, units, isolation, target and placement.
QK/V/OV patching at source tokens with a fixed H11 receiver; preregistered opposite-sign predictions of
F-copy, F-match, F-unit and F-role.
\item[Primary outcome.] Counterfactual effect-direction accuracy on unseen structures under genuine
interventions.
\item[Comparators.] The four function hypotheses against each other.
\item[No-go response.] If no simple function predicts, widen to a distributed mechanism; do not reverse-engineer
labels from top-activating examples.
\item[Artifact.] Function-by-test prediction scorecard.
\item[Gate.] G5.
\item[Status.] \textbf{OPEN}; kept, gated on G4.
\item[Revision.]
\begin{itemize}
\item The function hypotheses map onto Section~\ref{sec:hypotheses}: F-copy onto H1, F-match and F-unit onto H2,
F-role onto H3.
\item The receiver is a component found anew (X07 replacement), not H11. Cells are role-by-value minimal pairs that
change only the role relation or only the value, token-aligned.
\item Predictions state sign and size of $I_h$ under a declared claim type and are committed before the test on
held-out structures, values and the second model (checklist item~14). Only T2 or T3 can support a computation
claim.
\item If G3 reads H1 or H2, only the optional unit-matching study runs, framed against Liu and against Owusu,
Wiegreffe and Feldman.
\end{itemize}
\end{description}

% ------------------------------------------------------------------
\subsection{Card X09 --- Source-to-sink circuit and rescue}
\label{card:X09}
\begin{description}
\item[Purpose.] A discovery graph, then a separate causal test panel; independent edge discovery and true
patching. Source-to-H11-to-downstream edits, knockout and rescue; a faithfulness--sparsity curve with alternative
ablation baselines.
\item[Primary outcome.] Significant selective mediated recovery and causal-abstraction agreement, against random
and shortcut circuits.
\item[Comparators.] Random circuits; alternative ablation baselines.
\item[No-go response.] If the source path cannot be validated, the local effect remains real at best, but the
circuit claim is removed.
\item[Artifact.] Complete edge list, paths and counterfactual intervention trace.
\item[Gate.] G5.
\item[Status.] \textbf{OPEN}; kept, gated on G4.
\item[Revision.]
\begin{itemize}
\item The path runs through components found anew, not H11.
\item ``Mediated recovery'' is reported as absolute log-odds and $I_h$, never as normalised recovery.
\item Rescue starts within structure and value-matched, between the R and D twins; cross-structure donors come
later (checklist item~10).
\item Causal abstraction means a role variable localised by DAS, with interchange accuracy on held-out pairs
against neuron and random-subspace baselines plus the subspace-illusion check (checklist item~11).
\item Attribution patching is only a screen (item~7). Selection stability under bootstrap, with discovery and test
on disjoint structures (item~8). Direct and total effects with a frozen-LayerNorm control (item~5).
\end{itemize}
\end{description}

% ------------------------------------------------------------------
\subsection{Card X10 --- Architecture functional replication}
\label{card:X10}
\begin{description}
\item[Purpose.] In Llama, on independently selected, correctly matched panels (optionally also a same-backbone
reasoning-tuned model): find a functional analogue in an exploratory partition, then freeze the site and evaluate
the identical hypothesis and endpoint on independent clusters.
\item[Primary outcome.] Sign-consistent intervention effects for the function, not for the same head index.
\item[Comparators.] The primary model's frozen function; matched random components.
\item[No-go response.] If it does not replicate, publish an architecture-specific mechanism if it is deep enough,
and discuss alternative pathways and the limited scope.
\item[Artifact.] Model-specific mechanism table and matched denominators.
\item[Gate.] G6.
\item[Status.] \textbf{OPEN}; kept, gated on G4 (and on a G5 result to replicate).
\item[Revision.] The second architecture is Llama-3.1-8B-Instruct, searched only if its own behavioural test
passed. The reasoning variant is Qwen3-8B with thinking on and off within the same weights: T1 on a fixed trace,
T3 with on-policy resampling, never a hand-edited trace as primary evidence (checklist item~12). Qwen3-32B is used
only for a frozen-function transfer test: freeze the function, not the head index.
\end{description}

% ------------------------------------------------------------------
\subsection{Card X11 --- Selective mitigation challenge}
\label{card:X11}
\begin{description}
\item[Purpose.] On external OOD and role-reversal panels, compare a targeted edit, a generic edit, cue removal,
no intervention and solver routing. Freeze the edit and gate intensity on development data; match computational
cost. Evaluate correct-answer accuracy, correct/decoy odds, cue shift and harm.
\item[Primary outcome.] A constrained utility frontier and a paired improvement in correct outcomes.
\item[Comparators.] As listed under Purpose.
\item[No-go response.] If a simpler filter wins, frame the neural intervention as a causal diagnostic, not as
deployable superiority.
\item[Artifact.] Attack/repair matrix with generated-output examples.
\item[Gate.] G7.
\item[Status.] \textbf{OPEN}; kept, gated on G4 (and on G5 for a mechanism-guided edit).
\item[Revision.] The comparators add the cheap ladder: the relevance instruction, self-consistency and a longer
reasoning budget. The primary outcome is O1 (with O2). Harm is measured on R cells as O4, since an edit must not
stop the answer responding where the quantity is relevant; cost is O5. Every comparator is tagged oracle or
deployable. Correct/decoy odds become O6 (certified value against a prespecified misuse value).
\end{description}

% ------------------------------------------------------------------
\subsection{Card X12 --- External audit and claim freeze}
\label{card:X12}
\begin{description}
\item[Purpose.] With an untouched final cohort and the full source and provenance registry: independent review of
the three strongest claims, each adversarially tested with a predeclared falsifier; a one-time reveal of the final
cohort.
\item[Primary outcome.] A precisely scoped CL1--CL6 claim matrix (Section~\ref{sec:claims}) with raw reproducible artifacts and uncertainty.
\item[Comparators.] Each claim's predeclared falsifier.
\item[No-go response.] If the claims fail, contract them transparently or keep the result as a later experiment;
never tune on the external test and call it confirmation.
\item[Artifact.] Frozen manuscript evidence matrix and reproducibility bundle.
\item[Gate.] G6.
\item[Status.] \textbf{OPEN}; kept, gated on G4.
\item[Revision.]
\begin{itemize}
\item Claim levels follow Section~\ref{sec:claims}. The external cohort consists of independently authored or
adapted items, each with its own certificate.
\item The review-derived checks are added: a mechanical claim-wording check against the claim registry, a
provenance-consistency check (authorship and validation fields against the text), and the forbidden-phrase check.
\item Agent authorship and the validation tier (LLM panel or human expert) are disclosed. Without either, results
stay exploratory.
\end{itemize}
\end{description}

% ===== s08_statistics =====
\section{Statistical design: standards for inference, not decoration}
\label{sec:statistics}

\subsection{Independent unit, estimands and non-independence}
The independent unit for a primary claim is the independently authored symbolic structure or template cluster,
not a rendered prompt, paraphrase, numeric value, phrasing or decoder sample. Values, phrasings, roles, samples and
models are repeated measures nested in structures; counting them as independent is pseudoreplication. For
generalisation to new laws or domains, the law or domain cluster is the sampling unit, or the population is
declared fixed and inference conditional on it; with only a few domains, no asymptotic $p$-values are reported as
if there were many.

\textbf{Primary behavioural estimands} (Section~\ref{sec:rolevalue}), per model and structure, averaged over
structures:
\begin{itemize}
\item \emph{Inappropriate dependence} in stratum D (separately C, K), co-primary (F1): the tracking loss
$\bar T_N-\bar T_D$ against the non-numeric sentence over all attempted samples (the loss against the no-cue prompt
beside it), which any departure outside tolerance lowers whatever its sign; the misuse shares against their rates
under $N$ attribute it (secondary family).
\item \emph{Inadequate responsiveness} in stratum R, co-primary (F2): the neglect share (R samples answering the
irrelevant-role value) against the placebo match rate in the same cell, beside $\bar T_R$.
\item \emph{Role-by-value interaction}, co-primary (F3): the role $\times$ value term of a mixed-effects logistic
model of tracking; physics predicts no dependence of tracking on $z$ in either role.
\item \emph{Descriptive}: $\lambda_R$, $\lambda_D$ and $\Lambda$ with the multimodality guard of
Section~\ref{sec:rolevalue}; they gate nothing.
\item \emph{Presence}: O1 difference between the non-numeric sentence and a D cell, for comparison with the
numeric-salience literature.
\end{itemize}
Each estimand is reported per phrasing and pooled over phrasings; a role claim needs agreement across phrasings.
Population-average and population-specific results (Section~\ref{sec:populations}) are both reported; results on
reliably solved or model-positive items are selected on the outcome and do not generalise to all items.
Eligibility (completeness, feasibility, population membership) is recorded before any intervention or cue
condition is scored, and no-cue and relevant-role (R) performance is always measured beside the cue conditions.
\textbf{Causal estimands} are the interaction $I_h$ of Section~\ref{sec:causal} under a declared claim type; an
activation intervention is not the natural indirect effect of the quantity. \textbf{Monitoring} (G7 only):
out-of-sample paired difference in a proper score at a fixed information budget, AUROC as a ranking metric.

\subsection{Intervals, tests and their assumptions}
Samples and values are first reduced to structure-level estimates; the primary analysis is then a paired analysis
of the $J$ structure-level differences: a $t$-interval with $J-1$ degrees of freedom, and a sign-flip
randomisation test whose validity rests on the symmetry of the paired differences under the null (checked and
reported). The $k\times N$ draws are never pooled as independent (Miller 2024). Hierarchical modelling (mixed-effects
logistic or linear models with structure, value and phrasing effects, simplified when singular; Matuschek et al.)
and cluster resampling (bootstrap over structures, 2,000 resamples, seed 42, as in Stage~14) are complementary
views, not a ranking: at $J$ of a few dozen, CLT and bootstrap intervals can under-cover (Bowyer et al., ICML 2025),
and a clustered Bayesian interval is reported beside them. For binary rates: denominators always, Wilson or
clustered intervals; for model comparisons on the same structures: paired analyses. Calibration claims (G7)
follow the October rules: reliability curves with bin support, ECE with fixed declared bins plus an adaptive
variant, never ``best calibrated'' from one metric.

\textbf{Cue-induced changes} are read against the self-flip floor measured in G2 (two independent draws under the
same condition), never against zero.

\subsection{Effect sizes, equivalence and the language of null results}
Every contrast has a smallest effect of scientific interest (SESOI) stated in G1 \emph{with its reason}, not chosen
for a convenient sample size, and in the units of the estimand it applies to. For the D tracking loss and misuse shares the
reason is behavioural: the share of answers that may be lost to an irrelevant quantity before it matters (planning
value 0.10 in tracking); for the R neglect share above its placebo rate the same 0.10 (a tenth of R answers ignoring a quantity that
matters); for the F3 role $\times$ value interaction an odds ratio of 1.5 per value step on tracking (planning),
because a smaller value dependence would not change a reading of any single item; for descriptive $\lambda_D$ a leak
of a tenth of the relevant-role response; for tracking shares and O1 differences the reason states which
use the result serves (a reliability claim may care about a few percentage points; a costly mechanism study needs an
effect large enough to localise). Equivalence is claimed only by two one-sided tests within the SESOI, with
adequate power: with $\omega$ the standard deviation of the per-structure difference, about
$8.6\,\omega^2/\Delta^2$ structures give 80\% power when the true effect is zero; at $\omega=0.2$, $\Delta=0.10$
needs about 34 structures and $\Delta=0.05$ about 137. If the sample cannot support the margin, the outcome is
INCONCLUSIVE, and the sample size, not the margin, is what G1 revises before confirmation. Four outcomes are kept
distinct in all text: an effect (interval excludes zero and extends beyond the margin), a trivial effect (interval
excludes zero but lies inside the margin), equivalence (interval inside the margin and including zero), and
inconclusive (none of these). ``No effect'' is never written for an inconclusive result.

\textbf{Rare categories.} Misuse-curve and copy matches may be rare (Stage~14.20 had 2 wrong answers in 128 cells).
Their rates are estimated per structure with a Jeffreys Beta prior (zero counts give a finite rate with a
proper interval), combined across structures with a Beta-Binomial hierarchical model, and reported with
denominators; a contrast between rare rates is never read from a structure-pooled count.

\subsection{Multiple comparisons and research degrees of freedom}
\begin{itemize}
\item G3 is estimation; its scorecard is not a test.
\item G4 tests the role-by-value family as three co-primary contrasts, Holm-corrected at family-wise
$\alpha=0.05$: (F1) \emph{inappropriate dependence} in D, the tracking loss $\bar T_N-\bar T_D$ against the
non-numeric sentence (one-sided, since a departure in any direction lowers tracking); (F2) \emph{inadequate
responsiveness} in R, the neglect share against its placebo rate (one-sided); (F3) the \emph{role-by-value interaction}, the role $\times$ value term of a
mixed-effects logistic model of tracking (two-sided: physics predicts tracking that does not depend on $z$ in either
role). For H6-RD each of F1--F3 is instead tested for equivalence within its SESOI. $\Lambda$ and $\lambda$ are
reported as descriptive summaries and gate nothing. The misuse shares that attribute an F1 loss to particular
misuse curves are a secondary family, Holm-corrected across curves and tested in each curve's predicted direction,
only after F1 passes. The second model follows by fixed sequence after the first. At most one
secondary family chosen in G3 is tested after the primary family, Holm-corrected, and labelled as chosen on
EXPLORATION data.
\item Prespecified secondary scans use Benjamini--Hochberg and stay exploratory.
\item Selection over layers, heads, sites, positions, phrasings, metrics and thresholds is handled by discovery on
one set and a single evaluation on an independent set, never by reporting the best of many tests.
\item If queue delays force G4 into batches, one interim analysis at half the sample with Lan--DeMets
O'Brien--Fleming spending may be preregistered.
\end{itemize}

\subsection{Power and precision protocol}
Before G3 and again before G4, simulate from the G2 estimates (structure, value and phrasing variance components,
$\omega$ of the cue difference, within-item Bernoulli variance, which more samples reduce and between-structure
variance does not; Wang 2026), for $J\in\{24,40,60,80\}$ complete role-by-value structures and
$k\in\{3,5,8\}$, at least 1,000 simulated datasets each analysed exactly as planned, on simulated null data as well as effect data so
that the type-I error of the chosen analysis and resampling unit is shown to be calibrated, for the tracking loss, the role-by-value interaction
and the misuse, neglect and copy shares, which may be rare (Stage~14.20 had 2 wrong answers in 128 cells). The
simulation includes a scenario with 20--40\% loss to infeasible cells, validation failures, parse failures and
budget exhaustion. Analytic bracket: $J\approx7.85\,\omega^2/\delta^2$ for 80\% power at two-sided $\alpha=0.05$
(Miller 2024), about $6.2\,\omega^2/\delta^2$ for the one-sided G4 tests, with $\omega$ the standard deviation of the
per-structure value \emph{of the estimand being tested} (estimated in G2, not borrowed from a no-cue difference);
with $\omega=0.25$, $\delta=0.10$ needs about 49 (two-sided) or 39 (one-sided) structures and $\delta=0.15$ about 22
or 17. More samples per item cannot
substitute for more independent structures. If errors are genuinely rare, the legitimate choices are a harder,
independently justified regime or the study of correct selection (H6); a continuous diagnostic (O6) is never
promoted to proof of an error because the categorical endpoint lacks power. The dataset is built to 250 and then 500
structures (Section~\ref{sec:datasetplan}); the simulation decides how many of them each gate opens, and the sealed
remainder is kept for later gates.

\subsection{Negative and missing result policy}
A zero-row job, failed hook, timeout, budget-exhausted or format-invalid response, infeasible cell, ineligible item
or untested control is \NE{}, not a null effect. A missingness codebook records reason, first seen, planned run,
attempts, eligible population and intended analysis. Exclusions and completeness are reported by condition, role,
value, phrasing and model. Every primary contrast is reported three ways: all attempted samples with non-complete
records as ``not correct'' (O1), complete-case, and with budget-exhausted records counted as wrong (a bound; answers
at the length limit are almost always wrong in published reasoning evaluations, Hochlehnert et al.). Data-dependent
missing values are never replaced by zero or negative labels.

% ===== s09_compute =====
\section{Compute, human effort, storage and reproducibility architecture}
\label{sec:compute}

This section keeps the October edition's requirements that are still valid (measure before budgeting, data
contracts, a conceptual preflight, storage discipline, independent human validation) and replaces its generic
seven-step compute ladder by a budget staged on the revised gates (Section~\ref{sec:gates}). Tags: \tg{V} checked
in repository files or job records; \tg{ENV} measured cluster or stack constraint; \tg{I} planning arithmetic or
design choice. Every figure below that is not tagged \tg{V} or \tg{ENV} is a planning number, not a result.

\subsection{Environment as it stands (11 November 2026)}
\label{sec:compute-env}

\textbf{Hardware and scheduler} \tg{ENV}. The \texttt{pabitra} account is shared with other projects, and their
jobs count against the same per-user QoS caps: H100 1 GPU, 12 CPU, 300G; A100 2 GPUs; H200 2 GPUs, 12 CPU, 300G.
\texttt{sbatch --test-only} does not model these caps, so the account's other jobs (\texttt{squeue -u pabitra})
are added by hand before a partition is chosen; live Slurm output outranks any dated snapshot. Measured H100 queue
waits were about 38 hours (job 417416) to about 6 days (job 425243). Weight loading from \texttt{/scratch} took
about 7 minutes for Qwen3-32B (job 417416) and has taken up to about half an hour for 7--8B models; walltime is
budgeted for it. The October edition's warning stands: throughput is measured on Sharanga, never invented.

\textbf{Storage roots} \tg{V}. Models live under \texttt{\$\{PHYSMON\_MODEL\_ROOT\}} (default
\texttt{\$SCRATCH/models}, i.e.\ \texttt{/scratch/pabitra/models}), resolved at load time; activations, Slurm logs
and large intermediates live under \texttt{/scratch/pabitra/physmon}. The October 2026 scratch wipe destroyed the
Phase-I activation tensors; the analyses that depended on them are \NE{} and are not reconstructed from
stand-ins. Scratch is treated as non-reconstructable runtime storage, never as an archive.

\textbf{Software provenance} \tg{V}. Pinned model manifests (revision, per-file sha256, architecture and
compatibility checks) are in \texttt{docs/model\_manifests/}; the registry is \texttt{docs/model\_registry.yml}.
Job templates set \texttt{PYTHONNOUSERSITE=1} so that no user-site package leaks into the pinned environment
(\texttt{docs/environment\_lock.yml}; TransformerLens 3.4.0). \texttt{gpunode6} has no usable \texttt{git}, so
every job carries \texttt{SUBMIT\_COMMIT} (the full SHA from the submission node) and refuses to run without a
commit; all jobs of one study step run from one frozen commit. Every accepted job carries an
\texttt{AUTHORIZATION\_REF}; each job still needs the user's acceptance of its exact command.

\subsection{Models}
\label{sec:compute-models}

\begin{center}\small
\begin{tabularx}{\linewidth}{@{}P{3.4cm}P{3.6cm}L@{}}
\toprule
Model & Role & Status \tg{V} \\
\midrule
Qwen2.5-7B-Instruct & First in the planning order for primary behaviour and mechanism (the primary model is fixed by the G2 rule) & Restored after the wipe, revision pinned, manifest
(\texttt{qwen\_primary.json}); HF-versus-TransformerLens parity passed on the smoke prompt (job 399217) \\
Llama-3.1-8B-Instruct & Second family & Restored, pinned, manifest (\texttt{llama\_primary.json}); parity passed
(399217) \\
Qwen3-8B (optionally 14B) & Thinking on/off within one set of weights & Optional, open decision D3; not
downloaded \\
Qwen3-32B, thinking & Scale and ceiling check on a hard subset only & Present on scratch (bf16); \textbf{revision
not pinned, no manifest}; runs 417416, 425243, 435605 carry \texttt{model\_commit NOT\_RECORDED}; TransformerLens
unverified \\
Qwen3.5-397B NVFP4, DeepSeek-R1-distill, Mistral, Qwen2.5-32B judge & Not planned & Not needed for the primary
question; no commitment to the 397B model (as the October edition advised); never used in a recorded result \\
\bottomrule
\end{tabularx}
\end{center}

Before any further Qwen3-32B run its revision is identified and a manifest written. If the scratch file hashes
match a published revision, that is recorded as an after-the-fact identification, not as a pin at run time.
Qwen3-32B is used only in bf16 (TransformerLens 3.4 \texttt{boot()} defaults to float32), through plain forward
hooks or the TransformerBridge, with exact patching only, after a parity test.

\textbf{Generation engine (open decision D4).} Plain Hugging Face \texttt{generate} gave Qwen3-32B about 20 tokens/s
aggregate with no gain from batch 2 to batch 3, far below the memory-bandwidth bound. KV memory does not explain
this: at about 256~KiB per token, the 10--18~GB left after the bf16 weights hold roughly 40{,}000--70{,}000
tokens, i.e.\ several 6{,}000-token sequences. A faster engine is therefore expected to help, but is not measured
here. If vLLM is approved it is installed in a separate \texttt{uv} environment (its wheels bundle their own
torch), after a check of the cluster's CUDA driver, followed by a logit and greedy-output parity check against the
pinned Hugging Face stack on fixed prompts and a short on-cluster timing job, before any study run. Generation and
analysis are split: vLLM generates and saves token ids; activations and every teacher-forced quantity come from a
Hugging Face teacher-forced forward pass over the saved sequences, so hooks never run inside vLLM. Without vLLM
the plan runs on the 7--8B models.

\subsection{Measured rates}
\label{sec:compute-rates}

\begin{center}\small
\begin{tabularx}{\linewidth}{@{}P{4.2cm}P{2.2cm}L@{}}
\toprule
Configuration & Job & Measured \tg{V} \\
\midrule
Qwen2.5-7B, reasoning, 2{,}048 tokens, \textbf{greedy, unbatched}, HF & 401357 & 12.6~s per record \\
Llama-3.1-8B, same & 401357 & 14.7~s per record \\
Qwen3-32B thinking, HF bf16, 1 H100, batch 2, 8{,}192 tokens & 417416 & 20.8 tokens/s aggregate; median 4{,}983
tokens per record; 5 of 20 truncated \\
Qwen3-32B, batch 3, 12{,}288 tokens & 425243 & 20.2 tokens/s (batch 3 gave no gain); median 5{,}242 tokens; 9 of
78 truncated \\
Qwen3-32B, batch 3, cue cells & 435605 & 19.8 tokens/s; mean 5{,}632 tokens per record; 12 of 128 truncated \\
\bottomrule
\end{tabularx}
\end{center}

The 7--8B rates were measured with greedy, unbatched decoding; the revised contract samples $k$ times per cell
and batches (Section~\ref{sec:outcomes}), so these rates are an upper bracket for planning, to be replaced by the
G2 timing. Qwen3-32B runs at about 14 records per hour; harder items truncate more (5 of the 7
upper-university Stage~A problems).

\subsection{Staged budget for the revised gates}
\label{sec:compute-budget}

Generations per model $=$ structures $\times$ cells per structure $\times$ $k$ samples. A ``core cell'' below is
one role-by-value cell (role $\times$ value $\times$ phrasing); the minimum complete structure has R and D at three
values in two phrasings, i.e.\ 12 cells, plus the no-cue base and the non-numeric sentence: 14 cells. An added
ablation condition is assumed to be crossed with three values and two phrasings (6 cells). The HF column multiplies by the measured 12.6--14.7~s per record; the vLLM
column is \textbf{an estimate scaled from published single-H100 benchmarks} (thousands of output tokens/s for
7--8B models, hundreds for 32B in bf16; the first revision estimated about 1~hour for 5{,}120 7--8B
generations), not a measurement, and is replaced by the on-cluster parity and timing job.

\begin{center}\footnotesize
\begin{tabularx}{\linewidth}{@{}P{2.6cm}P{3.9cm}P{1.9cm}P{2.3cm}P{1.6cm}L@{}}
\toprule
Gate step & Planning size & Generations per model & 7--8B HF (measured rate) & 7--8B vLLM (estimate) & Qwen3-32B \\
\midrule
G2 calibration (difficulty band) & 130 EXPLORATION structures $\times$ $k_{sel}$ 4, then about 60 $\times$ $k_{est}$ 4, then about 20
structures $\times$ 7 cells (R and D at three values in one phrasing, non-numeric sentence) $\times$ 4 & about 1{,}320 & about 4.6--5.4~h & under 0.5~h & not planned \\
G2 design-validation pilot & 8--12 structures $\times$ 14 cells $\times$ $k$ 4--8 & 448--1{,}344 &
about 1.6--5.5~h & under 0.5~h & not planned \\
G3 behavioural discovery (core) & 24--40 complete structures $\times$ 14 cells $\times$ $k$ 3--8; plus a small
non-physics control sized in its card & 1{,}008--4{,}480 & about 3.5--18~h, split into jobs under 10~h & about
0.2--1~h & excluded \\
G3 targeted ablations (only after the core result) & per ablation condition: 24--40 structures $\times$ 6 cells
$\times$ $k$ 3--8 & 432--1{,}920 per condition & about 1.5--7.8~h per condition & under 0.5~h per condition &
excluded \\
G4 confirmation & $J$ from the power simulation (one-sided bracket $J\approx6.2\,\omega^2/\delta^2$: about 17--39
for $\omega=0.25$, $\delta=0.15$--$0.10$; planning 50 to absorb losses) $\times$ 14 cells $\times$ $k$ 8, plus the
second answer format on half of the structures & about 8{,}400 at $J=50$ & about 29--34~h in 4 jobs & about 2~h & hard
subset only \\
Qwen3-32B hard-subset check & 20 structures $\times$ 4 conditions $\times$ $k$ 4 & 320 at 8{,}000--12{,}000 tokens
(2.6--3.8~M tokens) & -- & -- & HF about 35--55~h (4--6 jobs); vLLM bf16 several hours (estimate) \\
G5 mechanism & items $\times$ sites $\times$ donor--recipient pairs & card-specific & hours (HF teacher forcing) &
n/a (no hooks in vLLM) & exact patching on a small set, after parity \\
G6, G7 & sized after G4 & -- & -- & -- & -- \\
\bottomrule
\end{tabularx}
\end{center}

\textbf{Rules for the budget.} (1) Jobs are sized to finish in under 10 hours with an internal time budget (the
runner stops cleanly; unrun records are \NE{}), because 1--6-day queue waits make long jobs that die at walltime
expensive; a step that exceeds that is split into condition blocks from one frozen commit. (2) Sampled and greedy
records are never pooled. (3) Qwen3-32B enters only as a hard-subset check; without vLLM it is run, if at all, on
that small subset. (4) The thinking budget is part of the contract (12{,}288 tokens truncated 9--12\,\% of
Qwen3-32B records), and truncation is an outcome (O3), not a reason to drop a record. (5) The October edition's
behavioural cost model remains the template: structures $\times$ values $\times$ phrasings $\times$ cells $\times$
models $\times$ (scored forwards $+$ generated samples), with the output-format wrappers recorded separately.

\subsection{Human effort}
\label{sec:compute-human}

Independent physics review (block B4) is the binding non-GPU cost; until human reviewers are available, the
qualified LLM panel of Section~\ref{sec:llmpanel} provides the interim tier (its models are open decision D2). Generating and certifying
structures is agent and CPU work; review is not. Planning figure: 15--20 minutes per structure per reviewer.

\begin{center}\small
\begin{tabularx}{\linewidth}{@{}P{3.6cm}P{2.4cm}L@{}}
\toprule
Step & Structures & Review effort per reviewer per pass (planning) \\
\midrule
G2 design-validation pilot & 8--12 & about 2--4~h; the measured time per structure replaces the planning figure
(this replaces the October edition's 20-family adjudication-time pilot) \\
G3 behavioural discovery & 24--40 & about 6--13~h if reviewed; without review G3 is exploratory by design \\
G4 confirmation & about 50 opened (from power) & about 12--17~h for two-pass human review of 50, plus adjudication; for now
replaced by the qualified LLM panel (below), human hours deferred to the later accuracy study \\
Dataset admission (M250, M500) & 250, then 500 & LLM panel on every structure and phrasing (Section~\ref{sec:datasetplan}); human review only in the later accuracy study \\
\bottomrule
\end{tabularx}
\end{center}

Until human reviewers are available, G4 items are validated by the qualified LLM panel
(Section~\ref{sec:llmpanel}): roughly 2--3 judge runs of two passes per structure, a few GPU-hours locally or an API
budget, plus recorded adjudication of split verdicts (adjudicated items reported separately). Without a qualified panel, G4 is reported as
\texttt{EXPLORATORY-REPLICATION}. The authoring agent's own review is not validation, and agent authorship and LLM
validation are disclosed. The October edition's staffing roles are kept (hypotheses and claim ledger; generator and symbolic
verifier; instrumentation and reproducibility; validator coordination; an independent reviewer who attacks
leakage and causal interpretation); roles can be combined, but independence checks cannot all be done by the
agent that wrote the data. Its 400-structure, 800-judgment human validation budget is replaced by the dataset plan of
Section~\ref{sec:datasetplan}: 250 and then 500 structures, each validated by the qualified LLM panel, with human
review limited to the later accuracy study of the panel.

\subsection{Storage}
\label{sec:compute-storage}

Kept from the October edition: persist normalised per-structure, per-cell raw records and complete model
provenance in the repository; store activation caches transiently on scratch with an object index and hashes;
keep only discovery-critical activation slices plus reproducible extraction scripts; use safetensors, zarr or
HDF5 (no pickled code in shared artifacts); pin dtype, shape and named hook coordinates; log when a cache came from
an earlier model revision; make every run restartable and idempotent (raw records flushed after every batch,
\texttt{run\_summary.json} with a sha256 of the records, used output directories refused).

Activation memory is recomputed for current models from
$n_{\text{items}}\times n_{\text{positions}}\times n_{\text{layers}}\times d_{\text{model}}\times$ bytes, not
copied from Phase~I. Illustration \tg{I}: one residual-stream site in bf16 is $3584\times2\times28\approx196$~KiB
per token for Qwen2.5-7B and $4096\times2\times32=256$~KiB per token for Llama-3.1-8B, i.e.\ roughly 0.4--0.5~GiB
for one 2{,}048-token trace at every layer. Reasoning traces therefore make all-position caching impractical;
G5 caches only the positions and sites a card names, recomputing from saved token ids when needed.

\subsection{Machine-readable data contracts}
\label{sec:contracts}

\textbf{Every raw record} carries the October fields still in use: \texttt{run\_id}, \texttt{phase\_id} (now the
gate), \texttt{exploratory\_or\_confirmatory}, \texttt{structure\_uuid}, \texttt{template\_cluster\_id},
\texttt{symbolic\_program\_hash}, \texttt{cue\_type}, \texttt{cue\_semantic\_role}, \texttt{cue\_numeric\_value},
\texttt{unit\_type}, \texttt{numeric\_proximity}, \texttt{digit\_overlap}, \texttt{cue\_position},
\texttt{isolation\_condition}, \texttt{model\_id}, \texttt{model\_commit}, \texttt{tokenizer\_commit},
\texttt{chat\_template}, \texttt{answer\_wrapper\_id}, \texttt{generated\_answer\_raw}, \texttt{parsed\_answer},
\texttt{parse\_status}, \texttt{correctness}, \texttt{seed}, \texttt{dtype}, \texttt{wall\_time},
\texttt{artifact\_checksum}. The October likelihood fields (\texttt{target\_answer\_ids},
\texttt{target\_span\_start/end}, \texttt{gold\_logprob\_raw/per\_token}, \texttt{alternative\_scores}) are kept
only on records that compute the O6 diagnostic.

\textbf{Added by the revised framework:} \texttt{stratum} (R, D, C, K), \texttt{role\_relation\_id},
\texttt{value\_index} and value $z$, \texttt{phrasing\_id} and \texttt{marker\_word\_present},
\texttt{certified\_response} $y^*_r(z)$, \texttt{tracking} $T$, \texttt{feasibility\_status},
\texttt{misuse\_matches} (\texttt{UNIQUE\_MATCH}/\texttt{MULTI\_MATCH}/\texttt{NO\_MATCH} with every matched
hypothesis), \texttt{completion\_status} (O3), \texttt{thinking\_closed}, \texttt{generated\_tokens},
\texttt{token\_budget}, \texttt{relative\_error}, the three grading views and the per-view
\texttt{source\_ambiguous} flag, \texttt{unit\_stated\_in\_prompt}, sampling parameters (temperature, top-$p$,
top-$k$, sample index), \texttt{contract\_version}, \texttt{data\_role}, \texttt{population},
\texttt{submit\_commit}, \texttt{model\_manifest\_sha256}, \texttt{engine} and its version (HF or vLLM), and
\texttt{oom\_fallback}; \texttt{created\_by}, \texttt{verified\_by} and \texttt{admission\_category} on every item.

\textbf{Causal records} add the October fields (\texttt{receiver\_id}, \texttt{donor\_id}, \texttt{edit\_site},
\texttt{edit\_semantics}, \texttt{scale}, \texttt{clean/corrupt\_definition}, \texttt{hook\_tensor\_shape},
\texttt{absolute\_Y\_control}, \texttt{absolute\_Y\_intervened}, \texttt{KL\_damage},
\texttt{model\_positive\_eligibility}, recipient generation outcomes) plus \texttt{claim\_type} (T1, T2, T3;
Section~\ref{sec:causal}), the declared metric, the clean cue effect beside the interaction $I_h$, the ablation
kind (resample, interchange, mean), noising or denoising, and the matched-random-component null. Every derived
table names its source manifest and calculation commit.

\textbf{Continuous-integration checks} (October list, kept): invariant answer across nuisance edits; correct
answer change under certified relevant edits; only the declared spans mutated; exact-duplicate and cluster-split
integrity; no train/test structure or cluster collisions; leakage-free preprocessing; out-of-fold scores;
identical denominators for paired models; no zero-row result read as a zero effect; hooks removed after each
run; no gradients accumulated; no confirmatory run with a development-fitted threshold. \textbf{Added:} every D,
C, K cell certified invariant and every R cell certified responsive over its declared range; feasibility recorded,
never imputed; phrasings without marker words present; a record not \texttt{COMPLETE} never graded
\texttt{WRONG}; greedy and sampled records never pooled; every model in a run has a manifest, or its records carry
\texttt{model\_commit NOT\_RECORDED}.

\subsection{Preflight skeleton}
\label{sec:preflight}

Conceptual validation contract, not a runnable implementation; it extends the October skeleton with the
manifest, engine-parity and feasibility steps.

\begin{verbatim}
# Conceptual validation contract (revised); not runnable code.
manifest = freeze_structural_splits(certified_role_by_value_structures)
assert no_cluster_overlap(manifest.discovery, manifest.confirmation)
assert no_cluster_overlap(manifest.confirmation, manifest.external)
record_feasibility(manifest)            # R, D, C, K per structure; never imputed
for model in selected_models:
    require_pinned_manifest(model)      # docs/model_manifests/, else NOT_RECORDED
    validate_native_chat_template(model)
    validate_hook_parity(model, sites=planned_sites, length=planned_trace_length)
    if engine(model) == "vllm":
        validate_engine_parity(model, reference="pinned HF stack")
    run_timing_job(model)               # replaces planning rates
for model in selected_models:
    for group in manifest.discovery:    # design-validation pilot first, then G3
        records = generate_all_preregistered_cells(model, group, k_samples)
        assert certified_invariance_and_response(records.certificates)
        emit_immutable_raw_records(records, submit_commit)
# Freeze choices using discovery ONLY.
# G4 family is fixed in G1; discovery may add one secondary family
protocol = freeze_g4_family(manifest.discovery)
power = simulate_power(protocol, variance_components_from_G2)
write_preregistration(protocol, power, git_sha=submit_commit)
# Confirmation panel is loaded ONCE by a frozen script.
result = evaluate_frozen_protocol(protocol, manifest.confirmation)
report_cluster_robust_effects_and_all_units(result)   # NOT_EVALUATED kept apart
# Unsuccessful results never silently reroute to a new confirmatory test.
if not result.passes_prespecified_gate:
    close_claim_and_open_new_discovery_branch(protocol, result)
\end{verbatim}

% ===== s10_conditional =====
\section{Conditional decisions: how to react if the evidence is unfavourable}
\label{sec:conditional}

The October edition's decision table is kept where its rule still applies, re-indexed to the revised gates
(Section~\ref{sec:gates}); rows that named the retired head-11 starting point or the retired signed likelihood
label are generalised to any component or likelihood metric. The revision's contingencies follow in a second
table. Every row is decided before the data it applies to are seen.

\subsection{Decision rules kept from the October edition}
\label{sec:conditional-kept}

{\small
\begin{longtable}{@{}P{3.3cm}P{3.5cm}P{4.2cm}P{4.0cm}@{}}
\toprule
Gate and observed outcome & Scientific interpretation & Required action & Prohibited reaction \\
\midrule
\endhead
\bottomrule
\endfoot
G0: raw outputs not reproducible & Evidence chain broken, or model or runtime drift & Fix code and provenance,
rerun upstream and invalidate dependents & Continue adding components or plots around unverified metrics \\
G0/G2: a score changes with an arbitrary base or anchor (the Phase-I signed label) & Label not
indexing-invariant & Use the retired label only as a historical record; anchor-free endpoints (O1--O4) &
Relabel retrospectively and say the original preregistration always meant the new metric \\
G2: wrapper or tokenisation explains an effect & String-scoring artifact or a strong interface interaction &
Restrict the claim; study interface sensitivity in the methodology paper (Section~\ref{sec:papers}) & Assert a
physics reasoning failure from gold log-probability alone \\
G2: generated answers stable, $p(y^*)$ drifts & Likelihood-only phenomenon & Report as an O6 diagnostic;
focus on reliability or calibration, not on answer corruption & Claim answers are corrupted or repaired by a
likelihood change \\
G1/G4: independent review rejects the semantics of items & Construct invalid for the affected items & Redesign
and revalidate; rerun affected findings & Treat a solver or free-symbol check as sufficient evidence of
natural-language correctness \\
G3: no role-by-value interaction & Physical-role selection not supported for that stratum & Test the
numeral-presence, plausibility and generic-perturbation alternatives (H1, H2, H0) & Add post-hoc cue types until
the desired sign appears \\
G3: effect in one stratum or condition, absent in others & Heterogeneous phenomenon & Deep-dive the stratum
that shows it; others become scope tests & Pool strata or conditions to claim a universal mechanism (strata are
never pooled, Section~\ref{sec:strata}) \\
G5: effect only under zero or mean ablation, none under resample or interchange & Off-distribution lesion or
compensation & Diagnose the lesion; withhold a causal-function claim & Call the component necessary in all
circumstances \\
G5: component effect only on gold-string likelihood (T1 metric) & Causal contribution to the scoring
distribution & Pursue a readout account, or test T2/T3 metrics & Claim an improved physics computation \\
G5: no single-feature account & Distributed or multiple mechanisms possible & Test subspaces, several components,
edge discovery; narrow the scope & Hand-label top-activating examples and declare mechanistic understanding \\
G5: path attribution fails rescue & Proposed route wrong or insufficient & Redefine on discovery data or fall
back to localisation & Publish attribution arrows as established information flow \\
G5/G6: second model uses a different component & Function may still recur & Find an analogue with independent
discovery and confirmation, or narrow the claim & Equate a same-layer miss with architectural absence \\
G6: no external functional replication & Architecture- or benchmark-specific mechanism & Limit generalisation;
strengthen the single-model case study and state the boundary & Keep universal LLM language \\
G6: a monitor fails cluster-held-out folds (side track) & Template recognition more plausible than a general
monitor & Drop deployment claims; new representation only after fresh discovery & Quote the best leave-one-family
score as out-of-distribution performance \\
G6: a monitor matches entropy under an adequately powered test & Hidden states add little for this task &
Report the negative; no deployment claim & Read non-significance as superiority or as equivalence without a
proper test \\
G7: a cheap filter or the relevance instruction dominates & Component-level repair not practical & Treat the
edit as a scientific intervention, not a product method & Exclude the cheap baseline from the comparison \\
G6: final external set contradicts the preferred claim & Mechanism or structure family not universal & Report
the boundary, stop tuning, collect a new confirmation set for a revised claim & Reuse an opened external cohort
until it passes \\
\end{longtable}
}

\subsection{Contingencies added by the revision}
\label{sec:conditional-new}

{\small
\begin{longtable}{@{}P{3.6cm}P{5.4cm}P{6.0cm}@{}}
\toprule
Situation & Required action & Prohibited reaction \\
\midrule
\endhead
\bottomrule
\endfoot
A model is outside the difficulty band after harder or easier variants & Report the model as uninformative at
that budget & Change the band after seeing results \\
No qualified LLM panel and no human reviewer & Results stay exploratory; G4 is \texttt{EXPLORATORY-REPLICATION};
any manuscript says so, with the agent-authorship disclosure & Call agent verification independent validation \\
LLM panel qualified, no human review yet & G4 is ``confirmatory under LLM-panel validation''; disclosed; human study
of panel accuracy planned (G6/P11) & Call panel validation expert validation \\
vLLM not approved (D4) & Run on the 7--8B models with Hugging Face; Qwen3-32B only for a small hard-subset
check, if at all & Extrapolate Qwen3-32B behaviour from single-sample Stage~14 runs \\
Role-by-value design infeasible for many structures & Run the role contrast on the feasible subset with its
count stated; calibration counts complete structures only & Edit structures after G1 to force feasibility, or
impute missing cells \\
Queue starvation & CPU work (G0, G1, foils, cards, new structures, analysis scripts tested on synthetic records)
continues; jobs stay under 10 hours & Enlarge jobs to save queue entries at the cost of losing work at walltime;
touch other projects' jobs \\
Errors are rare (ceiling) & Study correct selection (H6) and its mechanism, or move to a harder regime that is
justified independently of the observed error rate (difficulty, reasoning budget, stratum) and fixed before its
data & Tune items or cues until errors appear; promote the O6 candidate log-odds diagnostic to proof of error \\
A cell has zero evaluable records & Report it as \NE{} with the reason & Read it as a zero effect \\
Results inconclusive (interval includes both zero and the equivalence margin) & Report as inconclusive with the
achieved precision & Report inconclusive as H0; equivalence must be shown against a margin fixed before the
data \\
H0 shown with equivalence bounds, or H5 (cost only) & A certified null or a cost finding is a result; write it up
as such (Section~\ref{sec:papers}) & Rescue it with new endpoints chosen after the data \\
Non-physics control shows the same effect & The phenomenon is a general numeric or dependency effect, not one that
needs physical semantics & Present it as physics-specific \\
\end{longtable}
}

\textbf{Stop--pivot hierarchy} (October edition, kept): credibility of the measured effect outranks benchmark
size; causal interpretation outranks component count; reproducible out-of-cluster prediction outranks peak
discrimination scores; selective correctness improvement outranks ratio recovery. The plan is adaptive about which
branch is investigated, not about rewriting confirmatory statistics after outcomes.

% ===== s11_literature =====
\section{Literature-grounded positioning and novelty test}
\label{sec:literature}

The October map of sources to design (its entries L01--L25 remain in the bibliography,
Section~\ref{sec:bibliography}) is kept. This edition adds the work that appeared or was found by the searches of
11 November 2026 and restates the novelty claim more narrowly. Absence of a related work from a selective search
is not proof that none exists; every source must be read in full before it is used to claim novelty; ``first to
show'' is never written unless verified against the full adjacent literature. Venue status was established by
opening the venue or arXiv page; ``per arXiv'' means only the arXiv comment claims the venue.

\subsection{Closest work and what PhysMon must add}

\begin{center}\small
\begin{longtable}{@{}P{4.0cm}P{5.4cm}P{6.4cm}@{}}
\toprule
Work & What it already contributes & What PhysMon must add to be distinct \\
\midrule\endhead
PRISM-Physics (Zhao et al., ICLR 2026) & Physics solutions as causal DAGs of formulas; rule-based symbolic formula matching; process-level scoring of intermediate steps; human-aligned evaluation & Controlled interventions on a quantity's physical role and value, with certified response; internal causal analysis. Using verified physics dependencies or scoring intermediate steps is not itself new \\
GSM-DC (Yang et al., EMNLP 2025) & Symbolic dependency graphs with distractor nodes off the solution path; path-selection and arithmetic errors grow with depth & Physical laws with strata D/C/K (disconnected, cancelling, conditional), role-by-value response curves and role-swapped controls; not another distraction benchmark \\
Four-stage decomposition (Qu and Gomes, Findings of EMNLP 2026) & Mechanistic account of GSM-NoOp failure in Llama-3.3-70B, Qwen-2.5-14B and Gemma-2-9B: corruption localised to an ``operation planning'' band; Filler (same numbers) versus NoOp clauses; head ablation and amplification & A distinct functional mechanism, or a substantive generalisation beyond math word problems; a result that reduces to operation planning in physics templates is reported as such \\
Anchoring pathways (Owusu, Wiegreffe and Feldman, arXiv 2606.12818) & Irrelevant numbers shift judgements in 7--8B Qwen and Llama; logit-difference metric; edge-level attribution recovers the signal better than node-level; transfer across anchors but less across base and instruction-tuned models & A physically meaningful dependency-selection account beyond numerical anchoring: role, not only value, must predict the behaviour \\
The Readout Shortcut (Liu, arXiv 2605.22870) & 1--3B models copy the number nearest the answer delimiter; framing and position drive copying & Evidence that the phenomenon is not explained by positional readout copying: position is a controlled factor and T1 readout effects are separated from T2/T3 \\
GSM-Symbolic NoOp (Mirzadeh et al., ICLR 2025 per arXiv); GSM-IC (Shi et al., ICML 2023); GSM-Plus (Li et al., ACL 2024); MathCheck (Zhou et al., ICLR 2025) & Template-generated or manual distractors that do not change the answer; large accuracy drops & Physics-certified roles and the role-by-value design; distractibility as such is not a contribution \\
ProbleMathic (Anantheswaran et al., ICLR 2025 workshop) & Irrelevant numeric variables with a relevance taxonomy; added variables built not to share a unit with the originals & Unit and role manipulated as factors, with certificates \\
Numerical sensitivity (Sun et al., arXiv 2511.08022) & Sentences with numbers hurt far more than sentences without & The presence contrast is replication, not novelty \\
Inverse scaling in test-time compute (Gema et al., TMLR 2025 per arXiv); missing premises (Fan et al.) & Longer reasoning can lower accuracy with distractors; ill-posed questions inflate length; Qwen3-32B scaled positively on controlled ``Misleading Math'' & Length effects (H5) are known; a cost result must be framed as such \\
Physics of LMs 2.1 (Ye et al., ICLR 2025) & Probes show synthetic-trained models know which parameters are necessary & Pretrained models, physical roles, causal (not probing) evidence \\
Premise removal (Meadows et al., Findings of EMNLP 2024 per arXiv) & Physics inference often ignores physical context & Added-quantity role manipulation with certified response \\
Binding: Feng and Steinhardt (ICLR 2024); Gur-Arieh, Geva and Geiger (ICLR 2026 per arXiv); Wu, Geiger and Milli\`ere (ICML 2025 per arXiv) & Binding IDs; positional, lexical and reflexive mechanisms combine; variable dereferencing & Role binding is likely distributed; the causal standard does not assume one feature \\
Copy suppression (McDougall et al.); competition of mechanisms (Ortu et al., ACL 2024 per arXiv); Stochastic Chameleons (Cheng et al., ACL 2025) & Negative heads suppress copied tokens; competing context and recall pathways & ``Two pathways compete'' is not new; the physics-role variable must be \\
\bottomrule
\end{longtable}
\end{center}

Physics benchmarks found without irrelevant-quantity tests (per their abstracts): UGPhysics (ICML 2025), PhysReason
(ACL 2025), PHYSICS (Findings of ACL 2025), SciBench (ICML 2024), PhySense, PhysicsEval, SeePhys.

\subsection{What could actually be novel (conditional claims)}
\begin{description}
\item[N1 (revised).] A role-by-value behavioural account on certified physics strata: whether a model's answers
follow the certified dependence on a quantity's value (invariant for D, C, K; the certified response for R) across
phrasings, separated from numeral salience, plausibility, unit, position and marker wording, and contrasted with a
non-physics dependency control. \emph{Risk:} if the pattern matches GSM-DC or NoOp behaviour in math, the physics
framing adds little.
\item[N2.] A mechanistically predictive function for how physical role gates a quantity's use (or how correct
selection is implemented, H6), established under the causal standard of Section~\ref{sec:causal} on held-out
structures and a second model. \emph{Risk:} it reduces to operation planning (Qu and Gomes), anchoring (Owusu et
al.) or readout copying (Liu).
\item[N3 (demoted).] Whether a single-prompt internal monitor adds information beyond deployable comparators under
cluster-held-out evaluation. \emph{Risk:} small; G7 only.
\item[N4.] A selective intervention that improves certified invariance or accuracy without damaging relevant use
and beats cheap alternatives. \emph{Risk:} a raw accuracy gain without mechanistic linkage belongs to robust
training.
\item[N5 (new, modest).] A documented case of how teacher-forced string likelihood misled a physics-reliance
study (the Phase-I lesson), supported by the reviewers' reproduced audit and Stage~14; consistent with evidence that
what models encode and what they say can diverge (Yuchi, Du and Eisner, EACL 2026).
\end{description}

\textbf{Novelty falsification before any submission.} Re-run targeted searches (certified physics distractor;
physical relevance counterfactual; role-by-value; numerical competition heads; operation planning distractor;
anchoring circuits; readout copying; causal abstraction scientific reasoning); inspect experiment tables and code
of every close hit; keep \texttt{closest\_work\_matrix.csv} with year and venue, generator, formal guarantee,
behavioural target, models, localisation method, functional evidence, external validity and intervention.

% ===== s12_sprint =====
\section{The first execution sprint}
\label{sec:sprint}

The October sprint (tag and reproduce Phase~I; base-rotation attack; scoring interface; structure overlap;
validator qualification; typed generator; pilot) was executed as Stages 14.0--14.19 (Appendix~\ref{app:stage14}),
except that validator qualification stopped at materials and the user's review of the planted-defect bank, and the
physicist review was replaced by agent verification at the user's instruction of 26 October 2026. The next sprint
is ordered by which uncertainty would make later work meaningless. Only the parity, pilot and calibration items
need a GPU; nothing is submitted
without the user's acceptance of the exact command.

\begin{enumerate}
\item \textbf{G0 status sync.} Stage~14 status and roadmap through 14.20; ledger rows for jobs 435605 and 435606;
journal index.
\item \textbf{G0 review additions} (a)--(g) of Section~\ref{sec:gates}, as scripts with tests where code is
involved.
\item \textbf{Qwen3-32B provenance}: revision identification and manifest (read-only on scratch).
\item \textbf{G1 construct}: the strata taxonomy with certificates; the role-by-value generator (values, two or
more phrasings, misuse curves, feasibility records) with mutation tests; regenerated proximity conditions;
sampled and batched reasoning contract with the O3 status categories and the LaTeX final-line recovery rule for
O2 (prospective only); frozen analysis, scorecard and power-simulation scripts tested on synthetic records; SESOI
and equivalence margins with reasons; experiment cards.
\item \textbf{LLM review panel}: choose the judge models (D2), freeze prompts, extend the planted-defect bank with
fresh valid items and new planted defects plus a held-out portion, and qualify the judges (Section~\ref{sec:llmpanel}).
\item \textbf{M0 pilot structures} (8--12) with role-by-value cells, admitted under Section~\ref{sec:datasetplan}.
\item \textbf{A small non-physics dependency control set} (GSM-DC-style) built with the same role-by-value logic.
\item \textbf{G0 extended parity} at the planned hook sites and trace lengths for Qwen2.5-7B-Instruct and
Llama-3.1-8B-Instruct (one short GPU job).
\item \textbf{G2 design-validation pilot} on the M0 structures (one GPU job).
\item \textbf{M250 build} (CPU and LLM panel): upgrade the 80 v1 structures with role-by-value cells and re-validate
them (EXPLORATION), add 50 new EXPLORATION and 120 new sealed INTERNAL\_CONFIRMATION structures; composition
reported against the targets of Section~\ref{sec:datasetplan}.
\item \textbf{G2 calibration} on the two 7--8B models (one GPU job each).
\end{enumerate}

Decisions that change this sprint are listed in Appendix~\ref{app:disposition} (open decisions D1--D6).

% ===== s13_papers =====
\section{Proposed paper architectures depending on the evidence}
\label{sec:papers}

The October edition proposed three architectures: a causal mechanism paper (A, high bar), a measurement paper on
what certified counterfactuals measure (B), and a mechanism-guided robustness paper (C). They are mapped here onto
the revised hypotheses (Section~\ref{sec:hypotheses}) and claim hierarchy (Section~\ref{sec:claims}). Which paper
is written is decided by the evidence; the venue is chosen after the result, not before. None of these
architectures is a plan to obtain a particular outcome.

\begin{center}\small
\begin{tabularx}{\linewidth}{@{}P{3.4cm}LLL@{}}
\toprule
Architecture (outcome) & Core chain & Required evidence (gate) & Cannot be claimed from \\
\midrule
\textbf{P-H0} Certified null: physical dependency without numerical salience &
Role-by-value design $\to$ invariance in D, C, K and correct response in R, per stratum $\to$ equivalence against a
margin fixed in advance $\to$ the same on held-out structures and a second model &
G4 with the power simulation published; equivalence shown, not assumed; \NE{} cells reported apart &
An underpowered non-significant contrast; an inconclusive interval; zero evaluable rows \\
\textbf{P-H1/H2} Numeric salience (presence, plausibility, unit) &
Copy and substitution curves across roles and units $\to$ plausibility-matched and incompatible-unit controls
$\to$ behaviour framed against Liu (positional readout copying) and Owusu, Wiegreffe and Feldman (anchoring
pathways), stating what is beyond them &
G3 discovery confirmed at G4; mechanism optional (G5) and then judged against those works &
A single answer equal to a misuse value; a likelihood shift without an answer change; an effect those works
already explain, reported as new \\
\textbf{P-H3/H4} Relevance-selection failure or readout contamination, with mechanism (October A) &
Role-by-value departure following a misuse curve $\to$ T2 or T3 causal mediation with the interaction estimand
$\to$ within-structure rescue $\to$ causal abstraction (DAS with baselines) $\to$ held-out prediction $\to$ second
model &
G4 behaviourally, G5 for the mechanism, G6 for generality; ``shortcut reliance'' only under T2 or T3 &
One component knockout; an attention map; probe accuracy; a correlated direction; a T1 readout metric alone \\
\textbf{P-H5} Computational cost of irrelevant quantities &
O5 (length, budget exhaustion, thinking closure) rises while O1 and O4 are equivalent to no-cue $\to$ thinking on
versus off and budget sweep within one set of weights &
G4 with equivalence on O1 and O4; the D3 model if the within-weights contrast is claimed &
Mean versus median comparisons; truncated records relabelled wrong; single-sample runs \\
\textbf{P-H6} Mechanism of correct selection &
H6-RD (D and R; H6-extended adds C and K) $\to$ T2 or T3 localisation of the computation that excludes the
quantity (or uses it correctly in C and K) $\to$ held-out prediction &
G4 behaviourally, then G5; on the population of tasks solved reliably without a cue &
Decoding a role from a subspace (diagnostic only); intermediate use in C or K called an error \\
\textbf{P-M} Methodology: what teacher-forced string likelihood measures (October B) &
Phase-I label ambiguity $\to$ base, wrapper and tokenisation audit $\to$ divergence between teacher-forced
string likelihood (T1, O6) and sampled answers (O1--O4) $\to$ recommended protocol &
G0--G2 evidence on committed data, with the review findings; a generalisable diagnostic &
More variants of the same items; a minor score change \\
\textbf{P-R} Mechanism-guided selective robustness (October C; G7) &
Verified mechanism $\to$ cheap baselines (including the relevance instruction) and component-aware repair $\to$
held-out role reversals $\to$ quality, compute and damage trade-off &
Only after G5 and G6; real gains on O1 or O2 with external generalisation &
Flipping a teacher-forced score; an untested rule-based filter; excluding the cheap baseline \\
\bottomrule
\end{tabularx}
\end{center}

\textbf{Combinations.} More than one hypothesis can hold (Section~\ref{sec:hypotheses}); a paper states which held
in which stratum and model rather than choosing the most attractive. P-H0 and P-H5 are full results, not fallbacks.

\textbf{Editorial discipline} (October edition, kept). A main paper has one decisive thesis with appendices
supporting it; a programme may yield several publishable projects, and forcing them into one paper weakens
clarity. The withdrawn Phase-I paper's reviews (Appendix) are the reference for what a reviewer will check first:
whether the metric isolates the claimed function. This plan promises a scientifically credible process, not an
acceptance probability.

% ===== s14_bibliography =====
\section{Bibliography and primary-source links}
\label{sec:bibliography}

{\sloppy
\textbf{Provenance of the entries.} Entries \textbf{[L01]--[L30]} are the October edition's bibliography, kept with
its numbering, venues and URLs exactly as given there (located and cross-checked against publisher,
author-project or arXiv pages on 14 October 2026; ``preprint'' means not independently verified as peer-reviewed
at that date). Entries \textbf{[N01]--} were \textbf{added on 11 November 2026} from the revision's literature
search (three search agents). Their venue status is marked as the revision marks it: \textbf{[R]} the venue or
arXiv page was opened by a search agent and the entry confirmed; \textbf{per arXiv} only the arXiv comment
claims the venue; \textbf{[R, re-opened]} additionally re-opened by the planning agent (the two entries closest to
the novelty boundary). Entries listed by the revision without either mark are given as listed and are
\emph{unverified} here. Where the revision gives no URL, only the arXiv identifier is given, and where it gives
neither, only the citation as listed; no URL is supplied from memory. Every source must be read in full before it
is used to claim or deny novelty; absence from this list is not proof that no similar work exists.

\subsection*{Irrelevant context, distractors and robustness of reasoning}
\begin{enumerate}[label={},leftmargin=0pt]
\item \textbf{[L01]} I.~Mirzadeh et al.\ (2025). GSM-Symbolic: Understanding the Limitations of Mathematical
Reasoning in Large Language Models. ICLR.
\url{https://proceedings.iclr.cc/paper_files/paper/2025/hash/ec2e7a896f8250986b3907f57621ce94-Abstract-Conference.html}
(the revision lists its NoOp variant as ``ICLR 2025 per arXiv'' [R]).
\item \textbf{[L02]} M.~Yang et al.\ (2025). How Is LLM Reasoning Distracted by Irrelevant Context? An Analysis
Using a Controlled Benchmark. EMNLP, pp.~13329--13347. \url{https://aclanthology.org/2025.emnlp-main.674/}
(the revision's GSM-DC entry, ``Yang et al., EMNLP 2025'' [R], has the same first author and venue).
\item \textbf{[L03]} Z.~Cheng et al.\ (2025). Stochastic Chameleons: Irrelevant Context Hallucinations Reveal
Class-Based (Mis)Generalization in LLMs. ACL, pp.~30187--30214. \url{https://aclanthology.org/2025.acl-long.1458/}
\item \textbf{[N01]} Shi et al., Large Language Models Can Be Easily Distracted by Irrelevant Context (GSM-IC). ICML 2023, PMLR 202. \url{https://proceedings.mlr.press/v202/shi23a.html} [R] \emph{Added 11 Nov 2026.}
\item \textbf{[N02]} Li et al., GSM-Plus. ACL 2024. \url{https://aclanthology.org/2024.acl-long.163} [R] \emph{Added 11 Nov 2026.}
\item \textbf{[N03]} Zhou et al., Is Your Model Really A Good Math Reasoner? (MathCheck). ICLR 2025. \url{https://iclr.cc/virtual/2025/poster/28416} [R] \emph{Added 11 Nov 2026.}
\item \textbf{[N04]} Anantheswaran et al., Cutting Through the Noise (ProbleMathic). ICLR 2025 Workshop on Reasoning and Planning. arXiv 2406.15444. [R] \emph{Added 11 Nov 2026.}
\item \textbf{[N05]} Sun et al., numerical sensitivity (numeric versus non-numeric distractors). arXiv 2511.08022.
[R] \emph{Added 11 Nov 2026.}
\item \textbf{[N06]} Gema et al., inverse scaling in test-time compute. TMLR 2025, per arXiv. [R] \emph{Added 11 Nov
2026.}
\item \textbf{[N07]} Ye et al., Physics of Language Models, Part~2.1. ICLR 2025. arXiv 2407.20311. [R] \emph{Added 11 Nov 2026.}
\item \textbf{[N08]} Meadows et al., premise removal in physics derivations. Findings of EMNLP 2024, per arXiv. [R]
\emph{Added 11 Nov 2026.}
\item \textbf{[L21]} R.~Geirhos et al.\ (2020). Shortcut Learning in Deep Neural Networks. Nature Machine
Intelligence 2:665--673. \url{https://doi.org/10.1038/s42256-020-00257-z}
\end{enumerate}

\subsection*{Closest mechanistic accounts of distractor and numeral use}
\begin{enumerate}[label={},leftmargin=0pt]
\item \textbf{[N09]} Qu and Gomes, A Four-Stage Decomposition of Word-Problem Solving and Mechanistic Fragility in LLM Math Reasoning. Findings of EMNLP 2026. \url{https://aclanthology.org/2026.findings-emnlp.820/}
[R, re-opened] \emph{Added 11 Nov 2026.}
\item \textbf{[N10]} Owusu, Wiegreffe and Feldman, anchoring pathways. arXiv 2606.12818. [R, re-opened]
\emph{Added 11 Nov 2026.}
\item \textbf{[N11]} Liu, The Readout Shortcut. arXiv 2605.22870. [R] \emph{Added 11 Nov 2026.}
\item \textbf{[L04]} Y.~Nikankin et al.\ (2025). Arithmetic Without Algorithms: Language Models Solve Math with a
Bag of Heuristics. ICLR.
\url{https://proceedings.iclr.cc/paper_files/paper/2025/hash/8c5f30296296d2ae402ebbd09aaa9c12-Abstract-Conference.html}
\item \textbf{[L20]} S.~Naganna, T.~A.~Sijan, U.~Kalita (2026). Are Arithmetic Heuristic Neurons Form-Invariant? A
Mechanistic Analysis of Symbols, Text, and Code in LLMs. arXiv preprint. \url{https://arxiv.org/abs/2607.16693}
\end{enumerate}

\subsection*{Physics benchmarks and grading}
\begin{enumerate}[label={},leftmargin=0pt]
\item \textbf{[N12]} W.~Zhao, Q.~Ma, J.~Shi, S.~Wu, J.~Han, Y.~Xiao, S.-Y.~Chen, X.~Luo, L.~Schmidt, J.~Y.~Zou (2026).
PRISM-Physics: Causal DAG-Based Process Evaluation for Physics Reasoning. ICLR 2026. arXiv 2510.03185.
\url{https://proceedings.iclr.cc/paper_files/paper/2026/hash/1a379b51d7ec80b32d8c14f40c281ce1-Abstract-Conference.html}
\emph{Added 11 Nov 2026; proceedings page opened on 11 Nov 2026 while writing this section (title, authors and
venue confirmed); the arXiv page was not opened.}
\item \textbf{[L26]} Yale NLP. PHYSICS-Verified: Audited Advanced Physics Benchmark (public dataset card; inspected
14 October 2026). 1{,}034 qualifying-exam problems, 2{,}559 scored answers; newly adapted counterfactuals are not
automatically validated by this source. \url{https://huggingface.co/datasets/yale-nlp/physics-verified}
\item \textbf{[L30]} Yale NLP. PHYSICS Benchmark Repository and PHYSICS-Verified migration guidance (audited
source and reuse notes; inspected 14 October 2026). \url{https://github.com/yale-nlp/Physics} (the revision lists
PHYSICS as Findings of ACL 2025, unmarked).
\item \textbf{[L27]} X.~Xu et al.\ (2025). UGPhysics: A Comprehensive Benchmark for Undergraduate Physics
Reasoning with Large Language Models. ICML, PMLR 267:69849--69877.
\url{https://proceedings.mlr.press/v267/xu25ai.html}; official repository
\url{https://github.com/YangLabHKUST/UGPhysics}; grading conventions per arXiv 2502.00334 (revision).
\item \textbf{[L28]} X.~Zhang et al.\ (2025). PhysReason: A Comprehensive Benchmark towards Physics-Based
Reasoning. ACL 2025, pp.~16593--16615. \url{https://aclanthology.org/2025.acl-long.811/}; grading per arXiv
2502.12054 (revision).
\item \textbf{[L29]} X.~Wang et al.\ (2024). SciBench: Evaluating College-Level Scientific Problem-Solving
Abilities of Large Language Models. ICML 2024. \url{https://proceedings.mlr.press/v235/wang24z.html}
\item \textbf{[N13]} PhySense, PhysicsEval, SeePhys: physics benchmarks found by the revision's search without
irrelevant-quantity tests (per their abstracts); listed by name only, no venue or identifier given.
\emph{Added 11 Nov 2026; unverified here.}
\item \textbf{[N14]} Grading tools: Math-Verify; TinyV (arXiv 2505.14625); pint 0.26. \emph{Added 11 Nov 2026;
unmarked.}
\end{enumerate}

\subsection*{Causal and mechanistic interpretability methodology}
\begin{enumerate}[label={},leftmargin=0pt]
\item \textbf{[L05]} F.~Zhang and N.~Nanda (2024). Towards Best Practices of Activation Patching in Language
Models: Metrics and Methods. ICLR.
\url{https://proceedings.iclr.cc/paper_files/paper/2024/hash/06a52a54c8ee03cd86771136bc91eb1f-Abstract-Conference.html}
\item \textbf{[L22]} S.~Heimersheim and N.~Nanda (2024). How to Use and Interpret Activation Patching. arXiv
preprint. \url{https://arxiv.org/abs/2404.15255}
\item \textbf{[L06]} J.~Miller, B.~Chughtai, W.~Saunders (2024). Transformer Circuit Faithfulness Metrics are not
Robust. arXiv preprint. \url{https://arxiv.org/abs/2407.08734} (the revision lists COLM 2024, unmarked).
\item \textbf{[L09]} A.~Syed, C.~Rager, A.~Conmy (2023). Attribution Patching Outperforms Automated Circuit
Discovery. arXiv preprint. \url{https://arxiv.org/abs/2310.10348} (the revision lists BlackboxNLP 2024,
unmarked).
\item \textbf{[L07]} A.~Geiger et al.\ (2025). Causal Abstraction: A Theoretical Foundation for Mechanistic
Interpretability. JMLR 26(83):1--64. \url{https://www.jmlr.org/papers/v26/23-0058.html}
\item \textbf{[L08]} A.~Mueller et al.\ (2025). MIB: A Mechanistic Interpretability Benchmark. ICML / PMLR
267:45069--45108. \url{https://proceedings.mlr.press/v267/mueller25a.html}
\item \textbf{[L10]} T.~Bricken et al.\ (2025). Circuit Tracing: Revealing Computational Graphs in Language Models.
Transformer Circuits technical report.
\url{https://www.transformer-circuits.pub/2025/attribution-graphs/methods.html}
\item \textbf{[L11]} J.~Lindsey et al.\ (2025). On the Biology of a Large Language Model. Transformer Circuits
technical report. \url{https://www.transformer-circuits.pub/2025/attribution-graphs/biology.html}
\item \textbf{[L12]} T.~McGrath et al.\ (2023). The Hydra Effect: Emergent Self-repair in Language Model
Computations. arXiv preprint. \url{https://arxiv.org/abs/2307.15771}
\item \textbf{[L18]} J.~Chen, Y.~Luo, L.~Pan (2026). Mechanistic Data Attribution: Tracing the Training Origins of
Interpretable LLM Units. ICML / PMLR 306:15966--15998. \url{https://proceedings.mlr.press/v306/chen26cv.html}
\item \textbf{[L19]} A.~A.~Khan et al.\ (2026). Can Language Model Agents be Helpful Circuit Explainers in
Mechanistic Interpretability? arXiv preprint. \url{https://arxiv.org/abs/2606.24026}
\item \textbf{[L23]} TransformerLens (repository implementation). patching.py.
\url{https://github.com/TransformerLensOrg/TransformerLens/blob/main/transformer_lens/patching.py}
\item \textbf{[L24]} Anthropic (2025). Open-sourcing Circuit Tracing Tools.
\url{https://www.anthropic.com/research/open-source-circuit-tracing}
\item \textbf{[N15]} Added 11 Nov 2026 (revision's methodology list, unmarked): Kram\'ar et al., AtP* (arXiv
2403.00745); Rushing and Nanda (ICML 2024); Li and Janson (NeurIPS 2024); Grant et al., divergent representations
(ICLR 2026); Anani et al., Certified Circuits (ICML 2026); Geiger et al., DAS (CLeaR 2024); Makelov et al.\ (ICLR
2024); Bogdan et al., Thought Anchors (arXiv 2506.19143); Macar et al., Thought Branches (arXiv 2510.27484);
Yuan et al., Hidden Error Awareness (arXiv 2605.09502); Shi et al., hypothesis testing the circuit hypothesis
(arXiv 2410.13032); M\'eloux et al.\ (arXiv 2510.00845, 2512.18792).
\end{enumerate}

\subsection*{Binding and number representation (framing for G5)}
\begin{enumerate}[label={},leftmargin=0pt]
\item \textbf{[N16]} Added 11 Nov 2026: Feng and Steinhardt (ICLR 2024); Prakash et al.\ (ICLR 2024, per arXiv);
Gur-Arieh, Geva and Geiger (ICLR 2026, per arXiv); Wu, Geiger and Milli\`ere (ICML 2025, per arXiv); McDougall et
al., copy suppression (arXiv 2310.04625); Ortu et al.\ (ACL 2024, per arXiv); Levy and Geva (NAACL 2025, per
arXiv); Yuchi, Du and Eisner (EACL 2026); Wu, Shen and Jiang, I3C (NAACL 2024). Nikankin et al.\ is [L04].
\end{enumerate}

\subsection*{Uncertainty, hallucination detection and monitoring (G6 side track)}
\begin{enumerate}[label={},leftmargin=0pt]
\item \textbf{[L13]} J.~Kossen et al.\ (2024). Semantic Entropy Probes: Robust and Cheap Hallucination Detection
in LLMs. arXiv preprint. \url{https://arxiv.org/abs/2406.15927}
\item \textbf{[L14]} C.~Chen et al.\ (2024). INSIDE: LLMs' Internal States Retain the Power of Hallucination
Detection. ICLR. \url{https://mlanthology.org/iclr/2024/chen2024iclr-inside/}
\item \textbf{[L15]} R.~Khanmohammadi et al.\ (2025). Calibrating LLM Confidence by Probing Perturbed
Representation Stability. EMNLP. \url{https://aclanthology.org/2025.emnlp-main.530/}
\end{enumerate}

\subsection*{Statistics, evaluation noise and contamination}
\begin{enumerate}[label={},leftmargin=0pt]
\item \textbf{[N17]} Added 11 Nov 2026 (unmarked): Miller, ``Adding error bars to evals'' (arXiv 2411.00640);
Bowyer, Aitchison and Ivanova (ICML 2025, position paper on small-$N$ intervals); Wang, ``Measuring all the noises
of LLM evals'' (arXiv 2512.21326); Heineman et al., signal and noise (arXiv 2508.13144); Hochlehnert et al., ``A
sober look at progress in LM reasoning'' (arXiv 2504.07086); Matuschek et al.\ (J.\ Mem.\ Lang.\ 2017); Mancl and
DeRouen (Biometrics 2001); Lakens, Scheel and Isager (AMPPS 2018, TOST); Lan--DeMets group-sequential spending
(gsDesign).
\item \textbf{[L16]} Y.~Sun et al.\ (2025). The Emperor's New Clothes in Benchmarking? A Rigorous Examination of
Mitigation Strategies for LLM Benchmark Data Contamination. ICML / PMLR 267.
\url{https://proceedings.mlr.press/v267/sun25t.html}
\item \textbf{[L17]} S.~Chen et al.\ (2025). Benchmarking Large Language Models Under Data Contamination: A Survey
from Static to Dynamic Evaluation. EMNLP. \url{https://aclanthology.org/2025.emnlp-main.511/}
\end{enumerate}

\subsection*{Models, inference and tooling}
\begin{enumerate}[label={},leftmargin=0pt]
\item \textbf{[N18]} Added 11 Nov 2026 (unmarked): Qwen3 technical report (arXiv 2505.09388) and model cards
(sampling settings; licences: Qwen3 and Qwen2.5 Apache-2.0, Llama~3.1 community licence); vLLM installation
documentation and PyPI (release 0.33.0, 2026-11-02); vLLM quantized KV-cache documentation; TransformerLens 3.4.0
(2026-06-09, the pinned version) and 4.3.0 (2026-10-23) release information; circuit-tracer (transcoders for
Qwen3 0.6B--14B and Llama-3.1-8B-Instruct, none for Qwen2.5 or Qwen3-32B).
\end{enumerate}

\subsection*{Venue guidelines}
\begin{enumerate}[label={},leftmargin=0pt]
\item \textbf{[L25]} ICLR 2027. Reviewer Guidelines and Author Guidelines.
\url{https://iclr.cc/Conferences/2027/ReviewerGuidelines};
\url{https://iclr.cc/Conferences/2027/AuthorGuidelines}
\end{enumerate}

\textbf{Attached-source index (not public citations).} From the October edition: the submitted manuscript
\emph{Confidence Without Cause} (25~pp.); the June original proposal (30~pp.); the supplementary change log
(14~pp.); the historical June Part-II plan v1.1 (42~pp.); the supplied OpenReview/PAT LLM review (5~pp.); and a
separate 29-page \emph{Proved, Not Annotated} manuscript of undetermined relationship to the submission. Added in
this revision: the ICLR 2027 reviews of the withdrawn submission with their verification against committed data
(\texttt{docs/reviews/iclr2027\_submission50608/}) and the Stage~14 journal (\texttt{docs/journal/stage14/}).
These artifacts motivate the evidence ledger and critiques; this document is a proposed protocol, not an
attestation that suggested studies have been run.
\par}

% ===== s15_completion =====
\section{Definition of scientific completion}
\label{sec:completion}

A manuscript is scientifically ready only after the gate that its central claim needs has been passed
(Section~\ref{sec:gates}), and only at the level of the claim hierarchy (Section~\ref{sec:claims}) that the
evidence reaches: a behavioural role effect needs G4; a causal-mediation claim needs T3, or T2 with the controls of
Section~\ref{sec:causal}, at G5; a mechanism claim needs G5 with rescue, causal abstraction and held-out
prediction; a generalisation claim needs G6. Readiness is judged on independently verifiable artifacts, and the
manuscript must be able to answer every question below.

\begin{enumerate}
\item \textbf{Construct.} What model behaviour changed, in which outcome (O1--O6), in which stratum (R, D, C, K)
and population, and why must the physical answer be invariant (D, C, K) or respond by the certified amount (R)?
Which separate tests establish the natural-language physical validity of the items and their match to the encoded
graph, and was that validation independent (G1; B4 review at G4)?
\item \textbf{Alternatives.} Which weaker account was strongest --- generic perturbation (H0), numeral presence or
plausibility and unit (H1, H2), length or diversion (H5), string-scoring, format or difficulty artifacts --- and
where was it falsified, including by the non-physics control?
\item \textbf{Role by value.} Does the model keep its answers fixed where a quantity is irrelevant (no tracking loss, no misuse) and follow the certified response where it is relevant (no neglect), as the physics requires,
across values and across phrasings with and without marker words; are misuse hypotheses supported only by curves
across several values, with all matches kept?
\item \textbf{Computation.} If a mechanism is claimed: which computation was identified under a declared claim
type (T1, T2, T3), beyond the effect of removing a component, with the interaction estimand $I_h$ and the clean
cue effect reported side by side?
\item \textbf{Falsifiable predictions.} Which intervention predictions, committed before the test, would have
failed if the account were false, and what happened on held-out structures, values and the second model?
\item \textbf{Scope.} Which models, budgets, strata, conditions and thinking settings do and do not replicate, and
which models were uninformative because they fell outside the difficulty band?
\item \textbf{Monitoring and repair} (only if claimed, G7). Did a monitor add information beyond fair, cheap,
deployable baselines on unseen clusters; can an intervention improve O1 or O2 without disrupting legitimately
relevant quantities, against the relevance instruction and a cheap filter?
\item \textbf{Reproducibility.} Can an independent researcher reproduce successes and nulls, with denominators,
exclusions and \NE{} cells, from pinned model manifests, \texttt{SUBMIT\_COMMIT}-stamped raw records and the frozen
analysis scripts?
\item \textbf{Novelty.} Is the core insight different from GSM-DC, GSM-Symbolic, PRISM-Physics, Qu and Gomes,
Owusu, Wiegreffe and Feldman, Liu, numerical-heuristic accounts and existing causal interpretability work, and is
any overlap stated?
\item \textbf{Wording.} Does the manuscript say no more than the evidence supports, as checked by the wording
script against the claim registry and the evidence ledger, with agent authorship and agent verification
disclosed?
\end{enumerate}

\textbf{Consequences.} If (4) or (5) cannot be answered, the project can still produce rigorous measurement
science (a behavioural role paper, a certified null, a cost result or the methodology paper of
Section~\ref{sec:papers}), but it is not presented as a mechanistic explanation. If (1) fails, the work returns to
the construct and measurement definition, however strong other results look. A certified null (H0 with
equivalence shown against a margin fixed before the data and an adequately powered design) and a cost finding (H5)
are complete results; an inconclusive interval, an underpowered contrast or zero evaluable rows are not, and are
reported as what they are. Results without LLM-panel or human validation are complete only as exploratory results.

\appendix
% ===== a1_stage14 =====
\section{Stage 14 evidence (exploratory)}
\label{app:stage14}

All rows have data role EXPLORATION; nothing here is paper-eligible. Journal: \texttt{docs/journal/stage14/}
(the Stage~14.20 entries are in \texttt{14.19\_gold\_dataset.md}); status: \texttt{docs/stage14/status\_and\_roadmap.md};
integrity check: \texttt{scripts/v2/register\_stage14\_artifacts.py --verify}.

\begin{center}\small
\begin{longtable}{@{}P{1.5cm}P{8.8cm}P{5.5cm}@{}}
\toprule
Stage & Result \tg{V} & Bearing on this plan \\
\midrule\endhead
14.0--14.1 & Claim ledger and contract; X01 recomputed 51 Phase-I components from committed records; all reproduce except three: ``25 shared positives'' is the $\hat S$ overlap (0.5-nat overlap 22), the +0.067 combined gain cannot be recomputed from committed predictions, and the 71\% had no committed source (later 110/155) & Section~\ref{sec:ledger} \\
14.2 & X02 base rotation: 40--49\% (Qwen) and 41--56\% (Llama) of base-0 positives flip; legacy label retired by the user (J14.2-011) & Legacy endpoint retired \\
14.3 & Scoring contract with seven wrappers; Llama double BOS in 560/560 prompts; scored string lacks the ``Answer:'' prefix that 1,120/1,120 generations start with & Grading contract (O3) \\
14.4 & 34 structure-overlap clusters; 95/135 = 70.4\% exposure; template predictors lose signal under cluster-held-out folds, entropy features do not; monitor itself \NE{} (no activations) & G7 side track \\
14.5 & Validator materials; 33-item planted-defect bank reviewed by the user (27 defective, 6 valid controls); no validator qualified & B4 open (D2) \\
14.6, 14.15, 14.18 & Typed generator and certifier; structures S01--S30 (mutation tests: no mutant missed); S01--S20 approved in the user's walkthroughs & G1 builds on it \\
14.7 & Two-model pilot (422 records); without reasoning 3/211 (Qwen) and 6/211 (Llama) correct & Floor without reasoning \\
14.9 & bfloat16 reproducible; dtype noise SD 0.085--0.186 nats; same-unit cue raises the decoy by about 10--13.5 nats (Qwen) and 3.5--5.3 (Llama), with no detectable proximity dependence & Diagnostic only \\
14.10, 14.16 & X03: the sign of the presence effect holds in all seven wrappers (v2 structures) and in four wrappers on 60 legacy families; size and parse rates depend on format & Format as a factor \\
14.12 & Unit compatibility modulates the boost with the same digits & H2 \\
14.13 & Greedy reasoning at 2,048 tokens: Qwen 96/211, Llama 50/211 correct (43 Llama records ungraded); a sentence with no number flips answers about as often as numeric cues & Sampling and self-flip floor required \\
14.14, 14.17 & Every number in a prompt is boosted as an answer, a same-unit cue most; matching plausibility removes about 40\% (Qwen) and 53\% (Llama) of the same-unit effect; incompatible-unit and metadata effects become not distinguishable from plausibility; the remainder depends on unit compatibility & H1/H2 at the likelihood level \\
14.19 & Gold dataset v1: 260 samples from 80 structures in 11 domains (50 structures by domain agents), gold traces from the typed graph, agent verification in place of the physicist review; all \texttt{UNREVIEWED}; non-governing cues carry explicit disclaimers in the text & EXPLORATION only \\
14.20 & Qwen3-32B thinking, one sample each: 63/64 graded correct (9 budget-exhausted, 5 LaTeX-wrapped final lines); 32 candidate anchors $\times$ 4 cue kinds: 109 correct, 2 wrong, 17 \NE{}; mean 5,632 versus 4,509 tokens; resample (435606) cancelled, cue contrast \NE{}; 19.8--20.8 tokens/s & Ceiling on easy items; compute constraint \\
\bottomrule
\end{longtable}
\end{center}

% ===== a2_reviews =====
\section{ICLR 2027 review issues and where this plan answers them}
\label{app:reviews}

Issue numbers follow the consolidated register in \texttt{docs/reviews/iclr2027\_submission50608/analysis.md}
(section~4). ``Checked'' gives the reproduced check identifier.

\begin{center}\small
\begin{longtable}{@{}P{0.6cm}P{5.8cm}P{1.5cm}P{7.5cm}@{}}
\toprule
\# & Issue & Checked & Answer in this plan \\
\midrule\endhead
1 & Label scores the bare string, not the instructed format & C1 & Legacy label retired; O1--O3 on the model's own answers; grading contract (Section~\ref{sec:outcomes}) \\
2 & Label unstable across base and instruction & C2 & Retired; G0(a) cross-instruction stability record; self-flip floor \\
3 & A correct formula hint triggers the label & C3 & Formula-hint cell as a positive diagnostic (Section~\ref{sec:conditions}) \\
4 & ``Fitted entropy'' comparator is label-derived & C4, C5b & Oracle/deployable tags (Section~\ref{sec:status}, rule 6; G0(d)) \\
5 & Knockout one-sided; reversal not cancellation; cue-independent & C8 & Interaction estimand; claim types T1--T3; no normalised recovery (Section~\ref{sec:causal}); X07 retired as stated \\
6 & Heads chosen and tested on the same panel & -- & Discovery/test split; selection stability; held-out prediction (Section~\ref{sec:causal}) \\
7 & Wording exceeds evidence & -- & Claim hierarchy; G0(e) wording check; forbidden-phrase check \\
8 & Template-cluster leakage & C5c, C5d & Structure as unit; disjoint data roles; G4 on untouched clusters \\
9 & Cue C pooled with irrelevant cues & C6 & Convertible-unit family kept separate \\
10 & No human validation; agent authorship & C9 & Validation tiers (Section~\ref{sec:status}, rule 4); a qualified LLM panel from other model families (Section~\ref{sec:llmpanel}); later human study; G0(f) provenance consistency; disclosure \\
11 & Near-match claim overstated & C7 & Proximity as distinct conditions; role-by-value design \\
12 & Verifier scope and tolerance & C10 & Strata certificates; B1.5 audit if v1 is reused \\
13 & Scope: short plug-in calculations & -- & Problem-structure strata (B3.1); strata C and K require more than plug-in \\
14 & No downstream use & -- & G7 (repair against cheap alternatives) \\
15 & One string's probability is not confidence & C1--C3 & O1--O4; O6 only inside a causal claim type \\
16 & Threshold justification & C2b & Rule 5 of Section~\ref{sec:status}; G0(g) \\
17 & Similar AUROC is not shared information & -- & G7 only, with deployable comparators \\
18 & Winner's curse in layer sweeps; base-rate skew & -- & Discovery/test split; Section~\ref{sec:statistics} \\
19 & Inconsistent model-size statements & -- & Writing hygiene; claim-wording check \\
20 & DeepSeek scoring position & -- & Thinking models graded only after the end-of-thinking marker \\
\bottomrule
\end{longtable}
\end{center}

% ===== a3_disposition =====
\section{Disposition of the October edition and open decisions}
\label{app:disposition}

\subsection{Disposition}
\begin{center}\small
\begin{longtable}{@{}P{4.0cm}P{2.4cm}P{9.4cm}@{}}
\toprule
October element & Disposition & Note \\
\midrule\endhead
Rules of evidence (p.~3) & Kept & Section~\ref{sec:status}, with the strata corollary and six added rules \\
Data roles, shadow set & Kept & Unchanged \\
B0 threat model, novelty audit & Kept, partly done & Section~\ref{sec:literature} \\
B1 reproduce and break v1 & Closed & Stage~14 and the reproduced reviews; B1.5 only if v1 is reused \\
B2 construct choice & Decided & Section~\ref{sec:framework} \\
B3 typed generator & Extended & Strata, role-by-value, phrasings, misuse curves \\
B4 human validation & Kept; interim LLM-panel tier & Section~\ref{sec:llmpanel}; human study later; D2 \\
B5 external sets & Split & Non-physics control in G3; external physics in G6 \\
B6 pilot, power, roles & Restructured & G2 and G3, staged \\
B7 Phase-B go/no-go & Replaced & G3 scorecard and the fixed G4 family \\
P0 full reproduction & Narrowed & Parity kept; full regeneration off the critical path (D1) \\
P1 measurement validity & Largely done & Findings kept as diagnostics; P1.4 and P1.8 remain optional \\
P2 scale to 300--600 & Replaced & 250 structures, then 500 before G4 (Section~\ref{sec:datasetplan}) \\
P3 validation, parser, leakage & Kept & O3 status categories; parser audit \\
P4 behavioural discrimination & Central & G3/G4 \\
P5 monitoring & Demoted & G7 \\
P6--P8 & Kept, gated & Causal standard of Section~\ref{sec:causal} \\
P9 generalisation & Kept, gated & Freeze the function, not the index \\
P10 repair & Kept & G7 \\
P11 closure & Kept & External cohort in G6 \\
Inventory A01--F20 & Kept as candidates & Statuses in Section~\ref{sec:inventory} \\
X01--X03 & Done (exploratory) & \\
X04 & Becomes the role-by-value factorial & G3/G4 \\
X05 & Side track & G7 \\
X06 & Split & Internal part in G4, external part in G6 \\
X07 H11 lesion & Retired as stated & Cue-independent effect (C8b) \\
X08--X12 & Kept, gated & \\
Statistics (Section 8) & Revised & Structure-level paired analysis; SESOI with reasons; four-way result language \\
300--600 structure allocation & Replaced & 250 then 500 structures; gates open what their power simulation requires \\
\bottomrule
\end{longtable}
\end{center}

\subsection{Open decisions for the user}
Each default withholds an action; none is irreversible.

\begin{center}\small
\begin{tabularx}{\linewidth}{@{}P{0.7cm}P{6.6cm}L@{}}
\toprule
Id & Decision & Default until decided \\
\midrule
D1 & Keep the full P0 regeneration of Phase-I values off the critical path & Not run; possible later as a side track \\
D2 & Which 2--3 judge models form the LLM review panel (other families than the authoring model; local or via external APIs) & Decided in principle by the user (2026-11-11): LLM panel now, human study later. Until the models are chosen and qualified, G4 would be EXPLORATORY-REPLICATION \\
D3 & Download Qwen3-8B (optionally Qwen3-14B) with pinned revisions & Not downloaded \\
D4 & Install vLLM in a separate environment with parity and timing checks & Not installed; 7--8B models on Hugging Face \\
D5 & Resubmit the 32-anchor no-cue resample (job 435606 design) & Not resubmitted; lower priority than calibration; no ceiling is assumed, since that is what it would measure \\
D6 & Run the Phase-I monitor under cluster-held-out folds & Not scheduled \\
\bottomrule
\end{tabularx}
\end{center}

% ===== a4_revision_log =====
\section{Revision log}
\label{app:revisions}

\begin{description}
\item[14 October 2026 (v2.0).] Research redesign replacing the June 2026 Phase-II plan; Phase~B added; adopted as
plan of record by the user on 15 October 2026.
\item[11 November 2026, first draft of the revision.] Written as a separate Markdown revision layer
(\texttt{PhysMon\_Research\_v2.1.md}, commit \texttt{2cf08e2}) from the reviews, Stage~14 and three literature
searches. An independent agent review against the repository found ten factual errors (among them: hook parity
had already passed in job 399217; the Stage~14 reasoning runner is greedy, so sampling needs a new contract; a
wrong KV-memory explanation of the Qwen3-32B throughput) and twelve design flaws (no cell varied role alone; the
in-system contrast was confounded with unit; correctness conditioned on evaluability; near cues inside the
grading tolerance; a data-dependent equivalence margin; undefined gate thresholds), all corrected before commit.
\item[11 November 2026, this edition.] A second review by an external language model, supplied by the user and
checked by the agent, led to: the four role strata (disconnected, cancelling, conditional) with the rule that
intermediate use of a cancelling or conditional quantity is not an error; the role-by-value design with a response
scaled by the certified relevant-role response, several validated phrasings and misuse response curves; separate
operational, semantic and status outcomes, with truncation requiring budget exhaustion; near-answer cues kept in
correctness and flagged only for source attribution; all foil matches kept; causal claim types T1--T3; within-structure,
value-matched rescue first; three populations; gate fixes (formula hint not a gate; moderators secondary; G4 family
fixed in advance; scorecard instead of first-matching rule); staged discovery with an 8--12-structure design pilot;
a non-physics control in discovery; SESOI with reasons and four-way null language; PRISM-Physics added to the
closest work. An independent agent review of the assembled LaTeX draft then found that the first role-by-value estimand (a
log-elasticity interaction) was undefined for mixed sampled answers and zero golds and not sign-aligned across
structures, that a shared trace across roles is physically incoherent, that two values cannot separate misuse
curves, and that several sections still carried superseded rules; it was replaced by per-sample tracking, response
fractions scaled by the certified relevant-role response and a role-discrimination index, T1 restricted to
prompts that share a gold, three values per role, and consistent gate rules throughout.
\item[11 November 2026, refinements after a second external review.] The user supplied a further review by the same
external model and asked to rely on LLM reviewers for physics validation for now. Changes: the primary estimands
are per-sample tracking (whose loss catches a departure in either direction), misuse shares with signed,
direction-predicting misuse curves and an R neglect share; the response fractions $\lambda$ and the index $\Lambda$
are descriptive only, with a multimodality guard; G4 tests three co-primary contrasts (D tracking loss, R neglect,
role $\times$ value interaction on tracking) instead of gatekeeping on $\Lambda$; mechanism work uses two
interactions, $I_{\text{presence}}$ (against the non-numeric sentence) and $I_{\text{role}}$ (against the relevant
twin with the same numeral, on a behaviourally comparable outcome), and only the second supports a role-routing
claim; T1 scores the certified answer against a predeclared alternative on one explicitly identified prefix and
excludes no prefix; H6 is split into H6-RD and H6-extended; the primary model is fixed by a prespecified rule at the
end of G2; and a qualified panel of 2--3 independent LLM reviewers from other model families, qualified on the
user-labelled planted-defect bank, is the interim validation tier for G4, with human review as a later accuracy
study.
\item[11 November 2026, dataset size.] At the user's direction, DECOY-v2 dataset~v2 is built to 250 independent
structures (M250, after the G2 design pilot and before calibration) and expanded to 500 after G3 and before G4
(M500), with an admission standard, composition targets and a data-role allocation fixed in advance
(Section~\ref{sec:datasetplan}); this replaces the October 300--600 band and the earlier revision's ``only what the
gates need''. Dataset~v1 (80 structures, 260 instances) is upgraded and counted only as EXPLORATION. At the user's choice of 11 November 2026 the plan was consolidated into this single LaTeX source and
PDF; the Markdown revision layer was removed (preserved in Git history).
\end{description}

\end{document}
